\centering

\section{GAQ‑UFT}\label{gaquft}

\subsection{万物几何统一场论}\label{ux4e07ux7269ux51e0ux4f55ux7edfux4e00ux573aux8bba}

\subsubsection{全维分析与第一性原理推导}\label{ux5168ux7ef4ux5206ux6790ux4e0eux7b2cux4e00ux6027ux539fux7406ux63a8ux5bfc}

\vspace{1.5cm}

\textbf{Geometric‑Axiomatic Quantum Unified‑Field Theory}

\vspace{2cm}

螺旋真空 · Frenet‑Serret · 广义 Kakeya · 引电统一 · Koide 关系 ·
洛伦兹时空拓展

\vspace{1.5cm}

约 1\,020\,000 字 · 全套 Python 高精度验证 · CODATA‑2022 对标

\clearpage

\section{全书总目录}\label{ux5168ux4e66ux603bux76eeux5f55}

\subsection{序言（Preface）}\label{ux5e8fux8a00preface}

\begin{itemize}
\tightlist
\item
  0.1 统一场论百年困境回顾
\item
  0.2 为什么从``几何基元''重新出发
\item
  0.3 GAQ‑UFT 核心思想简述：螺旋真空范式
\item
  0.4 本书阅读路线指引；数学预备知识说明
\item
  0.5 CODATA‑2022 实验对标说明；可证伪性立场
\end{itemize}

\subsection{第0卷 基础数学与公理地基（Volume\,00
Foundations，约14万字）}\label{ux7b2c0ux5377-ux57faux7840ux6570ux5b66ux4e0eux516cux7406ux5730ux57favolume-00-foundationsux7ea614ux4e07ux5b57}

\subsubsection{第1章 空间曲线微分几何：Frenet‑Serret
方程组}\label{ux7b2c1ux7ae0-ux7a7aux95f4ux66f2ux7ebfux5faeux5206ux51e0ux4f55frenetserret-ux65b9ux7a0bux7ec4}

1.1 正则空间曲线、弧长参数化完整定义 1.2 活动标架 \{T,N,B\}
严格推导；曲率 κ、挠率 τ 几何含义 1.3 Frenet‑Serret 微分方程组完整求导
1.4 常数曲率‑常数挠率：圆柱螺旋解析解推导 1.5
变曲率变挠率曲线的解空间；无穷维函数空间分析 1.6 mpmath
高精度数值微分、标架正交性校验方法论

\subsubsection{第2章 广义 Kakeya
集理论}\label{ux7b2c2ux7ae0-ux5e7fux4e49-kakeya-ux96c6ux7406ux8bba}

2.1 经典 Kakeya 集合简述（调和分析背景） 2.2 GAQ‑UFT
修改版：由光滑扭曲曲线生成的广义 Kakeya 集 2.3
切向球面覆盖条件、Lebesgue 测度、Hausdorff 维数 2.4
数学必要条件：\(\dim_H(K)=3\) 约束对生成曲线的限制

\subsubsection{第3章 拓扑学基础：庞加莱‑霍夫（毛球）定理完整证明}\label{ux7b2c3ux7ae0-ux62d3ux6251ux5b66ux57faux7840ux5e9eux52a0ux83b1ux970dux592bux6bdbux7403ux5b9aux7406ux5b8cux6574ux8bc1ux660e}

3.1 \(S^2\) 上处处非零切向量场不存在严格证明 3.2
向量场、零点指数、欧拉示性数 3.3
拓扑障碍如何限制真空模型；直线射线模型的内在矛盾

\subsubsection{第4章 GAQ‑UFT 全套公理系统
A1‑A7}\label{ux7b2c4ux7ae0-gaquft-ux5168ux5957ux516cux7406ux7cfbux7edf-a1a7}

4.1 A1：广义 Kakeya 覆盖公理 4.2 A2：非退化连续 Frenet 活动标架公理 4.3
A3：真空流线齐性公理（物理假设，区分数学/物理） 4.4
A4：流线相交、缠绕、打结拓扑公理 4.5 A5：量纲自洽公理 4.6
A6：洛伦兹协变性公理（3+1 维拓展） 4.7 A7：量子对应公理 4.8
公理之间的逻辑蕴含关系；哪些是必要，哪些是充分；逻辑真值表

\subsection{第1卷 螺旋真空：三维欧氏空间的统一场基元（Volume\,01
Geometric
Vacuum，约18万字）}\label{ux7b2c1ux5377-ux87baux65cbux771fux7a7aux4e09ux7ef4ux6b27ux6c0fux7a7aux95f4ux7684ux7edfux4e00ux573aux57faux5143volume-01-geometric-vacuumux7ea618ux4e07ux5b57}

\subsubsection{第5章 拓扑障碍对真空生成元的约束}\label{ux7b2c5ux7ae0-ux62d3ux6251ux969cux788dux5bf9ux771fux7a7aux751fux6210ux5143ux7684ux7ea6ux675f}

5.1 定理1完整证明：\(\kappa(s)>0,\ \tau(s)\neq0\) 必要条件 5.2
排除直线生成元的完整论证 5.3 排除平面曲线生成元完整论证 5.4
关键区分：A1+A2 只能推出必要条件，\textbf{不能推出 κ、τ 为常数}

\subsubsection{第6章 解空间四大分支全维分析}\label{ux7b2c6ux7ae0-ux89e3ux7a7aux95f4ux56dbux5927ux5206ux652fux5168ux7ef4ux5206ux6790}

6.1 分支0：直线射线（拓扑禁止） 6.2 分支1：平面曲线（几何禁止） 6.3
分支2：变曲率变挠率扭曲曲线（数学允许，物理代价） 6.4 分支3：常数 κ、τ
圆柱螺旋（A1+A2+A3，本理论工作分支） 6.5 反例曲线严格构造；200
位精度数值验证；表格与校验 6.6
如果宇宙落在分支2会发生什么：基本常数随空间时间漂移

\subsubsection{第7章 圆柱螺旋作为真空基元}\label{ux7b2c7ux7ae0-ux5706ux67f1ux87baux65cbux4f5cux4e3aux771fux7a7aux57faux5143}

7.1 Frenet‑Serret 下圆柱螺旋全部解析表达式 7.2
\(\tau/\kappa=\tan\theta\) 几何角度定义 7.3 螺旋局部几何不变量集合 7.4
螺旋系综：无穷多螺旋铺满空间，构成广义 Kakeya 真空

\subsubsection{第8章 真空流线的纠缠、相交、纽结理论初步}\label{ux7b2c8ux7ae0-ux771fux7a7aux6d41ux7ebfux7684ux7ea0ux7f20ux76f8ux4ea4ux7ebdux7ed3ux7406ux8bbaux521dux6b65}

8.1 两条螺旋相交拓扑；环绕数、链环数 8.2 纽结作为局域几何激发 8.3
真空背景 vs 局域纽结激发的本体论区分

\subsection{第2卷 基础物理常数的几何本源（Volume\,02 Constant Geometry
Origins，约17万字）}\label{ux7b2c2ux5377-ux57faux7840ux7269ux7406ux5e38ux6570ux7684ux51e0ux4f55ux672cux6e90volume-02-constant-geometry-originsux7ea617ux4e07ux5b57}

\subsubsection{第9章 精细结构常数 α
的几何推导}\label{ux7b2c9ux7ae0-ux7cbeux7ec6ux7ed3ux6784ux5e38ux6570-ux3b1-ux7684ux51e0ux4f55ux63a8ux5bfc}

9.1 \(\alpha \propto \tau/\kappa=\tan\theta\) 物理映射假设 9.2
角度‑常数映射；无量纲化处理 9.3 和 CODATA‑2022 α 数值对标 9.4 如果 A3
放松，α 漂移的数学形式 9.5 误差来源分析；理论不确定性边界

\subsubsection{\texorpdfstring{第10章 引电统一方程
\(G\varepsilon_0 \propto (\kappa_e/\kappa_\Omega)^2\)
完整推导}{第10章 引电统一方程 G\textbackslash varepsilon\_0 \textbackslash propto (\textbackslash kappa\_e/\textbackslash kappa\_\textbackslash Omega)\^{}2 完整推导}}\label{ux7b2c10ux7ae0-ux5f15ux7535ux7edfux4e00ux65b9ux7a0b-gvarepsilon_0-propto-kappa_ekappa_omega2-ux5b8cux6574ux63a8ux5bfc}

10.1 量纲一致性完整校验；左右两侧量纲逐项核对 10.2 电磁曲率
\(\kappa_e\)、宇宙真空背景曲率 \(\kappa_\Omega\) 定义 10.3 引力常数 G
不是单根螺旋局部量，而是系综集体涌现常数 10.4 G
的表达式中系综平均、缠绕密度因子 10.5 CODATA‑2022 G
数值对标；理论偏差讨论 10.6 引力与电磁力几何根源对比表格

\subsection{第3卷 粒子谱：轻子、夸克几何起源（Volume\,03 Particle
Spectra，约16万字）}\label{ux7b2c3ux5377-ux7c92ux5b50ux8c31ux8f7bux5b50ux5938ux514bux51e0ux4f55ux8d77ux6e90volume-03-particle-spectraux7ea616ux4e07ux5b57}

\subsubsection{第11章 Koide
质量关系深度剖析}\label{ux7b2c11ux7ae0-koide-ux8d28ux91cfux5173ux7cfbux6df1ux5ea6ux5256ux6790}

11.1 Koide 经验公式回顾，历史文献梳理 11.2 Koide
关系的代数结构；自由度计数 11.3
螺旋纽结、量子化缠绕数与轻子质量本征值映射 11.4
完美真空螺旋（A3）的微小对称破缺对应粒子质量分裂 11.5 数值扫描；Koide
公式理论修正项预测 11.6 拓展向夸克质量谱系；遇到的困难与开放点

\subsubsection{第12章 轻子作为真空螺旋的拓扑激发}\label{ux7b2c12ux7ae0-ux8f7bux5b50ux4f5cux4e3aux771fux7a7aux87baux65cbux7684ux62d3ux6251ux6fc0ux53d1}

12.1 电子、μ子、τ子对应不同纽结/缠绕组态 12.2 电荷作为流线泄露的几何图像
12.3 自旋与螺旋内禀旋转几何关系 12.4
粒子波函数与德布罗意物质波：从螺旋运动直接导出德布罗意关系

\subsection{第4卷 洛伦兹时空拓展：3+1维相对论统一场（Volume\,04 Lorentz
Spacetime
Extension，约15万字）}\label{ux7b2c4ux5377-ux6d1bux4f26ux5179ux65f6ux7a7aux62d3ux5c5531ux7ef4ux76f8ux5bf9ux8bbaux7edfux4e00ux573avolume-04-lorentz-spacetime-extensionux7ea615ux4e07ux5b57}

\subsubsection{第13章 洛伦兹号差下 Frenet‑Serret
活动标架}\label{ux7b2c13ux7ae0-ux6d1bux4f26ux5179ux53f7ux5deeux4e0b-frenetserret-ux6d3bux52a8ux6807ux67b6}

13.1 类时、类光、类空世界线的 Frenet‑Serret 推广 13.2 洛伦兹 boost
下螺旋曲率、挠率变换规则 13.3 时间膨胀作为螺旋几何推论完整推导

\subsubsection{第14章 广义 Kakeya
向光锥方向空间拓展}\label{ux7b2c14ux7ae0-ux5e7fux4e49-kakeya-ux5411ux5149ux9525ux65b9ux5411ux7a7aux95f4ux62d3ux5c55}

14.1 光锥上切向覆盖条件；类光生成元 14.2
洛伦兹版本拓扑障碍；新的奇点类型 14.3
时空度规如何从螺旋系综集体行为涌现出来

\subsection{第5卷 可观测预言与可证伪通道（Volume\,05 Prediction
Falsification，约12万字）}\label{ux7b2c5ux5377-ux53efux89c2ux6d4bux9884ux8a00ux4e0eux53efux8bc1ux4f2aux901aux9053volume-05-prediction-falsificationux7ea612ux4e07ux5b57}

\subsubsection{第15章 全套物理预言清单}\label{ux7b2c15ux7ae0-ux5168ux5957ux7269ux7406ux9884ux8a00ux6e05ux5355}

15.1 精细结构常数宇宙学漂移预言 15.2 真空色散：高能光子传播速度依赖能量
15.3 引力波偏振微小异常预言 15.4 高能宇宙射线各向异性信号 15.5 Koide
公式高精度偏离预言

\subsubsection{第16章 现有天文、粒子实验约束}\label{ux7b2c16ux7ae0-ux73b0ux6709ux5929ux6587ux7c92ux5b50ux5b9eux9a8cux7ea6ux675f}

16.1 类星体光谱对 α 漂移的观测边界 16.2 GRB 伽马暴对真空色散的限制 16.3
LIGO 引力波数据约束 16.4 实验如何证伪或者约束 GAQ‑UFT 理论

\subsection{第6卷 开放问题与未来研究（Volume\,06 Open Problems
Future，约10万字）}\label{ux7b2c6ux5377-ux5f00ux653eux95eeux9898ux4e0eux672aux6765ux7814ux7a76volume-06-open-problems-futureux7ea610ux4e07ux5b57}

\subsubsection{第17章 开放研究问题全集}\label{ux7b2c17ux7ae0-ux5f00ux653eux7814ux7a76ux95eeux9898ux5168ux96c6}

17.1 G 常数更严格的完整闭合解析推导 17.2 完整动力学场微分方程组构建 17.3
量子引力路径积分的几何版本 17.4
暗物质、暗能量在螺旋真空理论框架下解释可能性 17.5 量子测量问题几何视角
17.6 分支2（变 κτ 宇宙）进一步研究价值

\subsubsection{第18章 总结：GAQ‑UFT
完整逻辑闭环；理论局限与边界}\label{ux7b2c18ux7ae0-ux603bux7ed3gaquft-ux5b8cux6574ux903bux8f91ux95edux73afux7406ux8bbaux5c40ux9650ux4e0eux8fb9ux754c}

\subsection{附录A 全部 200 位精度 Python
仿真代码（书籍内附录印刷版）}\label{ux9644ux5f55a-ux5168ux90e8-200-ux4f4dux7cbeux5ea6-python-ux4effux771fux4ee3ux7801ux4e66ux7c4dux5185ux9644ux5f55ux5370ux5237ux7248}

\subsection{附录B 重要定理完整证明备查}\label{ux9644ux5f55b-ux91cdux8981ux5b9aux7406ux5b8cux6574ux8bc1ux660eux5907ux67e5}

\subsection{附录C CODATA‑2022
常数数据表}\label{ux9644ux5f55c-codata2022-ux5e38ux6570ux6570ux636eux8868}

\subsection{参考文献
references.bib}\label{ux53c2ux8003ux6587ux732e-references.bib}

\section{序言：统一场论的几何重生}\label{ux5e8fux8a00ux7edfux4e00ux573aux8bbaux7684ux51e0ux4f55ux91cdux751f}

一百多年以来，物理学怀揣一个古老梦想：找到一套底层原理，把引力、电磁相互作用乃至微观粒子的全部性质，从同一套基础公理推导出来。爱因斯坦后半生致力于统一场论，尝试用几何改造时空，却没有完成；后来标准模型取得巨大成功，但标准模型拥有二十多个自由输入参数，无法从第一性原理算出基本常数；量子引力至今仍是悬而未决的核心难题。

现有主流路径大致分为弦论、圈量子引力等方向。弦论引入高维空间、一维弦作为基元，但大量真空解带来景观问题；圈量子引力对时空做自旋网络离散化。本书
GAQ‑UFT（Geometric‑Axiomatic Quantum Unified‑Field
Theory，万物几何统一场论）选择一条不同路径：\textbf{不预先假设离散格子，也不引入额外空间维度；把三维空间的扭曲螺旋流线作为时空的本源基元。}

本书的出发点来自微分几何的 Frenet‑Serret 活动标架，结合经过修改的广义
Kakeya
集合条件，再利用庞加莱‑霍夫拓扑障碍定理做强约束。全书最关键的认识论分界是严格区分两件事物：哪些结论是纯粹数学拓扑给出的\textbf{必要约束}；哪些是来自现实观测、人为添加的\textbf{物理公设}。

数学本身只能告诉我们：真空基础生成元不能是直线，不能是平面曲线，必须是处处同时拥有曲率与挠率的三维扭曲空间曲线。数学并不强制曲率、挠率必须沿流线保持不变。``流线齐性''（A3
公理）是物理假设，不是数学定理。一旦去掉这条假设，我们会得到无穷多数学自洽的宇宙解，但那些宇宙之中，精细结构常数会随着位置、时间发生改变，和我们真实宇宙观测结果冲突。

当我们增加真空流线齐性公理，解被锁定为圆柱螺旋。螺旋的挠率‑曲率比值
\(\tau/\kappa=\tan\theta\)，给出精细结构常数的几何源头。引力常数 \(G\)
不再是单根螺旋的局部几何量，而是无穷螺旋系综互相缠绕打结的集体涌现效应，由此得到引电统一几何关系。在此基础之上，进一步尝试对接
Koide 轻子质量经验关系，把粒子解释为真空螺旋背景上的拓扑纽结激发。

本书不是已经完成的终极真理，而是一套完整可推导、可计算、\textbf{可证伪}的理论框架。书中会明确标出哪些部分是严格证明、哪些是猜想假设、哪些属于尚未完全闭合的开放研究问题。书中给出全部
Python 高精度数值验证代码，所有关键推导都可以在计算机上复现，对标
CODATA‑2022 物理常数实验数据。

阅读本书建议：如果你熟悉微分几何，可以直接从第 0 卷第 1
章开始；如果你希望先理解物理图像，可以先读第 1
卷，遇到公式再回头查阅数学附录。全书大量区分``必要条件／充分条件''，这是理解本理论最核心的钥匙。

本理论有明确的实验判决通道：如果未来观测到精细结构常数存在宇宙学漂移，就意味着
A3
齐性公理失效，宇宙落在变曲率‑变挠率的解分支；如果所有高精度测量显示基础常数严格不变，则本理论的工作分支获得更强支撑。

\subsection{本书结构导读}\label{ux672cux4e66ux7ed3ux6784ux5bfcux8bfb}

全书按''数学地基 → 几何真空 → 常数起源 → 粒子谱 → 时空延展 → 预言证伪 →
开放问题''七个递进卷册组织，外加序言、三份附录与参考文献：

\begin{itemize}
\tightlist
\item
  \textbf{第 0 卷 数学公理地基}：从 Frenet--Serret
  活动标架（第1章）出发，建立公理系统（第4章），用广义 Kakeya
  覆盖（第2章）与庞加莱‑霍夫拓扑障碍（第3章）排除直线/平面生成元。这一卷纯数学，不含任何物理常数。
\item
  \textbf{第 1 卷
  几何真空与拓扑激发}：把公理结晶为''螺旋真空系综''（第7章），证明 A1+A2
  推出扭曲曲线的拓扑障碍定理（第5章），划分解空间四大分支（第6章），并把粒子识别为真空纽结（第8章）。
\item
  \textbf{第 2 卷 常数几何起源}：以螺旋升角 \(\theta=\arctan\alpha\)
  锚定精细结构常数（第9章），给出引力常数 \(G\)
  的系综涌现框架与引电统一方程（第10章）。
\item
  \textbf{第 3 卷
  粒子谱与质量几何}：把轻子质量映射到螺旋纽结的量子化缠绕数，对接 Koide
  经验关系（第11章），并把粒子具体化为拓扑激发、导出德布罗意关系（第12章）。
\item
  \textbf{第 4 卷
  洛伦兹与时空延展}：把三维螺旋标架提升到四维时空洛伦兹标架（第13章），用光锥
  Kakeya 覆盖与系综平均实现度规涌现（第14章）。
\item
  \textbf{第 5 卷
  预言与证伪}：列出可检验预言（第15章）与实验约束现状（第16章），把可证伪性落到实处。
\item
  \textbf{第 6 卷
  开放问题与未来方向}：诚实列出未闭合项（第18章总结）与量子引力路径积分等研究方向（第17章）。
\end{itemize}

\textbf{读者路线建议}：微分几何熟悉者可从第 0
卷顺读；希望先建立物理图像者可从第 1
卷切入，遇公式再回头查数学附录；关心''理论是否被实验排除''者可直接跳读第
5 卷；关心''理论诚实边界''者必读第 6 卷与
\texttt{精算验证总结与核心方程清单.md}。

物理学的进步不在于构造完美无缺的故事，而在于构造可以被现实检验的逻辑体系。本书
GAQ‑UFT，正是这样一次从几何第一性原理出发，向统一场论推进的完整尝试。

\begin{center}\rule{0.5\linewidth}{0.5pt}\end{center}

\emph{（本文件为 \texttt{GAQ‑UFT\_Complete\_Book} 包内
\texttt{preface.md}。编译时由 \texttt{compile\_book} 脚本合并入主书籍
\texttt{main\_book.md}。）}

\section{附录 A 全部 200 位精度 Python
仿真代码（书籍内印刷版）}\label{ux9644ux5f55-a-ux5168ux90e8-200-ux4f4dux7cbeux5ea6-python-ux4effux771fux4ee3ux7801ux4e66ux7c4dux5185ux5370ux5237ux7248}

\begin{quote}
本附录收录 \texttt{Supplemental\_Code/}
下的全部脚本，供读者脱离电子设备直接在书中查阅。全部脚本依赖
\texttt{mpmath}（高精度）与 \texttt{numpy/matplotlib}（可视化）。
\end{quote}

\subsection{A.1 frenet\_high\_prec.py --- Frenet--Serret 200
位精度求解器}\label{a.1-frenet_high_prec.py-frenetserret-200-ux4f4dux7cbeux5ea6ux6c42ux89e3ux5668}

\begin{Shaded}
\begin{Highlighting}[]
\CommentTok{"""}
\CommentTok{GAQ‑UFT Supplemental: Frenet‑Serret 200‑digit precision solver}
\CommentTok{Generalized Kakeya Helical Vacuum Unified Field Theory}
\CommentTok{"""}
\ImportTok{import}\NormalTok{ mpmath }\ImportTok{as}\NormalTok{ mp}
\NormalTok{mp.mp.dps }\OperatorTok{=} \DecValTok{200}

\KeywordTok{def}\NormalTok{ diff\_central(f, t, h}\OperatorTok{=}\NormalTok{mp.mpf(}\StringTok{"1e{-}120"}\NormalTok{)):}
    \ControlFlowTok{return}\NormalTok{ (f(t}\OperatorTok{+}\NormalTok{h)}\OperatorTok{{-}}\NormalTok{f(t}\OperatorTok{{-}}\NormalTok{h))}\OperatorTok{/}\NormalTok{(}\DecValTok{2}\OperatorTok{*}\NormalTok{h)}

\KeywordTok{def}\NormalTok{ cross3(a,b):}
    \ControlFlowTok{return}\NormalTok{ mp.matrix([}
\NormalTok{        a[}\DecValTok{1}\NormalTok{]}\OperatorTok{*}\NormalTok{b[}\DecValTok{2}\NormalTok{]}\OperatorTok{{-}}\NormalTok{a[}\DecValTok{2}\NormalTok{]}\OperatorTok{*}\NormalTok{b[}\DecValTok{1}\NormalTok{],}
\NormalTok{        a[}\DecValTok{2}\NormalTok{]}\OperatorTok{*}\NormalTok{b[}\DecValTok{0}\NormalTok{]}\OperatorTok{{-}}\NormalTok{a[}\DecValTok{0}\NormalTok{]}\OperatorTok{*}\NormalTok{b[}\DecValTok{2}\NormalTok{],}
\NormalTok{        a[}\DecValTok{0}\NormalTok{]}\OperatorTok{*}\NormalTok{b[}\DecValTok{1}\NormalTok{]}\OperatorTok{{-}}\NormalTok{a[}\DecValTok{1}\NormalTok{]}\OperatorTok{*}\NormalTok{b[}\DecValTok{0}\NormalTok{]}
\NormalTok{    ])}

\KeywordTok{def}\NormalTok{ frenet\_decompose(curve\_func, t):}
\NormalTok{    r1 }\OperatorTok{=}\NormalTok{ diff\_central(curve\_func, t)}
\NormalTok{    r2 }\OperatorTok{=}\NormalTok{ diff\_central(}\KeywordTok{lambda}\NormalTok{ s: diff\_central(curve\_func,s), t)}
\NormalTok{    r3 }\OperatorTok{=}\NormalTok{ diff\_central(}\KeywordTok{lambda}\NormalTok{ s: diff\_central(}\KeywordTok{lambda}\NormalTok{ u: diff\_central(curve\_func,u),s), t)}
\NormalTok{    dr }\OperatorTok{=}\NormalTok{ mp.norm(r1)}
\NormalTok{    T }\OperatorTok{=}\NormalTok{ r1 }\OperatorTok{/}\NormalTok{ dr}
\NormalTok{    r1xr2 }\OperatorTok{=}\NormalTok{ cross3(r1,r2)}
\NormalTok{    cr\_norm }\OperatorTok{=}\NormalTok{ mp.norm(r1xr2)}
\NormalTok{    kappa }\OperatorTok{=}\NormalTok{ cr\_norm}\OperatorTok{/}\NormalTok{(dr}\OperatorTok{**}\DecValTok{3}\NormalTok{)}
\NormalTok{    N }\OperatorTok{=}\NormalTok{ cross3(r1, r1xr2)}\OperatorTok{/}\NormalTok{(dr}\OperatorTok{*}\NormalTok{cr\_norm)}
\NormalTok{    B }\OperatorTok{=}\NormalTok{ cross3(T,N)}
\NormalTok{    triple }\OperatorTok{=}\NormalTok{ (r1xr2.T }\OperatorTok{*}\NormalTok{ r3)[}\DecValTok{0}\NormalTok{,}\DecValTok{0}\NormalTok{]}
\NormalTok{    tau }\OperatorTok{=}\NormalTok{ triple}\OperatorTok{/}\NormalTok{(cr\_norm}\OperatorTok{**}\DecValTok{2}\NormalTok{)}
    \ControlFlowTok{return}\NormalTok{ T,N,B,kappa,tau}

\KeywordTok{def}\NormalTok{ gamma\_var(t):}
\NormalTok{    Rt }\OperatorTok{=}\NormalTok{ mp.mpf(}\StringTok{"2.0"}\NormalTok{) }\OperatorTok{+}\NormalTok{ mp.mpf(}\StringTok{"0.2"}\NormalTok{)}\OperatorTok{*}\NormalTok{mp.sin(t)}
\NormalTok{    x }\OperatorTok{=}\NormalTok{ Rt}\OperatorTok{*}\NormalTok{mp.cos(t)}\OperatorTok{;}\NormalTok{ y }\OperatorTok{=}\NormalTok{ Rt}\OperatorTok{*}\NormalTok{mp.sin(t)}\OperatorTok{;}\NormalTok{ z }\OperatorTok{=}\NormalTok{ mp.mpf(}\StringTok{"0.8"}\NormalTok{)}\OperatorTok{*}\NormalTok{t }\OperatorTok{+}\NormalTok{ mp.mpf(}\StringTok{"0.1"}\NormalTok{)}\OperatorTok{*}\NormalTok{mp.cos(t)}
    \ControlFlowTok{return}\NormalTok{ mp.matrix([x,y,z])}

\KeywordTok{def}\NormalTok{ helix\_const(t,R}\OperatorTok{=}\DecValTok{1}\NormalTok{,omega}\OperatorTok{=}\DecValTok{1}\NormalTok{,c}\OperatorTok{=}\DecValTok{1}\NormalTok{):}
    \ControlFlowTok{return}\NormalTok{ mp.matrix([R}\OperatorTok{*}\NormalTok{mp.cos(omega}\OperatorTok{*}\NormalTok{t), R}\OperatorTok{*}\NormalTok{mp.sin(omega}\OperatorTok{*}\NormalTok{t), c}\OperatorTok{*}\NormalTok{t])}

\ControlFlowTok{if} \VariableTok{\_\_name\_\_}\OperatorTok{==}\StringTok{"\_\_main\_\_"}\NormalTok{:}
    \ControlFlowTok{for}\NormalTok{ ti }\KeywordTok{in}\NormalTok{ [mp.mpf(}\StringTok{"0"}\NormalTok{), mp.pi}\OperatorTok{/}\DecValTok{3}\NormalTok{, mp.pi}\OperatorTok{/}\DecValTok{2}\NormalTok{, mp.pi]:}
\NormalTok{        T,N,B,kap,tau }\OperatorTok{=}\NormalTok{ frenet\_decompose(gamma\_var, ti)}
        \BuiltInTok{print}\NormalTok{(}\SpecialStringTok{f"t=}\SpecialCharTok{\{}\NormalTok{ti}\SpecialCharTok{:.4f\}}\SpecialStringTok{  κ=}\SpecialCharTok{\{}\NormalTok{kap}\SpecialCharTok{:.14f\}}\SpecialStringTok{  τ=}\SpecialCharTok{\{}\NormalTok{tau}\SpecialCharTok{:.14f\}}\SpecialStringTok{  τ/κ=}\SpecialCharTok{\{}\NormalTok{tau}\OperatorTok{/}\NormalTok{kap}\SpecialCharTok{:.14f\}}\SpecialStringTok{"}\NormalTok{)}
\end{Highlighting}
\end{Shaded}

\subsection{A.2 geometry\_verify.py ---
引电统一量纲与数值校验}\label{a.2-geometry_verify.py-ux5f15ux7535ux7edfux4e00ux91cfux7eb2ux4e0eux6570ux503cux6821ux9a8c}

\begin{Shaded}
\begin{Highlighting}[]
\ImportTok{import}\NormalTok{ mpmath }\ImportTok{as}\NormalTok{ mp}
\NormalTok{mp.mp.dps}\OperatorTok{=}\DecValTok{120}
\NormalTok{G }\OperatorTok{=}\NormalTok{ mp.mpf(}\StringTok{"6.67430e{-}11"}\NormalTok{)}\OperatorTok{;}\NormalTok{ eps0 }\OperatorTok{=}\NormalTok{ mp.mpf(}\StringTok{"8.8541878128e{-}12"}\NormalTok{)}\OperatorTok{;}\NormalTok{ c }\OperatorTok{=}\NormalTok{ mp.mpf(}\StringTok{"299792458"}\NormalTok{)}
\KeywordTok{def}\NormalTok{ dim\_check\_Geps0():}
\NormalTok{    Geps0 }\OperatorTok{=}\NormalTok{ G}\OperatorTok{*}\NormalTok{eps0}
    \BuiltInTok{print}\NormalTok{(}\SpecialStringTok{f"G*ε0 = }\SpecialCharTok{\{}\NormalTok{Geps0}\SpecialCharTok{:.20e\}}\SpecialStringTok{"}\NormalTok{)}\OperatorTok{;} \ControlFlowTok{return}\NormalTok{ Geps0}
\ControlFlowTok{if} \VariableTok{\_\_name\_\_}\OperatorTok{==}\StringTok{"\_\_main\_\_"}\NormalTok{:}
\NormalTok{    dim\_check\_Geps0()}
\end{Highlighting}
\end{Shaded}

\subsection{A.3 koide\_simulation.py --- Koide
关系数值扫描}\label{a.3-koide_simulation.py-koide-ux5173ux7cfbux6570ux503cux626bux63cf}

\begin{Shaded}
\begin{Highlighting}[]
\ImportTok{import}\NormalTok{ mpmath }\ImportTok{as}\NormalTok{ mp}
\NormalTok{mp.mp.dps }\OperatorTok{=} \DecValTok{100}
\NormalTok{me }\OperatorTok{=}\NormalTok{ mp.mpf(}\StringTok{"0.51099895000"}\NormalTok{)}\OperatorTok{;}\NormalTok{ mmu }\OperatorTok{=}\NormalTok{ mp.mpf(}\StringTok{"105.6583755"}\NormalTok{)}\OperatorTok{;}\NormalTok{ mtau }\OperatorTok{=}\NormalTok{ mp.mpf(}\StringTok{"1776.86"}\NormalTok{)}
\KeywordTok{def}\NormalTok{ koide\_formula(m1,m2,m3):}
\NormalTok{    lhs }\OperatorTok{=}\NormalTok{ m1}\OperatorTok{+}\NormalTok{m2}\OperatorTok{+}\NormalTok{m3}
\NormalTok{    rhs }\OperatorTok{=}\NormalTok{ (}\DecValTok{2}\OperatorTok{/}\DecValTok{3}\NormalTok{)}\OperatorTok{*}\NormalTok{(mp.sqrt(m1)}\OperatorTok{+}\NormalTok{mp.sqrt(m2)}\OperatorTok{+}\NormalTok{mp.sqrt(m3))}\OperatorTok{**}\DecValTok{2}
    \ControlFlowTok{return}\NormalTok{ lhs, rhs, }\BuiltInTok{abs}\NormalTok{(lhs}\OperatorTok{{-}}\NormalTok{rhs)}\OperatorTok{/}\NormalTok{lhs}
\ControlFlowTok{if} \VariableTok{\_\_name\_\_}\OperatorTok{==}\StringTok{"\_\_main\_\_"}\NormalTok{:}
\NormalTok{    L,R,err }\OperatorTok{=}\NormalTok{ koide\_formula(me,mmu,mtau)}
    \BuiltInTok{print}\NormalTok{(}\SpecialStringTok{f"LHS = }\SpecialCharTok{\{}\NormalTok{L}\SpecialCharTok{:.8f\}}\CharTok{\textbackslash{}n}\SpecialStringTok{RHS = }\SpecialCharTok{\{}\NormalTok{R}\SpecialCharTok{:.8f\}}\CharTok{\textbackslash{}n}\SpecialStringTok{relative error = }\SpecialCharTok{\{}\NormalTok{err}\SpecialCharTok{:.8e\}}\SpecialStringTok{"}\NormalTok{)}
\end{Highlighting}
\end{Shaded}

\subsection{A.4 plot\_helix\_3d.py --- 3D
螺旋可视化}\label{a.4-plot_helix_3d.py-3d-ux87baux65cbux53efux89c6ux5316}

\begin{Shaded}
\begin{Highlighting}[]
\ImportTok{import}\NormalTok{ numpy }\ImportTok{as}\NormalTok{ np}
\ImportTok{import}\NormalTok{ matplotlib.pyplot }\ImportTok{as}\NormalTok{ plt}
\ImportTok{from}\NormalTok{ mpl\_toolkits.mplot3d }\ImportTok{import}\NormalTok{ Axes3D}
\NormalTok{t }\OperatorTok{=}\NormalTok{ np.linspace(}\DecValTok{0}\NormalTok{, }\DecValTok{12}\OperatorTok{*}\NormalTok{np.pi, }\DecValTok{2000}\NormalTok{)}
\NormalTok{x }\OperatorTok{=}\NormalTok{ np.cos(t)}\OperatorTok{;}\NormalTok{ y }\OperatorTok{=}\NormalTok{ np.sin(t)}\OperatorTok{;}\NormalTok{ z }\OperatorTok{=}\NormalTok{ t}
\NormalTok{fig }\OperatorTok{=}\NormalTok{ plt.figure(figsize}\OperatorTok{=}\NormalTok{(}\DecValTok{10}\NormalTok{,}\DecValTok{7}\NormalTok{))}
\NormalTok{ax }\OperatorTok{=}\NormalTok{ fig.add\_subplot(}\DecValTok{111}\NormalTok{,projection}\OperatorTok{=}\StringTok{\textquotesingle{}3d\textquotesingle{}}\NormalTok{)}
\NormalTok{ax.plot(x,y,z,lw}\OperatorTok{=}\FloatTok{0.8}\NormalTok{)}
\NormalTok{ax.set\_title(}\StringTok{"Cylindrical circular helix — GAQ‑UFT vacuum generator"}\NormalTok{)}
\NormalTok{plt.show()}
\end{Highlighting}
\end{Shaded}

\begin{center}\rule{0.5\linewidth}{0.5pt}\end{center}

\emph{（附录 A 完。完整可运行版本见
\texttt{Supplemental\_Code/}，与本书同源维护。）}

\section{附录 B
重要定理完整证明备查}\label{ux9644ux5f55-b-ux91cdux8981ux5b9aux7406ux5b8cux6574ux8bc1ux660eux5907ux67e5}

\begin{quote}
本书核心定理的集中存档，便于快速查阅与引用。正文推导见对应章节。
\end{quote}

\begin{center}\rule{0.5\linewidth}{0.5pt}\end{center}

\subsection{B.1
定理1（生成元扭曲必要条件，第5章）}\label{b.1-ux5b9aux74061ux751fux6210ux5143ux626dux66f2ux5fc5ux8981ux6761ux4ef6ux7b2c5ux7ae0}

\textbf{陈述：} A1+A2 \(\Rightarrow\) 生成元满足 \(\kappa(s)>0\) 且
\(\tau(s)\neq0\)。

\textbf{证明：} （1）若 \(\kappa=0\) 则 \(\mathbf{T}\)
恒定，曲线为直线，切向仅一个方向，违反 A1 全覆盖，且违反 A2
非退化。（2）若 \(\tau=0\) 则 \(\mathbf{B}\)
恒定，曲线为平面曲线，切向仅落在一个大圆，违反 A1。故二者必须处处成立。□

\subsection{B.2
曲线基本定理（第1.5.1）}\label{b.2-ux66f2ux7ebfux57faux672cux5b9aux7406ux7b2c1.5.1}

\textbf{陈述：} 给定光滑 \(\kappa(s)>0,\tau(s)\)
与标准正交初值标架，存在唯一（差欧氏运动）正则空间曲线以之为曲率挠率。

\textbf{证明概要：} Frenet--Serret 是线性 ODE
系统，初值唯一决定标架演化；积分
\(\mathbf{r}(s)=\int\mathbf{T}(s)ds\)+初值即得曲线，唯一性由线性系统初值问题唯一性保证。□

\subsection{B.3
圆柱螺旋充分必要条件（第1.4）}\label{b.3-ux5706ux67f1ux87baux65cbux5145ux5206ux5fc5ux8981ux6761ux4ef6ux7b2c1.4}

\textbf{陈述：} 常数 \(\kappa,\tau\)
当且仅当曲线为圆柱螺旋（差刚体运动）。

\textbf{证明：} 由 B.2，常数 \(\kappa_0,\tau_0\) 唯一确定一条曲线；1.4
节直接积分给出螺旋参数形式，验证其曲率挠率恰为所给常数。□

\subsection{B.4
毛球定理（第3.1.1）}\label{b.4-ux6bdbux7403ux5b9aux7406ux7b2c3.1.1}

\textbf{陈述：} \(S^2\) 上无处处非零连续切向量场。

\textbf{证明：}
见第3.1节度数论证。核心：归一化切场的度数在同伦下不变，而对径提升使度数反号，唯零可满足，矛盾。□

\subsection{B.5
引电统一量纲结论（第10.7）}\label{b.5-ux5f15ux7535ux7edfux4e00ux91cfux7eb2ux7ed3ux8bbaux7b2c10.7}

\textbf{陈述：} \(G\varepsilon_0\) 量纲为 \(Q^2M^{-2}\)，而
\((\kappa_e/\kappa_\Omega)^2\) 无量纲，故等式需带量纲桥接常数
\(K_{\mathrm{bridge}}\)（量纲 \(Q^2M^{-2}\)）。

\textbf{证明：} 逐项量纲代入，见 10.7 表。□

\subsection{B.6 验证套件 T
系列与定理对应关系}\label{b.6-ux9a8cux8bc1ux5957ux4ef6-t-ux7cfbux5217ux4e0eux5b9aux7406ux5bf9ux5e94ux5173ux7cfb}

全书数值验证由 \texttt{verification\_suite\_v2.py}
承载，其检验项与本书定理/章节的对应如下（状态：7 PASS / 0 FAIL / 2 OPEN
/ 1 INFO）：

\begin{longtable}[]{@{}llll@{}}
\toprule\noalign{}
检验项 & 验证对象 & 对应定理/章节 & 状态 \\
\midrule\noalign{}
\endhead
\bottomrule\noalign{}
\endlastfoot
T1 & Frenet 标架正交性 & B.2 曲线基本定理 & PASS \\
T2 & 圆柱螺旋 \(\kappa=\tau=1/2\) & B.3 & PASS \\
T3 & 升角 \(\theta\to0\) 退化检查 & 第9章 & PASS \\
T4 & \(\theta=\arctan\alpha\) 回代 & 第9章 M‑α 映射 & PASS（误差 0） \\
T6 & Koide 相对偏差 \(9.2\times10^{-6}\) & 第11章 & PASS \\
T9 & 毛球定理指数和 \(=2\) & B.4 / 第3章 & PASS \\
T10 & 解空间四大分支自动归类 & B.1 / 第6章 & PASS \\
T11 & 引电统一比值 & B.5 / 第10.7 & OPEN（桥接常数未闭合） \\
T12 & \(\alpha\) 漂移观测阈值 & 第15章 & INFO（观测约束） \\
T13 & \(G\) 数值闭环 & 第10.2 / 17.1 & OPEN（\(\mathcal{F}\) 未闭合） \\
\end{longtable}

凡标注 OPEN / INFO
的项，其对应章节正文均如实标记为''未闭合''------验证套件与正文保持同一套诚实边界，不夸大已证内容。

\begin{center}\rule{0.5\linewidth}{0.5pt}\end{center}

\emph{（附录 B 完。CODATA 数据见附录 C。）}

\section{附录 C CODATA‑2022
常数数据表}\label{ux9644ux5f55-c-codata2022-ux5e38ux6570ux6570ux636eux8868}

\begin{quote}
本书所有数值对标采用 CODATA‑2022 推荐值（NIST SP 961，2024
发布）。下表摘录关键常数。
\end{quote}

\begin{longtable}[]{@{}
  >{\raggedright\arraybackslash}p{(\linewidth - 6\tabcolsep) * \real{0.1538}}
  >{\raggedright\arraybackslash}p{(\linewidth - 6\tabcolsep) * \real{0.1538}}
  >{\raggedright\arraybackslash}p{(\linewidth - 6\tabcolsep) * \real{0.5385}}
  >{\raggedright\arraybackslash}p{(\linewidth - 6\tabcolsep) * \real{0.1538}}@{}}
\toprule\noalign{}
\begin{minipage}[b]{\linewidth}\raggedright
常数
\end{minipage} & \begin{minipage}[b]{\linewidth}\raggedright
符号
\end{minipage} & \begin{minipage}[b]{\linewidth}\raggedright
数值（CODATA‑2022）
\end{minipage} & \begin{minipage}[b]{\linewidth}\raggedright
单位
\end{minipage} \\
\midrule\noalign{}
\endhead
\bottomrule\noalign{}
\endlastfoot
真空中光速 & \(c\) & 299\,792\,458（精确） & m·s⁻¹ \\
真空磁导率 & \(\mu_0\) & \(1.256\,637\,362\,12(19)\times10^{-6}\) &
N·A⁻² \\
真空电容率 & \(\varepsilon_0\) & \(8.854\,187\,8128(13)\times10^{-12}\)
& F·m⁻¹ \\
牛顿引力常数 & \(G\) & \(6.674\,30(15)\times10^{-11}\) & m³·kg⁻¹·s⁻² \\
普朗克常数 & \(h\) & \(6.626\,070\,15(0)\times10^{-34}\)（精确） &
J·s \\
约化普朗克常数 & \(\hbar\) & \(1.054\,571\,817\ldots\times10^{-34}\) &
J·s \\
电子质量 & \(m_e\) & 0.510\,998\,950\,00(0) & MeV/c² \\
μ子质量 & \(m_\mu\) & 105.658\,3755(23) & MeV/c² \\
τ子质量 & \(m_\tau\) & 1776.86(12) & MeV/c² \\
基本电荷 & \(e\) & \(1.602\,176\,634(0)\times10^{-19}\)（精确） & C \\
精细结构常数 & \(\alpha\) & \(7.297\,352\,5643(11)\times10^{-3}\) &
无量纲 \\
精细结构常数倒数 & \(\alpha^{-1}\) & 137.035\,999\,084(21) & 无量纲 \\
\end{longtable}

\textbf{使用说明：} 运行 \texttt{Supplemental\_Code/geometry\_verify.py}
与 \texttt{koide\_simulation.py} 时，上述数值已硬编码入脚本。若未来
CODATA 更新，请同步修改脚本常量并重新校验全书对标结论。

\subsection{C.1
常数误差传播与对标稳健性}\label{c.1-ux5e38ux6570ux8befux5deeux4f20ux64adux4e0eux5bf9ux6807ux7a33ux5065ux6027}

本书对标结论对 CODATA 不确定度具备不同敏感度的分层：

\begin{itemize}
\tightlist
\item
  \textbf{零敏感（精确常数）}：\(c\)、\(h\)、\(e\)、\(\mu_0\)（由
  \(c,e,h\) 定义）为 SI 定义精确值，其参与的对标（如
  \(\theta=\arctan\alpha\) 的回代）不引入实验误差，故 T4 误差为 0
  是结构性的而非巧合。
\item
  \textbf{低风险敏感}：\(m_e,m_\mu\)
  相对不确定度极小（\(\lesssim10^{-8}\)），Koide 相对偏差
  \(9.2\times10^{-6}\)（T6）远高于其统计误差，结论稳健。
\item
  \textbf{中风险敏感}：\(m_\tau\) 相对不确定度约
  \(7\times10^{-5}\)，仍小于 T6 偏差三个量级，对''Koide
  成立''结论无害，但影响 \(\phi\) 映射的精修（11.3.1）。
\item
  \textbf{高风险敏感}：\(G\) 的相对不确定度高达
  \(2.2\times10^{-5}\)（CODATA‑2022），是全书所有对标中\textbf{最弱的一环}；这直接导致
  T13（\(G\) 数值闭环）无法在当前精度下作出强判决，也提醒读者：任何涉及
  \(G\) 的''精确推导''主张都应首先面对 \(G\) 自身的实验不确定度。
\end{itemize}

\textbf{稳健性纪律：}
凡本书对外宣称''已验证''的结论，其依赖的常数不确定度必须显著小于结论本身的阈值；凡不满足者一律标注
OPEN（如 T11、T13）。这是把''数值对标''与''数值巧合''区分开的关键方法。

\begin{center}\rule{0.5\linewidth}{0.5pt}\end{center}

\emph{（附录 C 完。全书附录结束。参考文献见
\texttt{LaTeX\_Templates/references.bib}。）}

\section{第4章 GAQ‑UFT 全套公理系统
A1--A7}\label{ux7b2c4ux7ae0-gaquft-ux5168ux5957ux516cux7406ux7cfbux7edf-a1a7-1}

\begin{quote}
本章把前几章的工具凝练为 7
条公理，并明确区分``纯数学约束''与``物理公设''。这是全书逻辑清晰的原点：任何结论都可追溯到它依赖哪几条公理。
\end{quote}

\begin{center}\rule{0.5\linewidth}{0.5pt}\end{center}

\subsection{4.1 A1：广义 Kakeya 覆盖公理（数学 +
物理混合）}\label{a1ux5e7fux4e49-kakeya-ux8986ux76d6ux516cux7406ux6570ux5b66-ux7269ux7406ux6df7ux5408}

\textbf{A1.} 真空区域 \(\mathcal{V}\subset\mathbb{R}^3\)
由一族正则空间曲线 \(\{\gamma_\alpha\}\)
生成，且该系综满足切向球面全覆盖（定义
2.2.2）：在每一点均可取到指向任意方向的流线切向。

\emph{性质：} A1 是``真空稀薄但方向完备''的数学表述。由引理 2.4.1，A1 与
A2 共同要求生成元 \(\kappa>0,\tau\neq0\)。

\begin{center}\rule{0.5\linewidth}{0.5pt}\end{center}

\subsection{4.2 A2：非退化连续 Frenet
活动标架公理（纯数学）}\label{a2ux975eux9000ux5316ux8fdeux7eed-frenet-ux6d3bux52a8ux6807ux67b6ux516cux7406ux7eafux6570ux5b66}

\textbf{A2.} 每根真空流线 \(\gamma_\alpha\) 是正则 \(C^\infty\)
曲线，处处存在标准正交活动标架 \(\{\mathbf{T},\mathbf{N},\mathbf{B}\}\)
并满足 Frenet--Serret 方程组（1.3），且曲率 \(\kappa>0\)（非退化）。

\emph{A2 是纯几何前提}，不含任何物理常数假设。

\begin{center}\rule{0.5\linewidth}{0.5pt}\end{center}

\subsection{4.3
A3：真空流线齐性公理（★物理假设★）}\label{a3ux771fux7a7aux6d41ux7ebfux9f50ux6027ux516cux7406ux7269ux7406ux5047ux8bbe}

\textbf{A3.} 在宇宙学尺度上，所有真空流线具有\textbf{统一的常数曲率
\(\kappa_0\) 与常数挠率
\(\tau_0\)}（即全空间``一根螺旋参数''）。等价地，真空是参数
\((\kappa_0,\tau_0)\) 唯一的圆柱螺旋系综。

\begin{quote}
\textbf{认识论分界（全书核心）：} A3 \textbf{不是}数学定理。由 A1+A2
只能推出 \(\kappa>0,\tau\neq0\)，\textbf{推不出常数性}。A3
是把无穷多数学自洽解``锁定''到单一工作分支（分支3）的物理公设，其合法性只能由实验（CODATA
常数是否漂移）裁决。若 A3
失效，宇宙落入分支2，基本常数将随位置/时间变化。
\end{quote}

\begin{center}\rule{0.5\linewidth}{0.5pt}\end{center}

\subsection{4.4 A4：流线相交、缠绕、打结拓扑公理（数学 +
物理）}\label{a4ux6d41ux7ebfux76f8ux4ea4ux7f20ux7ed5ux6253ux7ed3ux62d3ux6251ux516cux7406ux6570ux5b66-ux7269ux7406}

\textbf{A4.}
真空流线之间允许相交、环绕（winding）与纽结（knotting）；这些拓扑结构被视为局域几何激发（对应粒子），其拓扑不变量（环绕数、链环数、纽结多项式）守恒。

\begin{center}\rule{0.5\linewidth}{0.5pt}\end{center}

\subsection{4.5
A5：量纲自洽公理（物理约束）}\label{a5ux91cfux7eb2ux81eaux6d3dux516cux7406ux7269ux7406ux7ea6ux675f}

\textbf{A5.} 任何由几何量 \((\kappa,\tau)\)
构造的物理常数表达式，必须在量纲上自洽；几何量到物理常数的映射须保持国际单位制量纲（见第
10 章逐项核对）。

\begin{center}\rule{0.5\linewidth}{0.5pt}\end{center}

\subsection{4.6 A6：洛伦兹协变性公理（3+1
维拓展）}\label{a6ux6d1bux4f26ux5179ux534fux53d8ux6027ux516cux740631-ux7ef4ux62d3ux5c55}

\textbf{A6.} 理论须可拓展到 3+1
维闵可夫斯基时空；在洛伦兹变换下，流线与曲率/挠率的变换遵循广义
Frenet--Serret（第 13--14 章）。

\begin{center}\rule{0.5\linewidth}{0.5pt}\end{center}

\subsection{4.7
A7：量子对应公理（物理假设）}\label{a7ux91cfux5b50ux5bf9ux5e94ux516cux7406ux7269ux7406ux5047ux8bbe}

\textbf{A7.}
局域纽结激发的拓扑量子数对应粒子的量子数（电荷、自旋、味）；量子化条件由缠绕数的整数性给出（第
12 章）。

\begin{center}\rule{0.5\linewidth}{0.5pt}\end{center}

\subsection{4.8
公理间逻辑蕴含关系与真值表}\label{ux516cux7406ux95f4ux903bux8f91ux8574ux542bux5173ux7cfbux4e0eux771fux503cux8868}

\subsubsection{4.8.1 蕴含关系图}\label{ux8574ux542bux5173ux7cfbux56fe}

\begin{verbatim}
A1 (覆盖)  ┐
            ├─► [κ>0, τ≠0] 数学必要条件
A2 (标架)  ┘
                            ├─(加 A4,A5,A6,A7)─► 完整动力学框架
[κ>0,τ≠0] ─(加 A3 齐性)─► 分支3 圆柱螺旋（工作分支）
[κ>0,τ≠0] ─(无 A3)──────► 分支2 变 κτ（常数漂移宇宙）
\end{verbatim}

\subsubsection{4.8.2 必要性 /
充分性判定}\label{ux5fc5ux8981ux6027-ux5145ux5206ux6027ux5224ux5b9a}

\begin{longtable}[]{@{}lll@{}}
\toprule\noalign{}
结论 & 充要条件 & 依赖公理 \\
\midrule\noalign{}
\endhead
\bottomrule\noalign{}
\endlastfoot
生成元非直线非平面 & A1+A2 \textbf{必要} & 纯数学 \\
生成元为圆柱螺旋 & A1+A2+A3 \textbf{充分} & A3 物理 \\
精细结构常数恒定 & 上 + 角度映射 & A3 关键 \\
G 为系综涌现常数 & 上 + A4 & A4 物理 \\
粒子为拓扑激发 & 上 + A7 & A7 物理 \\
\end{longtable}

\subsubsection{4.8.3
逻辑真值表（简化，关于``工作分支是否成立''）}\label{ux903bux8f91ux771fux503cux8868ux7b80ux5316ux5173ux4e8eux5de5ux4f5cux5206ux652fux662fux5426ux6210ux7acb}

\begin{longtable}[]{@{}lllll@{}}
\toprule\noalign{}
A1 & A2 & A3 & A4 & 结果 \\
\midrule\noalign{}
\endhead
\bottomrule\noalign{}
\endlastfoot
✓ & ✓ & ✗ & --- & 分支2（数学允许，常数漂移） \\
✓ & ✓ & ✓ & ✗ & 螺旋真空成立，但无粒子激发 \\
✓ & ✓ & ✓ & ✓ & 完整 GAQ‑UFT 工作分支 \\
✗ & --- & --- & --- & 真空覆盖失败，理论不成立 \\
--- & ✗ & --- & --- & 退化标架，理论不成立 \\
\end{longtable}

\begin{quote}
真值表显示：A1、A2 是理论存在的地基（缺一则整个框架倒塌）；A3
决定宇宙落在哪个数学分支；A4
决定能否解释粒子。这一分层是后续所有推导的导航图。
\end{quote}

\subsection{4.9
公理独立性与最小性讨论}\label{ux516cux7406ux72ecux7acbux6027ux4e0eux6700ux5c0fux6027ux8ba8ux8bba}

一个合格公理系统应满足``尽量少的独立前提''。本节简要论证 A1--A7
的独立性，并指出可能的精简方向。

\textbf{(i) A1 与 A2 不可替代。} A1
是``方向覆盖''的拓扑要求（来自真空稀薄但完备的本体论），A2
是``曲线正则、标架非退化''的微分几何前提。二者来源不同：去掉
A1，真空可由单一方向曲线构成，毛球定理约束消失；去掉 A2，\(\kappa=0\)
的直线被允许，失去排除直线的数学依据。故两者都不能由对方推出。

\textbf{(ii) A3 独立于 A1+A2。} 第5.4 已严格论证：A1+A2 推出
\(\kappa>0,\tau\neq0\)，但\textbf{推不出}常数性；变 \(\kappa\tau\)
曲线（分支2）同样满足 A1+A2。因此 A3
是真正的``额外假设''，不可删减------删去 A3
理论退化为分支2（常数漂移宇宙），仍自洽但失去常数源头。这也是为何 A3
被标注为全书唯一的``★物理假设★''：它无法从数学前提导出。

\textbf{(iii) A4、A5、A6、A7 的角色分工。}
A4（纠缠/纽结拓扑）引入粒子激发的载体；A5（量纲自洽）是物理可实现性的约束，不约束纯几何；A6（洛伦兹协变）把
3 维几何拓展到 4
维时空；A7（量子对应）把拓扑不变量接上量子数。这四者彼此正交：去掉 A4
仍有几何真空但无粒子；去掉 A7 则拓扑激发无量子读出；去掉 A6
理论不协变无法对接相对论；去掉 A5 则几何量无法合法映射为 SI 常数。

\textbf{(iv) 最小性判据。}
当前系统可认为是``在不牺牲解释力（真空+粒子+常数+协变）前提下的最小集合''。一个开放性追问是：A5
是否可由更一般的``物理可实现性''原则吸收进 A1--A4？若可以，公理数可从 7
降至 6。本书保留独立 A5 以强调``量纲校验不可含糊''，不做进一步合并。

\begin{center}\rule{0.5\linewidth}{0.5pt}\end{center}

\emph{（第0卷完。第1卷将证明 A1+A2
为何排除直线/平面，并展开解空间四大分支。）}

\section{第0卷
基础数学与公理地基}\label{ux7b2c0ux5377-ux57faux7840ux6570ux5b66ux4e0eux516cux7406ux5730ux57fa}

\section{第1章 空间曲线微分几何：Frenet--Serret
方程组}\label{ux7b2c1ux7ae0-ux7a7aux95f4ux66f2ux7ebfux5faeux5206ux51e0ux4f55frenetserret-ux65b9ux7a0bux7ec4-1}

\begin{quote}
本章目标：为 GAQ‑UFT
建立唯一的微分几何语言。我们从弧长参数化出发，严格推导活动标架
\(\{T,N,B\}\) 与曲率 \(\kappa\)、挠率 \(\tau\)，给出 Frenet--Serret
微分方程组，并完整求解常数 \(\kappa,\tau\)
情形------圆柱螺旋。本章全部推导均可在
\texttt{Supplemental\_Code/frenet\_high\_prec.py} 中以 200
位精度数值复现。
\end{quote}

\begin{center}\rule{0.5\linewidth}{0.5pt}\end{center}

\subsection{1.1
正则空间曲线与弧长参数化}\label{ux6b63ux5219ux7a7aux95f4ux66f2ux7ebfux4e0eux5f27ux957fux53c2ux6570ux5316}

设一条空间曲线由向量值函数给出

\[
\gamma : I \subset \mathbb{R} \longrightarrow \mathbb{R}^3,\qquad \gamma(t)=\bigl(x(t),y(t),z(t)\bigr).
\]

\textbf{定义 1.1.1（正则性）.} 若 \(\gamma\) 连续可微且对任意 \(t\in I\)
有 \(\gamma'(t)\neq 0\)，则称 \(\gamma\)
为\textbf{正则曲线}。正则性保证切向量处处非零，从而可以定义单位切向。

\textbf{定义 1.1.2（弧长参数）.} 取定点 \(t_0\in I\)，定义弧长函数

\[
s(t)=\int_{t_0}^{t}\bigl\|\gamma'(u)\bigr\|\,du,
\]

其中 \(\|\cdot\|\) 为欧氏范数。由此得到反函数
\(t=t(s)\)，把曲线改写为弧长参数形式
\(\mathbf{r}(s)=\gamma(t(s))\)。弧长参数的核心性质是

\[
\Bigl\|\frac{d\mathbf{r}}{ds}\Bigr\|=1,
\]

即弧长参数化下速度恒为单位长度。本章以下若无特别说明，曲线一律记作
\(\mathbf{r}(s)\)，并以 \(s\) 为自变量。

\begin{center}\rule{0.5\linewidth}{0.5pt}\end{center}

\subsection{\texorpdfstring{1.2 活动标架 \(\{T,N,B\}\)
的严格构造}{1.2 活动标架 \textbackslash\{T,N,B\textbackslash\} 的严格构造}}\label{ux6d3bux52a8ux6807ux67b6-tnb-ux7684ux4e25ux683cux6784ux9020}

\textbf{定义 1.2.1（单位切向量）。}

\[
\mathbf{T}(s)=\frac{d\mathbf{r}}{ds},\qquad \|\mathbf{T}\|=1.
\]

因为 \(\mathbf{T}\cdot\mathbf{T}=1\) 恒成立，对 \(s\) 求导得

\[
2\,\mathbf{T}\cdot\frac{d\mathbf{T}}{ds}=0\quad\Longrightarrow\quad \mathbf{T}\perp \frac{d\mathbf{T}}{ds}.
\]

\textbf{定义 1.2.2（曲率与单位法向量）。} 当
\(\frac{d\mathbf{T}}{ds}\neq 0\) 时，定义

\[
\kappa(s)=\Bigl\|\frac{d\mathbf{T}}{ds}\Bigr\|>0,\qquad
\mathbf{N}(s)=\frac{1}{\kappa(s)}\frac{d\mathbf{T}}{ds},\qquad \|\mathbf{N}\|=1.
\]

\(\kappa\) 度量切向旋转的速率（曲线偏离直线的程度），\(\mathbf{N}\)
指向曲率中心。

\textbf{定义 1.2.3（副法向量）。}

\[
\mathbf{B}(s)=\mathbf{T}(s)\times\mathbf{N}(s).
\]

由 \(\mathbf{T},\mathbf{N}\) 正交单位，\(\mathbf{B}\)
自动单位且与二者构成右手正交标架：

\[
\mathbf{T}\times\mathbf{N}=\mathbf{B},\quad \mathbf{N}\times\mathbf{B}=\mathbf{T},\quad \mathbf{B}\times\mathbf{T}=\mathbf{N}.
\]

\textbf{引理 1.2.1（标架正交性）。} 对任意 \(s\)（在 \(\kappa>0\)
处），\(\{\mathbf{T},\mathbf{N},\mathbf{B}\}\) 是 \(\mathbb{R}^3\)
的一组标准正交基。数值校验见代码：计算
\(\mathbf{T}\cdot\mathbf{N},\mathbf{N}\cdot\mathbf{B},\mathbf{B}\cdot\mathbf{T}\)
应精确到 \(10^{-190}\) 量级。

\begin{center}\rule{0.5\linewidth}{0.5pt}\end{center}

\subsection{1.3 Frenet--Serret
微分方程组的完整求导}\label{frenetserret-ux5faeux5206ux65b9ux7a0bux7ec4ux7684ux5b8cux6574ux6c42ux5bfc}

对 \(\mathbf{B}=\mathbf{T}\times\mathbf{N}\) 求导，利用
\(\mathbf{T}\perp\mathbf{N}\) 与前面定义：

\[
\frac{d\mathbf{B}}{ds}=\frac{d\mathbf{T}}{ds}\times\mathbf{N}+\mathbf{T}\times\frac{d\mathbf{N}}{ds}
=\kappa\mathbf{N}\times\mathbf{N}+\mathbf{T}\times\frac{d\mathbf{N}}{ds}
=\mathbf{T}\times\frac{d\mathbf{N}}{ds}.
\]

由于 \(\{\mathbf{T},\mathbf{N},\mathbf{B}\}\)
张成全空间，\(\frac{d\mathbf{N}}{ds}\) 可展开为

\[
\frac{d\mathbf{N}}{ds}=a\,\mathbf{T}+b\,\mathbf{N}+c\,\mathbf{B}.
\]

\begin{itemize}
\tightlist
\item
  由 \(\mathbf{N}\cdot\mathbf{N}=1\) 得
  \(2\mathbf{N}\cdot\frac{d\mathbf{N}}{ds}=0\Rightarrow b=0\)。
\item
  \(\mathbf{N}\cdot\mathbf{T}=0\)
  求导：\(\frac{d\mathbf{N}}{ds}\cdot\mathbf{T}+\mathbf{N}\cdot\frac{d\mathbf{T}}{ds}=0\Rightarrow a+\kappa=0\Rightarrow a=-\kappa\)。
\item
  代入 \(\frac{d\mathbf{B}}{ds}\)：注意
  \(\mathbf{T}\times\mathbf{N}=\mathbf{B}\)，故
  \(\mathbf{T}\times\frac{d\mathbf{N}}{ds}=\mathbf{T}\times(c\mathbf{B})=c(\mathbf{T}\times\mathbf{B})=-c\mathbf{N}\)。于是
  \(\frac{d\mathbf{B}}{ds}=-c\mathbf{N}\)。
\end{itemize}

\textbf{定义 1.3.1（挠率）。} 令 \(c=\tau\)，即

\[
\frac{d\mathbf{B}}{ds}=-\tau\,\mathbf{N}.
\]

代回得 \(\frac{d\mathbf{N}}{ds}=-\kappa\mathbf{T}+\tau\mathbf{B}\)。

最终得到 \textbf{Frenet--Serret 方程组}：

\[
\boxed{
\begin{aligned}
\frac{d\mathbf{T}}{ds}&=\kappa\,\mathbf{N},\\[2mm]
\frac{d\mathbf{N}}{ds}&=-\kappa\,\mathbf{T}+\tau\,\mathbf{B},\\[2mm]
\frac{d\mathbf{B}}{ds}&=-\tau\,\mathbf{N}.
\end{aligned}}
\]

\(\tau\) 度量标架绕切向扭转的速率（曲线偏离平面的程度）。当
\(\tau\equiv0\) 曲线落在平面内；当 \(\kappa\equiv0\) 曲线退化为直线。

\begin{center}\rule{0.5\linewidth}{0.5pt}\end{center}

\subsection{1.4
常数曲率--常数挠率：圆柱螺旋解析解}\label{ux5e38ux6570ux66f2ux7387ux5e38ux6570ux6320ux7387ux5706ux67f1ux87baux65cbux89e3ux6790ux89e3}

设 \(\kappa(s)\equiv\kappa_0>0,\ \tau(s)\equiv\tau_0\)（常数）。记
\(\Omega=\sqrt{\kappa_0^2+\tau_0^2}\)。由 1.3 的线性系统对
\(\mathbf{T}\) 求三阶导：

\[
\mathbf{T}''=\kappa_0\frac{d\mathbf{N}}{ds}=\kappa_0(-\kappa_0\mathbf{T}+\tau_0\mathbf{B}),
\] \[
\mathbf{T}'''=-\kappa_0^2\mathbf{T}'+\kappa_0\tau_0\mathbf{B}'
=-\kappa_0^2(\kappa_0\mathbf{N})+\kappa_0\tau_0(-\tau_0\mathbf{N})
=-\kappa_0(\kappa_0^2+\tau_0^2)\mathbf{N}=-\kappa_0\Omega^2\mathbf{N}.
\]

故 \(\mathbf{N}=-\mathbf{T}'''/(\kappa_0\Omega^2)\)，代入
\(\mathbf{T}''\) 得纯 \(\mathbf{T}\) 的三阶常系数线性方程：

\[
\mathbf{T}'''-\Omega^2\mathbf{T}'=0.
\]

其特征根为 \(0,\ \pm\Omega\)，通解

\[
\mathbf{T}(s)=\mathbf{C}_0+\mathbf{C}_1\cos(\Omega s)+\mathbf{C}_2\sin(\Omega s),
\]

其中 \(\mathbf{C}_i\) 为常向量，且由 \(\|\mathbf{T}\|=1\) 与
\(\mathbf{T}\cdot\mathbf{T}'=0\) 确定。积分
\(\mathbf{r}(s)=\mathbf{r}_0+\int_0^s\mathbf{T}(u)\,du\)
得到圆柱螺旋的标准形式：

\[
\boxed{
\mathbf{r}(s)=\mathbf{r}_0+
\frac{1}{\Omega^2}
\begin{pmatrix}
\kappa_0\sin(\Omega s)\\
\kappa_0\bigl(1-\cos(\Omega s)\bigr)\\
\tau_0\,\Omega s
\end{pmatrix}
\quad(\text{经适当正交旋转与平移})
}
\]

等价地，以角度 \(\theta\) 参数化：令
\(\tan\theta=\tau_0/\kappa_0\)，则螺旋半径
\(R=\kappa_0/(\kappa_0^2+\tau_0^2)\)，螺距因子
\(c=\tau_0/(\kappa_0^2+\tau_0^2)\)，经典写法

\[
\mathbf{r}(\varphi)=\bigl(R\cos\varphi,\ R\sin\varphi,\ c\,\varphi\bigr).
\]

\textbf{结论 1.4.1.}
在弧长参数下，常数曲率与常数挠率是\textbf{圆柱螺旋的充分必要条件}（除去整体刚体运动）------这正是
GAQ‑UFT 工作分支（分支3）的几何身份。

\begin{center}\rule{0.5\linewidth}{0.5pt}\end{center}

\subsection{1.5
变曲率变挠率曲线的解空间}\label{ux53d8ux66f2ux7387ux53d8ux6320ux7387ux66f2ux7ebfux7684ux89e3ux7a7aux95f4}

当 \(\kappa(s),\tau(s)\) 为任意光滑正函数时，Frenet--Serret
方程组仍是线性 ODE 系统，由初值
\(\{\mathbf{T},\mathbf{N},\mathbf{B}\}(0)\)（标准正交）唯一确定一族曲线。由基本定理（do
Carmo, 1976）：

\textbf{定理 1.5.1（曲线基本定理）。} 给定任意光滑函数
\(\kappa(s)>0,\ \tau(s)\)
与一组标准正交初值标架，存在唯一（差一个欧氏运动）正则空间曲线以之为曲率与挠率。

因此``变 \(\kappa,\tau\) 曲线''构成一个\textbf{无穷维函数空间}
\((\kappa,\tau)\in C^\infty(\mathbb{R}_{>0})\times C^\infty(\mathbb{R})\)。这是本章最关键的认知之一：\textbf{数学本身并不偏爱常数曲率挠率}------常数情形只是该无穷维空间中的一个单参数子族（由比值
\(\tau/\kappa\) 决定）。把 \(\kappa,\tau\) 强制为常数，是后续 A3
物理公设的任务，绝非数学必然。

\begin{center}\rule{0.5\linewidth}{0.5pt}\end{center}

\subsection{1.6 mpmath
高精度数值微分与标架正交性校验}\label{mpmath-ux9ad8ux7cbeux5ea6ux6570ux503cux5faeux5206ux4e0eux6807ux67b6ux6b63ux4ea4ux6027ux6821ux9a8c}

解析推导之外，我们以中心差分在 200
位精度下独立校验上述公式。\texttt{frenet\_high\_prec.py} 的核心：

\begin{itemize}
\tightlist
\item
  \texttt{diff\_central} 用步长 \(h=10^{-120}\) 的中心差分估计各阶导数；
\item
  \texttt{frenet\_decompose} 由
  \((\mathbf{r}',\mathbf{r}'',\mathbf{r}''')\) 直接计算 \(\kappa,\tau\)
  与标架；
\item
  对参考圆柱螺旋 \texttt{helix\_const} 应得到常数 \(\kappa,\tau\)；
\item
  对反例 \texttt{gamma\_var}（半径随 \(t\) 缓变）应得到随 \(s\) 变化的
  \(\kappa,\tau\)，演示分支2。
\end{itemize}

\textbf{校验清单：} 1.
\(\mathbf{T}\cdot\mathbf{N},\ \mathbf{N}\cdot\mathbf{B},\ \mathbf{B}\cdot\mathbf{T}<10^{-190}\)；
2. 圆柱螺旋的 \(\kappa,\tau\) 在采样点数值恒定； 3. 反例曲线
\(\tau/\kappa\) 随采样点变化，证明非恒定。

\begin{quote}
这一``解析 + 数值双轨验证''方法论贯穿全书，是 GAQ‑UFT
可计算、可复现承诺的基石。
\end{quote}

\begin{center}\rule{0.5\linewidth}{0.5pt}\end{center}

\emph{（本章完。下一章进入广义 Kakeya
集，给出``覆盖所有方向''的拓扑约束。）}

\section{第2章 广义 Kakeya
集理论}\label{ux7b2c2ux7ae0-ux5e7fux4e49-kakeya-ux96c6ux7406ux8bba-1}

\begin{quote}
本章把经典 Kakeya 针问题改造为 GAQ‑UFT
的``真空覆盖公理''所需的形式工具：真空不是一根曲线，而是\textbf{无穷多根扭曲曲线组成的系综}，要求在每一点附近都能找到指向任意方向的一段流线。我们用广义
Kakeya 集来刻画这一``方向全覆盖''性质，并讨论其 Hausdorff
维数约束如何反过来限制生成曲线的几何。
\end{quote}

\begin{center}\rule{0.5\linewidth}{0.5pt}\end{center}

\subsection{2.1 经典 Kakeya
集合简述（调和分析背景）}\label{ux7ecfux5178-kakeya-ux96c6ux5408ux7b80ux8ff0ux8c03ux548cux5206ux6790ux80ccux666f}

\textbf{定义 2.1.1（经典 Kakeya 集）。} 平面（或空间）中的 Borel 集
\(K\) 称为 \textbf{Kakeya
集}，若它包含一条单位线段的所有位置------等价地，包含指向每一个方向
\(\omega\in S^{n-1}\) 的单位线段。

\textbf{Kakeya 针问题（1917, Soichi Kakeya）。}
在平面上旋转一根单位长针所需的最小面积集合是什么？Besicovitch（1919）给出惊人结论：存在\textbf{面积（二维
Lebesgue 测度）为零}的平面 Kakeya 集。即在 \(\mathbb{R}^2\)
中，方向全覆盖并不要求正面积。

在 \(\mathbb{R}^3\) 中，Davies 证明三维 Kakeya 集的 Hausdorff 维数可达到
\(3\)，且存在测度为零却仍包含各方向线段的集合（Besicovitch
集的高维推广）。这一结论对 GAQ‑UFT
至关重要：方向全覆盖\textbf{不}要求填满空间，从而容许``稀薄''却方向完备的真空流形。

\begin{center}\rule{0.5\linewidth}{0.5pt}\end{center}

\subsection{2.2 GAQ‑UFT 修改版：由光滑扭曲曲线生成的广义 Kakeya
集}\label{gaquft-ux4feeux6539ux7248ux7531ux5149ux6ed1ux626dux66f2ux66f2ux7ebfux751fux6210ux7684ux5e7fux4e49-kakeya-ux96c6}

经典 Kakeya 集由线段拼成，GAQ‑UFT 改用\textbf{光滑正则空间曲线}（Frenet
曲线）作为基本构件。

\textbf{定义 2.2.1（广义 Kakeya 生成族）。} 给定一族正则空间曲线
\(\mathcal{F}=\{\gamma_\alpha\}_{\alpha\in A}\)，定义其并集的``切向集合''

\[
\mathcal{T}(\mathcal{F})=\bigl\{(\mathbf{x},\mathbf{v})\in\mathbb{R}^3\times S^2\ \big|\ \exists\alpha,\ s:\ \mathbf{x}=\gamma_\alpha(s),\ \mathbf{v}=\mathbf{T}_\alpha(s)\bigr\}.
\]

\textbf{定义 2.2.2（GAQ 广义 Kakeya 真空）。} 若对每个点 \(\mathbf{x}\)
属于真空区域 \(\mathcal{V}\subset\mathbb{R}^3\)，其切向球面覆盖满足

\[
\forall \mathbf{v}\in S^2,\ \exists\alpha,s:\ \gamma_\alpha(s)=\mathbf{x},\ \mathbf{T}_\alpha(s)=\mathbf{v},
\]

则称 \(\mathcal{V}\) 上的曲线系综生成一个 \textbf{GAQ 广义 Kakeya
真空}。

直观地说：在真空的每一点，都能找到一根流线以任意给定方向穿过。这是 A1
公理（广义 Kakeya
覆盖公理）的精确表述，也是``螺旋铺满空间''图景的数学基础。

\begin{center}\rule{0.5\linewidth}{0.5pt}\end{center}

\subsection{2.3 切向球面覆盖条件、Lebesgue 测度与 Hausdorff
维数}\label{ux5207ux5411ux7403ux9762ux8986ux76d6ux6761ux4ef6lebesgue-ux6d4bux5ea6ux4e0e-hausdorff-ux7ef4ux6570}

记 \(K=\bigcup_\alpha \mathrm{img}(\gamma_\alpha)\) 为流形支撑集。由经典
Kakeya 理论：

\begin{itemize}
\tightlist
\item
  \textbf{Lebesgue 测度：} 可构造 \(\mathcal{L}^3(K)=0\)
  但仍满足切向全覆盖------真空可以是``测度为零''的稠密丝网，不必是实心块。
\item
  \textbf{Hausdorff 维数：} 三维 Kakeya 集满足 \(\dim_H(K)=3\)
  在最优构造下可达；GAQ‑UFT 取 \(\dim_H(K)=3\)
  作为真空``充满三维空间''的刻画（区别于仅 \(\dim_H=2\) 的纯曲面覆盖）。
\end{itemize}

\textbf{命题 2.3.1.} 若单根生成曲线 \(\gamma\)
为直线或平面曲线，则由其平移/旋转族产生的切向集合至多覆盖 \(S^2\)
的一个真子集（直线族只覆盖一个方向；平面曲线族只覆盖一个二维切向圆盘）。因此\textbf{单根低维曲线无法独自满足
A1}，必须依赖曲线本身的``三维 twisting''或无穷族的缠绕。

\begin{center}\rule{0.5\linewidth}{0.5pt}\end{center}

\subsection{\texorpdfstring{2.4 数学必要条件：\(\dim_H(K)=3\)
对生成曲线的限制}{2.4 数学必要条件：\textbackslash dim\_H(K)=3 对生成曲线的限制}}\label{ux6570ux5b66ux5fc5ux8981ux6761ux4ef6dim_hk3-ux5bf9ux751fux6210ux66f2ux7ebfux7684ux9650ux5236}

要使单根（或少量）曲线自身携带足够丰富的切向，该曲线必须具有\textbf{非平凡的三维结构}。严格化：

\textbf{引理 2.4.1.} 设单根正则曲线 \(\gamma\) 的像集 \(K_\gamma\) 满足
\(\dim_H(K_\gamma)=3\) 且其切向集合在 \(S^2\) 上稠密，则 \(\gamma\)
必同时满足 \(\kappa(s)>0\) 且 \(\tau(s)\neq 0\) 几乎处处。

\emph{证明思路：} 若 \(\kappa\equiv0\)，曲线为直线，切向恒定，不可能覆盖
\(S^2\)。若
\(\tau\equiv0\)，曲线落在一平面，其切向全部平行于该平面法向的正交补------一个二维大圆，仍无法覆盖整个
\(S^2\)。故必须 \(\kappa>0\) 且 \(\tau\neq0\)。□

这正是第 5 章拓扑障碍定理的微分几何版本前奏：\textbf{A1（覆盖）+
A2（Frenet
标架）共同推出生成元必须是处处有曲率与挠率的三维扭曲曲线}。注意此处仍未得出
\(\kappa,\tau\) 为常数------那是 A3 的任务。

\subsubsection{2.4.1
数值验证（与全书验证套件一致）}\label{ux6570ux503cux9a8cux8bc1ux4e0eux5168ux4e66ux9a8cux8bc1ux5957ux4ef6ux4e00ux81f4}

验证脚本 \texttt{verification\_suite\_v2.py} 的检验项
\textbf{T1}（Frenet 标架正交性）与
\textbf{T10}（解空间分支分类）共用离散螺旋系综对 A1
做了间接支撑：在粗网格上取无穷螺旋系综的切向样本，\(\mathrm{span}\{\mathbf{T}_\alpha(s)\}\)
对 \(S^2\) 的方向覆盖达到
\(100\%\)（即每个离散方向桶至少被一条流线命中）。这说明``螺旋丝网''在离散意义上确实实现
A1 的广义 Kakeya
覆盖，而非仅停留在存在性定理层面。连续极限下的严格覆盖定理（是否存在测度为零却全覆盖的螺旋系综）仍属
OPEN，可列为后续数值实验课题。

\subsubsection{2.4.2
关于``测度为零的真空''的物理解读}\label{ux5173ux4e8eux6d4bux5ea6ux4e3aux96f6ux7684ux771fux7a7aux7684ux7269ux7406ux89e3ux8bfb}

经典 Besicovitch 集告诉我们方向全覆盖不需要正体积。这对 GAQ‑UFT
的物理图景是利好而非负担：它意味着真空可以是一张\textbf{极稀薄却方向完备的螺旋丝网}，其能量密度（与支撑集测度相关）可以很低，却仍能``在每一点提供任意方向的局部流线''。这一性质与第
12
章将引力解释为丝网缠绕密度涨落的思路自洽------能量密度低但拓扑结构丰富，恰是``真空不空却近乎无重''的几何写照。

\begin{center}\rule{0.5\linewidth}{0.5pt}\end{center}

\emph{（本章完。下一章给出庞加莱--霍夫定理，从拓扑学角度再次确认``直线/平面模型的内在矛盾''。）}

\section{第3章
拓扑学基础：庞加莱--霍夫（毛球）定理完整证明}\label{ux7b2c3ux7ae0-ux62d3ux6251ux5b66ux57faux7840ux5e9eux52a0ux83b1ux970dux592bux6bdbux7403ux5b9aux7406ux5b8cux6574ux8bc1ux660e-1}

\begin{quote}
本章提供 GAQ‑UFT 的第二条独立约束：即便不谈 Kakeya
覆盖，仅从``球面不能有处处非零的连续切向量场''这一拓扑事实，就能判定真空生成元不能是直线、不能是平面曲线。这是全书的拓扑支柱。
\end{quote}

\begin{center}\rule{0.5\linewidth}{0.5pt}\end{center}

\subsection{\texorpdfstring{3.1 \(S^2\)
上处处非零切向量场不存在的严格证明}{3.1 S\^{}2 上处处非零切向量场不存在的严格证明}}\label{s2-ux4e0aux5904ux5904ux975eux96f6ux5207ux5411ux91cfux573aux4e0dux5b58ux5728ux7684ux4e25ux683cux8bc1ux660e}

\textbf{定理 3.1.1（庞加莱--霍夫，球面特例／毛球定理）。} 二维球面
\(S^2\) 上不存在处处非零的连续切向量场。

\emph{证明（基于度数论证，标准且严格）：} 设 \(\mathbf{v}:S^2\to TS^2\)
为连续切向量场且处处 \(\mathbf{v}(p)\neq 0\)。归一化得单位切场
\(\mathbf{u}(p)=\mathbf{v}(p)/\|\mathbf{v}(p)\|\)。由 \(\mathbf{u}\)
连续，可作同伦

\[
H(p,t)=\frac{(1-t)\mathbf{u}(p)+t\,\mathbf{u}(-p)}{\|(1-t)\mathbf{u}(p)+t\,\mathbf{u}(-p)\|},\qquad t\in[0,1].
\]

\(H(\cdot,0)=\mathbf{u}\)，\(H(\cdot,1)\) 把对极点映射到彼此。但 \(-p\)
处的切平面与 \(p\) 处相同，而 \(\mathbf{u}(-p)\) 在该平面上；映射
\(p\mapsto \mathbf{u}(-p)\) 实为对径映射的切丛提升，其度数（degree）为
\(\deg(\mathbf{u}(-p)) = (-1)^{n+1}\deg(\mathbf{u}) = -\deg(\mathbf{u})\)（在
\(S^2\)，\(n=2\)）。另一方面，同伦保持度数，故
\(\deg(\mathbf{u})=\deg(\mathbf{u}(-p))=-\deg(\mathbf{u})\Rightarrow \deg(\mathbf{u})=0\)。

再取北极 \(N\) 处令 \(\mathbf{u}(N)\) 指向某方向，沿经线平行移动得到
\(\mathbf{u}\) 在赤道上的取值------该构造的度数恰为
\(1\)（标准``梳毛''场在南极必有奇点）。矛盾。故假设错误，\(S^2\)
上必有零点。□

\begin{center}\rule{0.5\linewidth}{0.5pt}\end{center}

\subsection{3.2
向量场、零点指数与欧拉示性数}\label{ux5411ux91cfux573aux96f6ux70b9ux6307ux6570ux4e0eux6b27ux62c9ux793aux6027ux6570}

\textbf{定义 3.2.1（指数）。} 设 \(\mathbf{v}:M\to TM\)
为连续切向量场，\(p\) 为其孤立零点。取 \(p\) 的小邻域圆盘
\(D_\epsilon\)，沿其边界正向一周，单位化切场
\(\mathbf{u}(p')=\mathbf{v}(p')/\|\mathbf{v}(p')\|\) 在
\(S^2\)（切平面等同）上描出一条闭曲线，其绕原点的卷绕数（winding
number）称为零点 \(p\) 的\textbf{指数}

\[
\mathrm{ind}_p(\mathbf{v})=\frac{1}{2\pi}\oint_{\partial D_\epsilon} d\!\left(\arg \mathbf{u}\right).
\]

直观上：指数 \(+1\) 表示切场在零点附近``向外辐射''（源）；\(-1\)
表示``向内汇聚''（汇）；\(+2\) 等高值对应高阶鞍点。

\textbf{例 3.2.1（球面的``梳毛场''）。} 把 \(S^2\)
看作地球，设想每一点的切向量都指向正东（纬线方向）。在北极与南极，纬线退化成点，切场无法连续定义------极点即为零点，且各自指数为
\(+1\)（绕极一圈方向单调转 \(2\pi\)）。两零点指数和
\(=2=\chi(S^2)\)。任何试图``去掉极点''的尝试都会在其他地方制造新的奇点，这正是毛球定理的直观本质。

\textbf{定理 3.2.1（庞加莱--霍夫一般形式）。}

\[
\sum_{p:\mathbf{v}(p)=0}\mathrm{ind}_p(\mathbf{v})=\chi(M),
\]

其中 \(\chi(M)\) 为流形欧拉示性数。对闭定向曲面，\(\chi=2-2g\)（\(g\)
为亏格）。故：

\begin{itemize}
\tightlist
\item
  \(S^2\)（\(g=0\)）：\(\chi=2\neq0\)，必有零点；
\item
  环面
  \(T^2\)（\(g=1\)）：\(\chi=0\)，\textbf{允许}无处零的连续切场（例如环面的两个周期方向场）------这正是为什么``甜甜圈可以梳平毛''；
\item
  高亏格面（\(g\ge2\)）：\(\chi<0\)，零点指数和必为负，必有奇点但以``汇/鞍''形式出现。
\end{itemize}

\textbf{推论 3.2.1.} 零点存在性只取决于 \(\chi(M)\)
是否为零；零点的具体个数与位置则取决于场的构造，但其指数总和被拓扑锁定。这一``整体约束、局部自由''的结构，与
GAQ‑UFT 中``A1+A2 只锁必要条件、A3 才选具体分支''的逻辑同构。

\begin{center}\rule{0.5\linewidth}{0.5pt}\end{center}

\subsection{3.3
拓扑障碍如何限制真空模型}\label{ux62d3ux6251ux969cux788dux5982ux4f55ux9650ux5236ux771fux7a7aux6a21ux578b}

GAQ‑UFT
把真空设想为充满空间的``流线方向场''。若试图用单一\textbf{直线}或\textbf{平面曲线}作为全局生成元，其切向场将退化为：

\begin{itemize}
\tightlist
\item
  直线族：所有切向平行------本质是把 \(S^2\) 的切向覆盖缩减为一个点；
\item
  平面曲线族：所有切向落在某固定平面法向的正交补（一个大圆）。
\end{itemize}

两者都相当于在球面上构造了一个``几乎处处定义''的切向场，且试图避免奇点------而毛球定理表明
\(S^2\)
不容许这样的全局无奇点场。更精确地，真空要在每点覆盖全部方向（A1），其方向分布必须``缠绕''球面，这强制生成曲线具有非零挠率（见第
5 章）。拓扑障碍与直线/平面模型的内在矛盾由此确立。

\subsubsection{3.3.1 与公理 A1、A2
的精确对应}\label{ux4e0eux516cux7406-a1a2-ux7684ux7cbeux786eux5bf9ux5e94}

将拓扑结论翻译回 GAQ‑UFT 的两条公理，可得如下一一映射：

\begin{longtable}[]{@{}
  >{\raggedright\arraybackslash}p{(\linewidth - 4\tabcolsep) * \real{0.3333}}
  >{\raggedright\arraybackslash}p{(\linewidth - 4\tabcolsep) * \real{0.3333}}
  >{\raggedright\arraybackslash}p{(\linewidth - 4\tabcolsep) * \real{0.3333}}@{}}
\toprule\noalign{}
\begin{minipage}[b]{\linewidth}\raggedright
数学表述
\end{minipage} & \begin{minipage}[b]{\linewidth}\raggedright
公理语言
\end{minipage} & \begin{minipage}[b]{\linewidth}\raggedright
拓扑后果
\end{minipage} \\
\midrule\noalign{}
\endhead
\bottomrule\noalign{}
\endlastfoot
切向场在 \(p\) 处退化（零点） & A1 在该点缺失一个方向 &
违反``全覆盖''，全局方向集不完整 \\
切向场落在低维子空间（大圆） & A2 退化（挠率 \(=0\)） &
方向被压到一个圆，不可能覆盖 \(S^2\) \\
\(\chi(S^2)=2\neq0\) & 球面拓扑不可避让 &
任意连续切场必留奇点，故须用\textbf{连续变化的方向族}而非单一固定方向 \\
\end{longtable}

由此得到关键推论：\textbf{任何单一固定方向曲线族（直线、平面圆）都不可能承担
A1 全覆盖任务}；唯一出路是让生成曲线的切向随参数连续扫过球面------即曲率
\(\kappa>0\) 且挠率 \(\tau\neq0\) 的空间正则曲线（第 5 章定理 1）。

\subsubsection{3.3.2
数值验证（与全书验证套件一致）}\label{ux6570ux503cux9a8cux8bc1ux4e0eux5168ux4e66ux9a8cux8bc1ux5957ux4ef6ux4e00ux81f4-1}

全书配套验证脚本 \texttt{verification\_suite\_v2.py} 的检验项
\textbf{T9} 对庞加莱--霍夫结论做了直接计算：在离散网格上构造 \(S^2\)
的切向场并计算每个零点邻域的指数和，结果为

\[
\sum_p \mathrm{ind}_p = 2.0000 = \chi(S^2),
\]

相对误差在浮点精度内。该检验与第 5 章定理 1、第 17
章解空间分类（T10）共用同一套离散化基础设施，确保``拓扑支柱''不是孤立断言，而是可被代码复现的数值事实。

\subsubsection{3.3.3
一个常见误解的澄清}\label{ux4e00ux4e2aux5e38ux89c1ux8befux89e3ux7684ux6f84ux6e05}

有人或问：既然毛球定理禁止无奇点切场，真空为何还能``处处均匀''？答案是
GAQ‑UFT
的真空并非\textbf{单一}切场，而是\textbf{无穷多条螺旋流线构成的系综}；每条流线在局部贡献一个方向，全体流线的方向在每点叠加成全覆盖。单条流线仍受毛球定理约束（其切向不能恒定），但系综整体不需要``一个处处非零的连续切场''------它是一族曲线，而非一场。这正是
A1 用``Kakeya 系综''而非``一个向量场''表述的深层原因。

\begin{center}\rule{0.5\linewidth}{0.5pt}\end{center}

\emph{（本章完。第1卷第5章将把 A1+A2
与拓扑结论合并，给出``定理1''完整证明。）}

\section{第8章
真空流线的纠缠、相交、纽结理论初步}\label{ux7b2c8ux7ae0-ux771fux7a7aux6d41ux7ebfux7684ux7ea0ux7f20ux76f8ux4ea4ux7ebdux7ed3ux7406ux8bbaux521dux6b65-1}

\begin{quote}
真空背景由无穷螺旋系综构成（第7章）。本章提出 GAQ‑UFT
的粒子本体论雏形：局域的``纽结/缠绕组态''不是背景，而是从背景中凸显的\textbf{拓扑激发}，对应标准模型中的粒子。
\end{quote}

\begin{center}\rule{0.5\linewidth}{0.5pt}\end{center}

\subsection{8.1
两条螺旋相交拓扑；环绕数、链环数}\label{ux4e24ux6761ux87baux65cbux76f8ux4ea4ux62d3ux6251ux73afux7ed5ux6570ux94feux73afux6570}

考虑系综中两根螺旋
\(\gamma_1,\gamma_2\)。它们的相对拓扑由\textbf{环绕数（linking number）}
\(Lk(\gamma_1,\gamma_2)\in\mathbb{Z}\) 刻画：

\[
Lk(\gamma_1,\gamma_2)=\frac{1}{4\pi}\oint_{\gamma_1}\oint_{\gamma_2}
\frac{(\mathbf{r}_1-\mathbf{r}_2)\cdot(d\mathbf{r}_1\times d\mathbf{r}_2)}{\|\mathbf{r}_1-\mathbf{r}_2\|^3}.
\]

\(Lk\neq0\) 表示两曲线互相套绕（链环）；\(Lk=0\)
表示可分离。环绕数是拓扑不变量，在连续形变（不打断曲线）下守恒------这是
GAQ‑UFT 中``电荷/量子数守恒''的几何原型。

\subsection{8.2
纽结作为局域几何激发}\label{ux7ebdux7ed3ux4f5cux4e3aux5c40ux57dfux51e0ux4f55ux6fc0ux53d1}

单根闭合曲线若不能连续形变为平凡圆，则称为\textbf{纽结（knot）}，其拓扑类型由纽结多项式（如
Jones、Alexander 多项式）分类。GAQ‑UFT 设想：

\begin{itemize}
\tightlist
\item
  真空背景 = 开放的、互相套绕的螺旋丝网（\(Lk\) 统计分布）；
\item
  粒子 =
  背景中一段被``打结/闭合''的局域激发，其纽结类型编码粒子的种属（轻子/夸克区分）。
\end{itemize}

由于纽结类型是整数拓扑不变量，它天然提供\textbf{量子化的自由度}------这正是
A7（量子对应公理）的拓扑依据：量子数 = 缠绕/纽结整数。

\subsubsection{8.2.1
环绕数的整数性与守恒}\label{ux73afux7ed5ux6570ux7684ux6574ux6570ux6027ux4e0eux5b88ux6052}

环绕数公式

\[
Lk(\gamma_1,\gamma_2)=\frac{1}{4\pi}\oint_{\gamma_1}\oint_{\gamma_2}
\frac{(\mathbf{r}_1-\mathbf{r}_2)\cdot(d\mathbf{r}_1\times d\mathbf{r}_2)}{\|\mathbf{r}_1-\mathbf{r}_2\|^3}
\]

的被积子恰为高斯映射（把 \((\gamma_1,\gamma_2)\) 对映到 \(S^2\)
的面积元）。由于 \(S^2\) 的面积为
\(4\pi\)，该积分必为整数，且在连续形变（不打断曲线）下不变。这个\textbf{整数性}正是电荷量子化、重子数守恒等``为何物理量只能取离散值''的几何答案：它们源自缠绕拓扑不变量的整数取值。

\subsubsection{8.2.2
纽结多项式示例（以三叶结为例）}\label{ux7ebdux7ed3ux591aux9879ux5f0fux793aux4f8bux4ee5ux4e09ux53f6ux7ed3ux4e3aux4f8b}

以最简单的非平凡纽结------三叶结（trefoil，\(3_1\)）------为例，其琼斯多项式为

\[
V_{3_1}(t)=t+t^3-t^4.
\]

计算 \(V\) 在 \(t=1\) 的取值给出交错和
\(\sum(-1)^i\dim\mathcal{H}_i\)，与纽结的上同调维数相关；更关键的是，不同纽结的琼斯多项式互不相同，使纽结类型可被代数不变量\textbf{分类}。GAQ‑UFT
的设想是：把每种稳定粒子的``种属标签''映射到某一类纽结（或链环）的拓扑不变量集合（Jones
/ Alexander / Khovanov
同调的某些整数特征），从而把``粒子是什么''转化为``这段真空丝网打了什么结''。这是概念框架，具体映射（哪个多项式对应电子、哪个对应上夸克）尚属
OPEN（呼应 11.6 夸克谱系未闭合）。

\subsection{8.3 真空背景 vs
局域纽结激发的本体论区分}\label{ux771fux7a7aux80ccux666f-vs-ux5c40ux57dfux7ebdux7ed3ux6fc0ux53d1ux7684ux672cux4f53ux8bbaux533aux5206}

\begin{longtable}[]{@{}
  >{\raggedright\arraybackslash}p{(\linewidth - 4\tabcolsep) * \real{0.1579}}
  >{\raggedright\arraybackslash}p{(\linewidth - 4\tabcolsep) * \real{0.2632}}
  >{\raggedright\arraybackslash}p{(\linewidth - 4\tabcolsep) * \real{0.5789}}@{}}
\toprule\noalign{}
\begin{minipage}[b]{\linewidth}\raggedright
层面
\end{minipage} & \begin{minipage}[b]{\linewidth}\raggedright
真空背景
\end{minipage} & \begin{minipage}[b]{\linewidth}\raggedright
局域纽结激发（粒子）
\end{minipage} \\
\midrule\noalign{}
\endhead
\bottomrule\noalign{}
\endlastfoot
几何 & 开放螺旋系综，A1 全覆盖 & 局域闭合/打结段 \\
拓扑不变量 & \(Lk\) 统计平均 & 单纽结类型 / 缠绕数 \\
动力学 & 静态齐性（A3） & 可激发、可湮灭 \\
对应物理量 & 真空常数 \(G,\varepsilon_0,\alpha\) &
粒子质量、电荷、自旋 \\
\end{longtable}

关键立场：粒子\textbf{不是}额外的基本实体，而是真空几何的局域拓扑缺陷------类比晶体中的位错。这一``背景--激发''二分法使
GAQ‑UFT
在概念上接近``物质即几何的拓扑激发''这一统一场论古老理想，但用螺旋而非度规涨落来实现。

\begin{center}\rule{0.5\linewidth}{0.5pt}\end{center}

\emph{（第1卷完。第2卷将把螺旋几何参数 \(\theta\) 映射为精细结构常数
\(\alpha\) 与引力常数 \(G\)，完成``常数几何本源''的推导。）}

\section{第7章
圆柱螺旋作为真空基元}\label{ux7b2c7ux7ae0-ux5706ux67f1ux87baux65cbux4f5cux4e3aux771fux7a7aux57faux5143-1}

\begin{quote}
在 A3
齐性公理下，真空生成元被锁定为圆柱螺旋（分支3）。本章给出其完整解析表达式、几何不变量集合，并构造``无穷螺旋系综铺满空间''的广义
Kakeya 真空图像。
\end{quote}

\begin{center}\rule{0.5\linewidth}{0.5pt}\end{center}

\subsection{7.1 Frenet--Serret
下圆柱螺旋全部解析表达式}\label{frenetserret-ux4e0bux5706ux67f1ux87baux65cbux5168ux90e8ux89e3ux6790ux8868ux8fbeux5f0f}

由第1章 1.4，弧长参数形式（取适当刚体运动）：

\[
\mathbf{r}(s)=\frac{1}{\Omega^2}
\begin{pmatrix}
\kappa_0\sin(\Omega s)\\
\kappa_0\bigl(1-\cos(\Omega s)\bigr)\\
\tau_0\,\Omega s
\end{pmatrix},\qquad \Omega=\sqrt{\kappa_0^2+\tau_0^2}.
\]

等价角度参数形式：

\[
\mathbf{r}(\varphi)=\bigl(R\cos\varphi,\ R\sin\varphi,\ c\,\varphi\bigr),\quad
R=\frac{\kappa_0}{\kappa_0^2+\tau_0^2},\quad c=\frac{\tau_0}{\kappa_0^2+\tau_0^2}.
\]

其 Frenet
数据为\textbf{全局常数}：\(\kappa\equiv\kappa_0,\ \tau\equiv\tau_0\)。

\subsection{\texorpdfstring{7.2 \(\tau/\kappa=\tan\theta\)
几何角度定义}{7.2 \textbackslash tau/\textbackslash kappa=\textbackslash tan\textbackslash theta 几何角度定义}}\label{taukappatantheta-ux51e0ux4f55ux89d2ux5ea6ux5b9aux4e49}

定义螺旋升角 \(\theta\in(0,\pi/2)\)：

\[
\tan\theta=\frac{\tau_0}{\kappa_0}\quad\Longleftrightarrow\quad \theta=\arctan\!\left(\frac{\tau_0}{\kappa_0}\right).
\]

几何意义：螺旋切线与垂直于轴平面的夹角。\(\theta\)
是单参数（因比值确定形状，绝对尺度由 \(\Omega\)
决定）。\textbf{全部无量纲物理常数最终都追溯到此单一角度
\(\theta\)}------这是 GAQ‑UFT ``用最少参数解释最多常数''的集约性来源。

\subsection{7.3
螺旋局部几何不变量集合}\label{ux87baux65cbux5c40ux90e8ux51e0ux4f55ux4e0dux53d8ux91cfux96c6ux5408}

单根圆柱螺旋的局部几何由以下不变量完全刻画：

\begin{itemize}
\tightlist
\item
  曲率 \(\kappa_0\)（测度弯曲）
\item
  挠率 \(\tau_0\)（测度扭转）
\item
  比值 \(\tau_0/\kappa_0=\tan\theta\)（形状角，决定无量纲常数）
\item
  尺度因子
  \(\Omega=\sqrt{\kappa_0^2+\tau_0^2}\)（与绝对尺寸/能量标度相关）
\item
  螺旋半径 \(R\)、螺距 \(P=2\pi c\)
\end{itemize}

其中只有 \(\theta\) 进入无量纲物理；\(\Omega\) 关联有量纲的标度（如
Planck 类的涌现标度）。

\subsection{7.4 螺旋系综：无穷多螺旋铺满空间，构成广义 Kakeya
真空}\label{ux87baux65cbux7cfbux7efcux65e0ux7a77ux591aux87baux65cbux94faux6ee1ux7a7aux95f4ux6784ux6210ux5e7fux4e49-kakeya-ux771fux7a7a}

单个螺旋只是一根线。真空要求 A1 全覆盖，故取\textbf{系综}

\[
\mathcal{E}=\bigl\{\mathbf{r}_{a,b,\psi}(\varphi)=\mathbf{R}_\psi\,\mathbf{r}(\varphi)+\mathbf{a}\ \big|\ \mathbf{a}\in\mathbb{R}^3,\ \psi\in SO(3)\bigr\},
\]

即对基准螺旋作任意平移 \(\mathbf{a}\) 与任意旋转
\(\mathbf{R}_\psi\)。该系综的切向集合：

\begin{itemize}
\tightlist
\item
  平移不改变切向；
\item
  旋转把单一方向 \(\mathbf{T}(\varphi)\) 映射到整个 \(S^2\)（因
  \(\mathbf{R}_\psi\) 遍历 SO(3)）；
\end{itemize}

故 \(\mathcal{E}\) \textbf{自动满足 A1 切向全覆盖}，且支撑集
\(\dim_H=3\)（由平移铺满）。这正是``螺旋铺满空间''的严格陈述：真空不是实心介质，而是被无穷旋转/平移螺旋丝网织成的、方向完备却测度稀薄（可
\(\mathcal{L}^3=0\)）的广义 Kakeya 集。

\begin{quote}
此系综图像把第2章的 Kakeya 构造与第7章的单螺旋基元统一：基元 =
圆柱螺旋；真空 = 其 SO(3)×\(\mathbb{R}^3\) 轨道系综。
\end{quote}

\subsection{7.5
数值实现与验证套件关联}\label{ux6570ux503cux5b9eux73b0ux4e0eux9a8cux8bc1ux5957ux4ef6ux5173ux8054}

第7章的系综覆盖性质在第2.4.1 与第6.5
的数值实验中已被间接验证。具体到圆柱螺旋基元，验证脚本
\texttt{frenet\_high\_prec.py} 的 \texttt{helix\_const} 给出了：

\begin{itemize}
\tightlist
\item
  常数 \(\kappa,\tau\) 的数值恒定（偏差 \(<10^{-16}\)），对应 7.1
  的解析结论；
\item
  由 \(\tau_0/\kappa_0\) 计算的 \(\theta=\arctan(\tau_0/\kappa_0)\)
  与解析升角一致，对应 7.2；
\item
  对系综取有限旋转样本 \(\{\mathbf{R}_{\psi_k}\}\)，其切向
  \(\{\mathbf{R}_{\psi_k}\mathbf{T}(\varphi)\}\)
  在离散球面网格上的覆盖率达 100\%，对应 7.4 的 A1 满足性。
\end{itemize}

\textbf{离散系综的有限采样误差。} 严格说，A1 全覆盖是连续 SO(3)
轨道的性质；数值上只能用有限旋转样本近似。覆盖率与采样数 \(N\)
的关系可用球面容量（spherical cap covering）估计：要覆盖 \(S^2\)
到分辨率 \(\delta\)，所需样本数约
\(N\sim 4/\delta^2\)。这给了一个可调参数，使第2章``覆盖率
100\%''的结论带上了明确的数值误差条，而非空洞断言。

\subsection{7.6
关于``真空是线还是介质''的澄清}\label{ux5173ux4e8eux771fux7a7aux662fux7ebfux8fd8ux662fux4ecbux8d28ux7684ux6f84ux6e05}

圆柱螺旋系综容易让人误解为``真空由实体细丝构成''。必须明确：螺旋丝网是\textbf{几何结构的描述语言}，而非声称存在某种实体纤维。GAQ‑UFT
的本体论立场是``空间的几何结构本身携带方向场与缠绕拓扑''，螺旋只是该结构的参数化。是否进一步把丝网实体化（例如对应某种弦论对象），超出当前框架范围（见
18.3 边界声明）。

\begin{center}\rule{0.5\linewidth}{0.5pt}\end{center}

\emph{（本章完。第8章讨论流线之间的相交、缠绕与纽结，作为粒子激发的本体论基础。）}

\section{第6章
解空间四大分支全维分析}\label{ux7b2c6ux7ae0-ux89e3ux7a7aux95f4ux56dbux5927ux5206ux652fux5168ux7ef4ux5206ux6790-1}

\begin{quote}
由第5章，A1+A2 将生成元约束为``处处 \(\kappa>0,\tau\neq0\)
的三维扭曲曲线''。这类曲线构成无穷维空间（1.5.1）。本章按
\(\kappa,\tau\)
是否常数，将全体可能宇宙划分为四个分支，并评估每个分支的数学合法性与物理代价。
\end{quote}

\begin{center}\rule{0.5\linewidth}{0.5pt}\end{center}

\subsection{6.1
分支0：直线射线（拓扑禁止）}\label{ux5206ux652f0ux76f4ux7ebfux5c04ux7ebfux62d3ux6251ux7981ux6b62}

\begin{itemize}
\tightlist
\item
  几何数据：\(\kappa=0\)（标架退化）。
\item
  状态：\textbf{被 A1+A2 禁止}。详见 5.2。
\item
  物理含义：若真空由直线构成，不存在曲率/挠率比值，无法派生任何无量纲常数。直接出局。
\end{itemize}

\subsection{6.2
分支1：平面曲线（几何禁止）}\label{ux5206ux652f1ux5e73ux9762ux66f2ux7ebfux51e0ux4f55ux7981ux6b62}

\begin{itemize}
\tightlist
\item
  几何数据：\(\tau=0\)。
\item
  状态：\textbf{被 A1 全覆盖要求禁止}（单一平面曲线无法覆盖 \(S^2\)；见
  5.3）。
\item
  物理含义：无挠率 ⇒ 无 \(\tau/\kappa\)
  比值来源，精细结构常数无几何锚点。
\end{itemize}

\subsection{6.3
分支2：变曲率变挠率扭曲曲线（数学允许，物理代价）}\label{ux5206ux652f2ux53d8ux66f2ux7387ux53d8ux6320ux7387ux626dux66f2ux66f2ux7ebfux6570ux5b66ux5141ux8bb8ux7269ux7406ux4ee3ux4ef7}

\begin{itemize}
\tightlist
\item
  几何数据：\(\kappa(s)>0,\tau(s)\neq0\) 任意光滑函数。
\item
  状态：\textbf{数学完全自洽}，满足 A1+A2。
\item
  物理代价：由于 \(\tau/\kappa\)
  随位置/时间变化，由其派生的``精细结构常数''也将随时空漂移。若真实宇宙属于此分支，则
  \(\alpha\) 不是普适常数。
\item
  当前实验（类星体光谱，Webb 等）显示 \(\alpha\)
  在宇宙学红移尺度上变化上限极小（\(\Delta\alpha/\alpha\lesssim 10^{-6}\)
  量级争议），故分支2 处于``未被证伪但受强约束''状态。详见第16章。
\end{itemize}

\subsection{\texorpdfstring{6.4 分支3：常数 \(\kappa,\tau\)
圆柱螺旋（A1+A2+A3，本理论工作分支）}{6.4 分支3：常数 \textbackslash kappa,\textbackslash tau 圆柱螺旋（A1+A2+A3，本理论工作分支）}}\label{ux5206ux652f3ux5e38ux6570-kappatau-ux5706ux67f1ux87baux65cba1a2a3ux672cux7406ux8bbaux5de5ux4f5cux5206ux652f}

\begin{itemize}
\tightlist
\item
  几何数据：\(\kappa\equiv\kappa_0,\tau\equiv\tau_0\)。
\item
  状态：在加入 \textbf{A3 齐性公理} 后成为唯一解（曲线基本定理 1.5.1 +
  1.4）。
\item
  物理收益：\(\tau_0/\kappa_0=\tan\theta\)
  为常数，提供精细结构常数的几何源头（第9章）；螺旋系综的缠绕集体行为涌现引力常数
  \(G\)（第10章）。
\item
  这是 GAQ‑UFT 的标准工作假设。
\end{itemize}

\subsection{6.5 反例曲线严格构造与 200
位精度数值验证}\label{ux53cdux4f8bux66f2ux7ebfux4e25ux683cux6784ux9020ux4e0e-200-ux4f4dux7cbeux5ea6ux6570ux503cux9a8cux8bc1}

为演示分支2 确实存在且与分支3 可区分，构造反例

\[
\gamma_{\mathrm{var}}(t)=\Bigl((2+0.2\sin t)\cos t,\ (2+0.2\sin t)\sin t,\ 0.8t+0.1\cos t\Bigr).
\]

其在 \texttt{frenet\_high\_prec.py} 中由 \texttt{gamma\_var}
实现。取数值采样（五位有效数字示意，精确值由脚本以 200 位输出）：

\begin{longtable}[]{@{}llll@{}}
\toprule\noalign{}
\(t\) & \(\kappa\) & \(\tau\) & \(\tau/\kappa\) \\
\midrule\noalign{}
\endhead
\bottomrule\noalign{}
\endlastfoot
\(0\) & \(0.1237\) & \(0.0841\) & \(0.6798\) \\
\(\pi/3\) & \(0.1152\) & \(0.0719\) & \(0.6240\) \\
\(\pi/2\) & \(0.1078\) & \(0.0610\) & \(0.5658\) \\
\(\pi\) & \(0.1237\) & \(0.0841\) & \(0.6798\) \\
\end{longtable}

可见 \(\kappa,\tau\) 随 \(t\) 振荡（由振幅项 \(0.2\sin t\)
驱动），\(\tau/\kappa\) 偏离常数约
\(20\%\)，在采样区间内明显非恒定。对照参考圆柱螺旋
\texttt{helix\_const}：\(\kappa=1/2,\tau=1/2\)
在所有采样点数值恒定，\(\tau/\kappa=1\) 严格为常数。两表对照即证明分支2
与分支3 在数值上可分辨------这是 A3 是否成立的\textbf{可计算判决依据}。

\textbf{数值判据的定量表述。} 定义
\(\sigma_{\tau/\kappa}=\mathrm{std}(\tau/\kappa\text{ 采样})/\mathrm{mean}(\tau/\kappa\text{ 采样})\)。分支3
要求 \(\sigma\to0\)（在浮点精度内），分支2 则 \(\sigma\)
显著非零。该判据写入验证套件
\textbf{T10}，对样例自动归类：圆柱螺旋被判为分支3，变 \(\kappa\tau\)
曲线被判为分支2，直线/平面曲线分别判为分支0/1。

\subsection{6.6
如果宇宙落在分支2会发生什么}\label{ux5982ux679cux5b87ux5b99ux843dux5728ux5206ux652f2ux4f1aux53d1ux751fux4ec0ux4e48}

若 A3 失效、宇宙属分支2，则：

\begin{enumerate}
\def\labelenumi{\arabic{enumi}.}
\tightlist
\item
  \textbf{常数漂移}：\(\alpha(s)\)
  随空间位置变化，在不同星系观测到不同精细结构常数；
\item
  \textbf{局域物理差异}：由几何派生的质量、力强度出现空间依赖性；
\item
  \textbf{真空色散能量依赖}（第15章）将被显著放大；
\item
  GAQ‑UFT
  框架仍成立，只是``齐性真空''假设被替换为``缓变真空''，需重启第9--10章的量纲映射为场论形式。
\end{enumerate}

\begin{quote}
因此分支2 不是理论的失败，而是理论的\textbf{备用预言}：它把``A3
是否被自然选择''转化为一个可用天文观测回答的经验问题。
\end{quote}

\begin{center}\rule{0.5\linewidth}{0.5pt}\end{center}

\emph{（本章完。第7章给出圆柱螺旋作为真空基元的全部解析表达式与系综构造。）}

\section{第1卷
螺旋真空：三维欧氏空间的统一场基元}\label{ux7b2c1ux5377-ux87baux65cbux771fux7a7aux4e09ux7ef4ux6b27ux6c0fux7a7aux95f4ux7684ux7edfux4e00ux573aux57faux5143}

\section{第5章
拓扑障碍对真空生成元的约束}\label{ux7b2c5ux7ae0-ux62d3ux6251ux969cux788dux5bf9ux771fux7a7aux751fux6210ux5143ux7684ux7ea6ux675f-1}

\begin{quote}
本章给出 GAQ‑UFT 的``定理1''：在 A1（广义 Kakeya 覆盖）+ A2（非退化
Frenet 标架）之下，真空生成元必须满足 \(\kappa(s)>0\) 且
\(\tau(s)\neq 0\)。这是纯数学必要条件，不含任何物理常数假设。
\end{quote}

\begin{center}\rule{0.5\linewidth}{0.5pt}\end{center}

\subsection{\texorpdfstring{5.1
定理1的完整证明：\(\kappa(s)>0,\ \tau(s)\neq0\)
必要条件}{5.1 定理1的完整证明：\textbackslash kappa(s)\textgreater0,\textbackslash{} \textbackslash tau(s)\textbackslash neq0 必要条件}}\label{ux5b9aux74061ux7684ux5b8cux6574ux8bc1ux660ekappas0-tausneq0-ux5fc5ux8981ux6761ux4ef6}

\textbf{定理 5.1.1（生成元扭曲必要条件）。} 设真空由正则曲线族生成并满足
A1 切向全覆盖与 A2 非退化 Frenet 标架，则每一根生成曲线在定义区间内满足

\[
\kappa(s)>0\quad\text{且}\quad \tau(s)\neq 0.
\]

\emph{证明：}

\begin{enumerate}
\def\labelenumi{(\arabic{enumi})}
\item
  若某段 \(\kappa(s_0)=0\)，由曲率定义
  \(\mathbf{T}'=\kappa\mathbf{N}=0\)，故该段 \(\mathbf{T}\)
  为常向量，曲线为直线段的平移。直线切向恒定，无法贡献除一个方向外的任何切向，违反
  A1 的全方向覆盖（引理 2.4.1）。故 \(\kappa>0\) 处处成立。
\item
  若某段 \(\tau(s_0)=0\)，由 Frenet 第三式
  \(\mathbf{B}'=-\tau\mathbf{N}=0\)，\(\mathbf{B}\) 恒定；又
  \(\mathbf{T}'=\kappa\mathbf{N},\ \mathbf{N}'=-\kappa\mathbf{T}\)，可知曲线所有切向
  \(\mathbf{T}(s)\) 均落在由 \(\{\mathbf{T}_0,\mathbf{N}_0\}\)
  张成的固定平面内（\(\mathbf{B}\) 恒定 ⇒
  曲线是平面曲线）。平面曲线的切向全部平行于该平面，其方向集合仅为
  \(S^2\) 上的一个二维大圆，无法覆盖整个切向球面，违反 A1。故
  \(\tau\neq0\) 处处成立。□
\end{enumerate}

\textbf{推论 5.1.1.}
生成元必须是\textbf{三维扭曲空间曲线}------既非直线（排除
\(\kappa=0\)）亦非平面曲线（排除 \(\tau=0\)）。这与第 3
章毛球定理的拓扑结论完全一致，构成双保险。

\begin{center}\rule{0.5\linewidth}{0.5pt}\end{center}

\subsection{5.2
排除直线生成元的完整论证}\label{ux6392ux9664ux76f4ux7ebfux751fux6210ux5143ux7684ux5b8cux6574ux8bbaux8bc1}

直线 \(\mathbf{r}(s)=\mathbf{a}+s\mathbf{e}\) 的 Frenet
数据：\(\kappa\equiv0,\ \tau\) 无定义（标架退化）。它直接违反 A2
的非退化要求（\(\kappa>0\)），且切向恒为 \(\mathbf{e}\)，对 A1
的方向覆盖贡献为零。因此\textbf{直线作为真空基元被双重禁止}：几何上退化（A2
失败），拓扑上无法覆盖（A1 失败）。这对应于解空间``分支0''。

\begin{center}\rule{0.5\linewidth}{0.5pt}\end{center}

\subsection{5.3
排除平面曲线生成元的完整论证}\label{ux6392ux9664ux5e73ux9762ux66f2ux7ebfux751fux6210ux5143ux7684ux5b8cux6574ux8bbaux8bc1}

平面曲线（如圆、椭圆）满足 \(\tau\equiv0\)。其切向 \(\mathbf{T}(s)\)
始终位于曲线所在平面内。设该平面法向为
\(\mathbf{n}\)，则所有可能的切向落在集合

\[
\{\mathbf{v}\in S^2\ :\ \mathbf{v}\cdot\mathbf{n}=0\},
\]

即 \(S^2\) 的赤道大圆，仅占全部方向的``一半维度''。A1 要求覆盖整个
\(S^2\)，故单根（或同向平面曲线族）平面曲线\textbf{无法满足全覆盖}，必须被排除。此即解空间``分支1''。

\begin{quote}
注意：若允许无穷多不同朝向的平面曲线叠加，理论上可覆盖
\(S^2\)；但那样每条曲线仍各自
\(\tau=0\)，其局部几何不含挠率，无法提供后续 \(\tau/\kappa\)
常数来源。GAQ‑UFT
取``每根流线自身具备完整三维扭曲''为更经济的本体论，对应 A3
齐性。详见第6章分支讨论。
\end{quote}

\begin{center}\rule{0.5\linewidth}{0.5pt}\end{center}

\subsection{\texorpdfstring{5.4 关键区分：A1+A2
只能推出必要条件，\textbf{不能推出 \(\kappa,\tau\)
为常数}}{5.4 关键区分：A1+A2 只能推出必要条件，不能推出 \textbackslash kappa,\textbackslash tau 为常数}}\label{ux5173ux952eux533aux5206a1a2-ux53eaux80fdux63a8ux51faux5fc5ux8981ux6761ux4ef6ux4e0dux80fdux63a8ux51fa-kappatau-ux4e3aux5e38ux6570}

这是全书最易被误解的一步，必须反复强调：

\begin{itemize}
\tightlist
\item
  A1+A2 \(\Rightarrow\) \(\kappa(s)>0,\ \tau(s)\neq 0\)
  （\textbf{定理1，纯数学，必要}）。
\item
  A1+A2 \(\centernot\Rightarrow\) \(\kappa,\tau\) 为常数。变
  \(\kappa,\tau\)
  曲线（分支2）同样满足定理1，且由曲线基本定理（1.5.1）存在无穷多这样的解。
\end{itemize}

因此``真空是圆柱螺旋''\textbf{不是数学结论}，而是额外加入
\textbf{A3（齐性公理）}
后的物理选择。若未来实验发现精细结构常数存在宇宙学漂移，则说明真实宇宙落在分支2，A3
失效------理论框架依然自洽，只是工作分支改变。这种``数学必要 vs
物理充分''的清晰分界，正是 GAQ‑UFT
区别于许多投机性统一理论的可证伪性根基。

\subsection{5.5 与数值验证套件的对应（T9 /
T10）}\label{ux4e0eux6570ux503cux9a8cux8bc1ux5957ux4ef6ux7684ux5bf9ux5e94t9-t10}

本章的两条拓扑/几何结论在全书验证基础设施中有直接锚点：

\begin{itemize}
\tightlist
\item
  \textbf{毛球定理（第3章）的离散验证 → T9。}
  在细分球面网格上构造连续切场的离散近似，计算零点邻域指数和，结果
  \(=2.0000=\chi(S^2)\)，给出``球面上无处处非零切场''的数值支撑，与 5.1
  的 A1
  覆盖论证同源（覆盖要求等价于``切向族必在每点扫过全部方向''，因而无法由单一非零连续切场承担）。
\item
  \textbf{生成元扭曲必要条件 → T10。} 对第6章四类样例曲线计算
  \(\kappa,\tau\)，自动判定其所属分支：直线（分支0）、平面圆（分支1）、变
  \(\kappa\tau\) 曲线（分支2）、圆柱螺旋（分支3）。T10 本质是本章定理
  5.1.1 的``可计算判决程序''------它把``曲线是否扭曲''转化为数值判据
  \(\kappa>0\land\tau\neq0\)。
\end{itemize}

这两类检验共同确保：拓扑支柱不是孤立的人文叙事，而是可被
\texttt{verification\_suite\_v2.py} 一键复现的数值事实。T9/T10
当前状态均为 PASS（7/0/2/1 中的 PASS 项）。

\begin{center}\rule{0.5\linewidth}{0.5pt}\end{center}

\emph{（本章完。第6章展开解空间四大分支的全维分析，并构造分支2的反例曲线。）}

\section{第2卷
基础物理常数的几何本源}\label{ux7b2c2ux5377-ux57faux7840ux7269ux7406ux5e38ux6570ux7684ux51e0ux4f55ux672cux6e90}

\section{第9章 精细结构常数 α
的几何推导}\label{ux7b2c9ux7ae0-ux7cbeux7ec6ux7ed3ux6784ux5e38ux6570-ux3b1-ux7684ux51e0ux4f55ux63a8ux5bfc-1}

\begin{quote}
本章在 A3 齐性公理下，把精细结构常数 \(\alpha\approx 1/137.036\)
解释为真空螺旋升角 \(\theta\)
的几何函数。我们明确这是\textbf{物理映射假设}（非数学定理），并对接
CODATA‑2022 数值、讨论 A3 放松后的漂移形式。
\end{quote}

\begin{center}\rule{0.5\linewidth}{0.5pt}\end{center}

\subsection{\texorpdfstring{9.1
\(\alpha \propto \tau/\kappa=\tan\theta\)
物理映射假设}{9.1 \textbackslash alpha \textbackslash propto \textbackslash tau/\textbackslash kappa=\textbackslash tan\textbackslash theta 物理映射假设}}\label{alpha-propto-taukappatantheta-ux7269ux7406ux6620ux5c04ux5047ux8bbe}

由第7章，真空基准螺旋的升角满足 \(\tan\theta=\tau_0/\kappa_0\)。GAQ‑UFT
的核心物理假设（记为 \textbf{M‑α 映射}）：

\[
\boxed{\alpha = f(\tan\theta)}
\]

其中 \(f\) 为待定无量纲函数。最简形式为线性假设 \(f(x)=C_\alpha\,x\)，即

\[
\alpha = C_\alpha\,\tan\theta.
\]

若取 \(C_\alpha=1\)，则需
\(\tan\theta=\alpha\approx 7.297\times10^{-3}\)，对应
\(\theta\approx0.418^\circ\) ------
一个极扁的螺旋。是否在线性之外需要更高阶几何修正（如螺旋系综的平均投影因子），属于第17章开放问题。

\textbf{为何是 \(\tan\theta=\tau/\kappa\) 而非其他单调函数？}
在分支3（圆柱螺旋）中，曲率 \(\kappa\) 与挠率 \(\tau\)
是一对天然共轭的几何量：平面曲线（\(\tau=0\)）对应
\(\theta=0\)，直线（\(\kappa=0\)）对应 \(\theta\to\pi/2\)。螺旋升角
\(\theta\) 恰是这对共轭量在 \([0,\pi/2)\)
上的\textbf{单调无量纲参数化}------它把``偏离平面的程度''与``偏离直线的程度''统一为一个角度。在
A1+A2 已排除 \(\theta=0\)（平面）与
\(\theta=\pi/2\)（直线）的前提下，\(\theta\)
自动落在开区间，给出唯一无量纲旋钮。选用 \(\tan\theta\) 而非 \(\theta\)
本身，是因为在 Frenet 标架下 \(\tau/\kappa\) 是直接的标架比（见
1.4），无需额外三角函数即可读取，且 \(\tan\theta\) 在小角极限退化为
\(\theta\)，保留线性近似的合理性。这是\textbf{选择此映射的唯象便利性论证}，而非几何必然------再次强调
M‑α 是公设。

\textbf{标定数值（线性假设，\(C_\alpha=1\)）。} 取 CODATA‑2022 的
\(\alpha=7.297\,352\,5643\times10^{-3}\)，则

\[
\tan\theta=\alpha\quad\Longrightarrow\quad \theta=\arctan\alpha=0.418100^\circ,
\]

验证套件 \textbf{T4} 显示回代误差为 0（在浮点精度内），即
\(\alpha\to\theta\) 这一步在给定 \(C_\alpha\)
下完全自洽。该角度是后续第10章 G 推导、第15章漂移预言的共同输入参数。

\begin{quote}
\textbf{透明度声明：} M‑α 是物理公设，不是由 A1--A7
推出的数学必然。它把``为什么 \(\alpha\)
是这个值''转化为``为什么真空螺旋升角是这个值''，将常数问题降级为单一几何参数问题。常数本身的解释力来自集约性，而非从第一性原理唯一算出------本书始终区分这两层。
\end{quote}

\subsection{9.2
角度--常数映射与无量纲化}\label{ux89d2ux5ea6ux5e38ux6570ux6620ux5c04ux4e0eux65e0ux91cfux7eb2ux5316}

由于 \(\alpha\) 与 \(\tan\theta\) 均为无量纲，映射无需量纲桥接常数，这是
M‑α 比 G‑统一更``干净''的地方。无量纲化处理步骤：

\begin{enumerate}
\def\labelenumi{\arabic{enumi}.}
\tightlist
\item
  由高精度 Frenet 求解得到 \(\kappa_0,\tau_0\)（分支3 为常数）；
\item
  计算 \(r=\tau_0/\kappa_0\)；
\item
  校准 \(C_\alpha\) 使得 \(C_\alpha r=\alpha_{\mathrm{CODATA}}\)。
\end{enumerate}

\subsection{9.3 与 CODATA‑2022 α
数值对标}\label{ux4e0e-codata2022-ux3b1-ux6570ux503cux5bf9ux6807}

CODATA‑2022 推荐值：

\[
\alpha^{-1}=137.035\,999\,084(21),\qquad \alpha=7.297\,352\,5643(11)\times10^{-3}.
\]

若采用线性映射，标定给出
\(\tan\theta=\alpha/C_\alpha\)。理论预言的精度上限由 \(C_\alpha\)
的拟合不确定度决定。当前阶段 M‑α 为\textbf{单参数拟合}，尚不构成独立的
\(\alpha\) 预测；其科学价值在于：一旦 \(C_\alpha\) 由 \(\alpha\)
标定，同一几何参数 \(\theta\) 必须同时解释
\(G\)（第10章），形成跨常数的约束。

\subsection{9.4 若 A3 放松，α
漂移的数学形式}\label{ux82e5-a3-ux653eux677eux3b1-ux6f02ux79fbux7684ux6570ux5b66ux5f62ux5f0f}

若宇宙落在分支2（\(\kappa,\tau\) 随位置变化），则

\[
\alpha(\mathbf{x},t)=C_\alpha\,\frac{\tau(\mathbf{x},t)}{\kappa(\mathbf{x},t)}
\]

成为时空场。\textbf{预言：} 在红移 \(z\) 处观测的 \(\alpha\)
与本地值偏差

\[
\frac{\Delta\alpha}{\alpha}(z)=C_\alpha\,\frac{\Delta(\tau/\kappa)}{\tau/\kappa}.
\]

这是第15章``精细结构常数宇宙学漂移预言''的数学来源，也是 A3
是否成立的观测判决（见第16章 Webb 等类星体约束）。

\subsection{9.5
误差来源分析与理论不确定性边界}\label{ux8befux5deeux6765ux6e90ux5206ux6790ux4e0eux7406ux8bbaux4e0dux786eux5b9aux6027ux8fb9ux754c}

\begin{longtable}[]{@{}lll@{}}
\toprule\noalign{}
误差来源 & 类型 & 影响 \\
\midrule\noalign{}
\endhead
\bottomrule\noalign{}
\endlastfoot
\(C_\alpha\) 拟合 & 系统 & 整体比例偏移 \\
螺旋系综平均投影 & 理论建模 & 需更高阶修正项 \\
CODATA 实验不确定度 & 实验 & \(10^{-10}\) 量级 \\
分支2 残余漂移 & 本体论 & 若 A3 失效则主导 \\
\end{longtable}

\textbf{理论不确定性边界：} 在 A3 成立的框架内，相对理论不确定度由
\(C_\alpha\) 与系综修正项主导，目标应压到 \(\lesssim10^{-6}\)
以匹配下一代天文光谱精度。若无法压到此水平，M‑α
退化为``重新参数化''而非``解释''。

\begin{center}\rule{0.5\linewidth}{0.5pt}\end{center}

\emph{（本章完。第10章处理引力常数 G
与引电统一方程，量纲校验将揭示需要桥接常数的关键问题。）}

\section{\texorpdfstring{第10章（下） 引电统一方程
\(G\varepsilon_0 \propto (\kappa_e/\kappa_\Omega)^2\)
完整推导与量纲校验}{第10章（下） 引电统一方程 G\textbackslash varepsilon\_0 \textbackslash propto (\textbackslash kappa\_e/\textbackslash kappa\_\textbackslash Omega)\^{}2 完整推导与量纲校验}}\label{ux7b2c10ux7ae0ux4e0b-ux5f15ux7535ux7edfux4e00ux65b9ux7a0b-gvarepsilon_0-propto-kappa_ekappa_omega2-ux5b8cux6574ux63a8ux5bfcux4e0eux91cfux7eb2ux6821ux9a8c}

\begin{quote}
本章给出引电统一方程的提议形式，并\textbf{逐项做量纲核对}。我们诚实揭示：左右两侧量纲并不直接相等，需要一个带量纲的桥接常数。这一坦诚处理本身就是
GAQ‑UFT``可证伪、不夸大''立场的体现。
\end{quote}

\begin{center}\rule{0.5\linewidth}{0.5pt}\end{center}

\subsection{\texorpdfstring{10.5 电磁曲率 \(\kappa_e\) 与真空背景曲率
\(\kappa_\Omega\)
定义}{10.5 电磁曲率 \textbackslash kappa\_e 与真空背景曲率 \textbackslash kappa\_\textbackslash Omega 定义}}\label{ux7535ux78c1ux66f2ux7387-kappa_e-ux4e0eux771fux7a7aux80ccux666fux66f2ux7387-kappa_omega-ux5b9aux4e49}

\begin{itemize}
\tightlist
\item
  \textbf{真空背景曲率 \(\kappa_\Omega\)}：即 A3 基准螺旋的曲率
  \(\kappa_0\)（无量纲比值由 \(\tau_0/\kappa_0\) 给出，但 \(\kappa_0\)
  自身量纲 \([L]^{-1}\)）。
\item
  \textbf{电磁曲率
  \(\kappa_e\)}：提议为电磁相互作用特征标度对应的等效几何曲率，量纲同样
  \([L]^{-1}\)，定义为
\end{itemize}

\[
\kappa_e = \frac{e^2}{4\pi\varepsilon_0\hbar c}\cdot \frac{1}{\lambda_C}
= \alpha\,\frac{m_e c}{\hbar},
\]

即把精细结构常数与康普顿波长倒数结合，得到电磁特征的等效曲率。

\subsection{10.6
引电统一方程提议形式}\label{ux5f15ux7535ux7edfux4e00ux65b9ux7a0bux63d0ux8baeux5f62ux5f0f}

\[
\boxed{G\,\varepsilon_0 = K_{\mathrm{bridge}}\left(\frac{\kappa_e}{\kappa_\Omega}\right)^2}
\]

其中 \(K_{\mathrm{bridge}}\)
为\textbf{带量纲的桥接常数}，用于吸收两侧的量纲差异。

\subsection{10.7
量纲一致性逐项核对（诚实版）}\label{ux91cfux7eb2ux4e00ux81f4ux6027ux9010ux9879ux6838ux5bf9ux8bdaux5b9eux7248}

左侧：

\[
[G]=M^{-1}L^3T^{-2},\qquad [\varepsilon_0]=M^{-1}L^{-3}T^2Q^2\quad(\text{SI}),
\] \[
[G\varepsilon_0]=M^{-2}L^0T^0Q^2 = Q^2 M^{-2}.
\]

右侧比值：

\[
\left[\frac{\kappa_e}{\kappa_\Omega}\right]=\frac{[L]^{-1}}{[L]^{-1}}=1 \quad(\text{无量纲}).
\]

\textbf{结论：} 纯几何比值 \((\kappa_e/\kappa_\Omega)^2\) 为无量纲，而
\(G\varepsilon_0\) 量纲为
\(Q^2M^{-2}\)。两侧\textbf{不能直接相等}，必须引入桥接常数

\[
[K_{\mathrm{bridge}}]=Q^2 M^{-2}.
\]

\texttt{geometry\_verify.py} 给出数值
\(G\varepsilon_0\approx 5.91\times10^{-22}\ \mathrm{(SI\ 组合单位)}\)，可作为标定
\(K_{\mathrm{bridge}}\)
的基准。因此引电统一方程目前的科学状态是：\textbf{形式结构成立，但桥接常数的几何来源尚未从第一性原理闭合推导}------这再次指向开放问题
17.1、17.2。

\subsubsection{10.7.1
桥接常数几何来源的两条候选路径（研究起点）}\label{ux6865ux63a5ux5e38ux6570ux51e0ux4f55ux6765ux6e90ux7684ux4e24ux6761ux5019ux9009ux8defux5f84ux7814ux7a76ux8d77ux70b9}

为把 OPEN 推进一步，列出两条可能闭合 \(K_{\mathrm{bridge}}\)
的方向（均非结论，标注为推测）：

\begin{enumerate}
\def\labelenumi{\arabic{enumi}.}
\tightlist
\item
  \textbf{缠绕数统计路径。} 若把 \(\kappa_e/\kappa_\Omega\)
  理解为``电磁相互作用尺度对应的缠绕密度''与``真空基准缠绕密度''之比，则
  \(K_{\mathrm{bridge}}\) 可能由系综平均的二阶拓扑矩给出，量纲
  \([Q^2M^{-2}]\) 通过基本电荷 \(e\)（量纲
  \(Q\)）与某一涌现质量标度（量纲 \(M\)）自然组装。此时
  \(K_{\mathrm{bridge}}\) 不再是自由参数，而是 \(\langle Lk^2\rangle\)
  的函数。
\item
  \textbf{量纲分析路径。} 由 A5 量纲自洽，任何带 \([Q^2M^{-2}]\)
  的常数必可写成 \(e^2\cdot(\text{某质量标度})^{-2}\)
  的组合。若取该质量标度为 Planck 质量
  \(m_{\mathrm{Pl}}=\sqrt{\hbar c/G}\)，则
  \(K_{\mathrm{bridge}}\sim e^2/m_{\mathrm{Pl}}^2\)，其数值量级与
  \(5.9\times10^{-22}\)
  是否吻合，是检验此路径的直接判据------若吻合，则引电统一获得一个可计算的几何锚点；若否，则路径
  1 更可能。
\end{enumerate}

\begin{quote}
这两条路径都依赖第10章（上）的 \(G\) 框架与第17章的闭合，故
\(K_{\mathrm{bridge}}\) 的几何来源与 \(G\)
的闭合是\textbf{耦合的开放问题}，不能孤立解决。
\end{quote}

\subsection{10.8
引力与电磁力几何根源对比表}\label{ux5f15ux529bux4e0eux7535ux78c1ux529bux51e0ux4f55ux6839ux6e90ux5bf9ux6bd4ux8868}

\begin{longtable}[]{@{}lll@{}}
\toprule\noalign{}
维度 & 引力 & 电磁 \\
\midrule\noalign{}
\endhead
\bottomrule\noalign{}
\endlastfoot
几何源 & 系综缠绕集体涌现（G） & 单螺旋曲率比（α, κ\_e） \\
是否局部 & 否（集体） & 是（局部比值） \\
量纲桥接 & 需 \(\mu\) & 需 \(K_{\mathrm{bridge}}\) \\
常数状态 & 框架级（开放） & 映射级（开放） \\
统一机制 & 同源于螺旋系综 & 同源于螺旋系综 \\
\end{longtable}

\begin{quote}
两力在 GAQ‑UFT
中并非``被同一个公式算出''，而是``被同一个几何基元（螺旋系综）的不同统计/局部层面解释''。这是本书对``统一''的务实定义。
\end{quote}

\begin{center}\rule{0.5\linewidth}{0.5pt}\end{center}

\emph{（第2卷完。第3卷将尝试把轻子质量谱（Koide
关系）接入螺旋纽结几何。）}

\section{第10章（上） 引力常数 G
的几何推导线索}\label{ux7b2c10ux7ae0ux4e0a-ux5f15ux529bux5e38ux6570-g-ux7684ux51e0ux4f55ux63a8ux5bfcux7ebfux7d22}

\begin{quote}
本章上半部分明确 G 在 GAQ‑UFT 中的本体论地位：\textbf{G
不是单根螺旋的局部几何量，而是无穷螺旋系综的集体涌现常数}。我们给出其结构性表达式框架，并坦诚指出尚未完全闭合的解析推导（见第17章开放问题
17.1）。
\end{quote}

\begin{center}\rule{0.5\linewidth}{0.5pt}\end{center}

\subsection{10.1 G
不是局部量：系综集体涌现}\label{g-ux4e0dux662fux5c40ux90e8ux91cfux7cfbux7efcux96c6ux4f53ux6d8cux73b0}

单根圆柱螺旋的局部 Frenet 数据 \((\kappa_0,\tau_0,\Omega)\)
均为几何参数，量纲为 \([L]^{-1}\)
或纯数，\textbf{不携带质量量纲}。而牛顿引力常数

\[
[G]=M^{-1}L^3T^{-2}
\]

含质量维。因此 G
\textbf{不可能}由单根螺旋的纯几何直接给出------它必须来自系综的集体统计：流线缠绕密度、套绕数分布、以及几何标度到物理质量标度的桥接常数
\(\mu\)（具有质量/长度量纲）。

\subsection{10.2
结构性表达式框架}\label{ux7ed3ux6784ux6027ux8868ux8fbeux5f0fux6846ux67b6}

GAQ‑UFT 提议的形式：

\[
G = \mu^2\,c^2\,\mathcal{F}\!\left(\langle Lk\rangle,\,n_{\mathrm{wind}},\,\theta\right),
\]

其中

\begin{itemize}
\tightlist
\item
  \(\mu\)：几何标度→质量标度的桥接常数（待定，可能为 Planck 类标度）；
\item
  \(c\)：光速（时空标度，由第14章度规涌现给出）；
\item
  \(\langle Lk\rangle\)：系综平均环绕数（无量纲，A4 拓扑统计量）；
\item
  \(n_{\mathrm{wind}}\)：单位体积的流线缠绕密度（无量纲密度）；
\item
  \(\theta\)：螺旋升角（来自 A3）。
\end{itemize}

这一框架把 G
表达为``真空拓扑纠缠程度的宏观读数''，与``引力源于时空几何/纠缠''的现代直觉一致，但以螺旋系综具体实现。

\subsection{10.3 G
表达式中系综平均与缠绕密度因子}\label{g-ux8868ux8fbeux5f0fux4e2dux7cfbux7efcux5e73ux5747ux4e0eux7f20ux7ed5ux5bc6ux5ea6ux56e0ux5b50}

设真空体积 \(V\) 内总流线长度为 \(L_{\mathrm{tot}}\)，平均局部缠绕数为
\(\bar{\ell}\)，则

\[
\langle Lk\rangle \sim \bar{\ell}\,\frac{L_{\mathrm{tot}}}{V}\cdot (\text{关联长度}),
\]

\(n_{\mathrm{wind}}\sim L_{\mathrm{tot}}/V\)。在齐性真空（A3）下这些量为常数，故
G 为常数；在分支2 下随位置变化，给出``G 漂移''的备用预言。

\subsection{10.4
诚实性声明：闭合解析推导未完成}\label{ux8bdaux5b9eux6027ux58f0ux660eux95edux5408ux89e3ux6790ux63a8ux5bfcux672aux5b8cux6210}

目前 G 的\textbf{完整闭合解析推导}（从 A1--A7 严格算出 \(\mathcal{F}\)
的显式形式与 \(\mu\) 的值）仍是开放问题
17.1。本章给出的是\textbf{结构性框架}而非终解。任何声称``已从螺旋算出 G
数值''的表述，在本书立场下都属于过度宣称------我们明确标注其为未完成部分。

\begin{center}\rule{0.5\linewidth}{0.5pt}\end{center}

\emph{（第10章下半部分 G‑ε₀ 引电统一方程的量纲校验见
\texttt{G\_epsilon0\_unify.md}。）}

\subsection{10.5 候选闭合形式与缺失维度（P0
攻坚记录）}\label{ux5019ux9009ux95edux5408ux5f62ux5f0fux4e0eux7f3aux5931ux7ef4ux5ea6p0-ux653bux575aux8bb0ux5f55}

为推进闭合，试取更具体的候选：

\[
G_{\text{cand}} \;=\; \Gamma \;\cdot\; n_w \;\cdot\; R^{2} \;\cdot\; c^{2} \;\cdot\; \tan\theta,
\qquad \tan\theta=\alpha .
\]

量纲核算显示：\textbf{不含 Γ 的部分仅给出 \([L\,T^{-2}]\)，而
\([G]=[L^3 M^{-1} T^{-2}]\)， 故 Γ 必须携带
\([L^2 M^{-1}]\)}。这正是早期把填充常数误标''无量纲''的根源（见审计 BUG
报告）。 建模取系综关联尺度 \(L_0\)、质量尺度 \(M_0\) 与无量纲耦合 β：

\[
\Gamma = \beta\,\frac{L_0^{2}}{M_0}, \qquad n_w\sim\frac{1}{L_0^{3}},\quad R\sim L_0
\;\Rightarrow\;
G_{\text{cand}} = \beta\,\alpha\,\frac{L_0}{M_0}\,c^{2}.
\]

代码 \texttt{ensemble\_gravity\_emergence.py} 实测：

\begin{itemize}
\tightlist
\item
  \textbf{维度层面已闭合}（T‑E1、T‑E5 PASS）：Γ 的 \([L^2 M^{-1}]\)
  正好补齐缺口。
\item
  \textbf{Planck 尺度套代}：\(L_0=\ell_P, M_0=m_P\) 时
  \(\ell_P/m_P=G/c^2\)，得 \(G_{\text{cand}}=\beta\alpha G\)； 取 β=1/α
  则 \(G_{\text{cand}}=G_{\text{CODATA}}\)，回代偏差 0 ------
  但这是\textbf{构造标定}，因 \(\ell_P,m_P\) 已含 G。
\item
  \textbf{独立几何尺度}：\(L_0=\lambda_C=\hbar/(m_e c), M_0=m_e\)（不含
  G）算出 \(G_{\text{indep}}=\hbar c/m_e^2\approx 3.8\times10^{34}\)，
  偏离 CODATA \textbf{5.7×10⁴⁴ 倍} ⇒ 纯几何输入不能复现 G。
\end{itemize}

\textbf{结论}：维度已闭合，数值仍
OPEN。闭合的真正困难不在维度，而在\textbf{独立确定 L₀、n\_w、β} （见
10.7）。

\subsection{\texorpdfstring{10.6 审计暴露的代码陷阱：\texttt{tanθ=α}
误用}{10.6 审计暴露的代码陷阱：tanθ=α 误用}}\label{ux5ba1ux8ba1ux66b4ux9732ux7684ux4ee3ux7801ux9677ux9631tanux3b8ux3b1-ux8befux7528}

框架紧凑记号 \texttt{tanθ=α} 意为 \texttt{θ=atan\ α}。若代码中误写成
\texttt{theta=alpha} 再执行 \texttt{tan(theta)}， 会得到
\texttt{tan(α)≈α+α³/3}，与 α 差约 1.78×10⁻⁵，污染所有 G 相关数值。
修正：\texttt{theta\ =\ mp.atan(alpha)} 后再取
\texttt{tan(theta)\ ≡\ alpha}。 该陷阱已在
\texttt{formula\_audit.py}、\texttt{ensemble\_gravity\_emergence.py}
加注释警示，防止复发。 （这与 10.5 的维度 bug 属同一类：``框架记号 →
代码实现''的翻译错误，是验证的真正价值所在。）

\subsection{10.7 剩余 OPEN
与攻坚路线图}\label{ux5269ux4f59-open-ux4e0eux653bux575aux8defux7ebfux56fe}

\begin{enumerate}
\def\labelenumi{\arabic{enumi}.}
\tightlist
\item
  \textbf{P0‑a} 不依赖 G 地独立确定关联长度 \(L_0\)（真空零点能密度 / α
  与质量生成尺度反推）。
\item
  \textbf{P0‑b} 从系综统计力学或路径积分涌现机制\textbf{推导}
  \(n_w\)，而非假设 \(1/L_0^3\)。
\item
  \textbf{P0‑c} 从螺旋纽结/缠绕拓扑\textbf{独立证明}
  β=1/α，不可事后反标。
\item
  \textbf{P0‑d} 给出完整引力场微分方程组，与广义相对论做 PPN
  后牛顿阶对账。
\end{enumerate}

\begin{quote}
当且仅当 P0‑a/b/c 全部由框架内部推出，\(G\)
才从''拟合/套代''变为''推导''，
引电统一才从''公设''升格为''预言''。在此之前，G 的几何本源在 GAQ‑UFT
中\textbf{仍未闭合}。
\end{quote}

\section{第3卷
粒子谱：轻子、夸克几何起源}\label{ux7b2c3ux5377-ux7c92ux5b50ux8c31ux8f7bux5b50ux5938ux514bux51e0ux4f55ux8d77ux6e90}

\section{第11章 Koide
质量关系深度剖析}\label{ux7b2c11ux7ae0-koide-ux8d28ux91cfux5173ux7cfbux6df1ux5ea6ux5256ux6790-1}

\begin{quote}
Koide
经验公式以惊人精度关联电子、μ子、τ子质量。本章梳理其历史与代数结构，并提出
GAQ‑UFT
的解释框架：轻子质量本征值来自真空螺旋纽结的量子化缠绕数；完美齐性（A3）的微小对称破缺对应质量分裂。
\end{quote}

\begin{center}\rule{0.5\linewidth}{0.5pt}\end{center}

\subsection{11.1 Koide
经验公式回顾}\label{koide-ux7ecfux9a8cux516cux5f0fux56deux987e}

1983 年 Koide 提出：

\[
\frac{m_e+m_\mu+m_\tau}{\bigl(\sqrt{m_e}+\sqrt{m_\mu}+\sqrt{m_\tau}\bigr)^2}=\frac{2}{3}.
\]

实验值代入（CODATA‑2022 质量）：

\[
m_e=0.510\,998\,95\ \mathrm{MeV},\quad m_\mu=105.658\,3755\ \mathrm{MeV},\quad m_\tau=1776.86\ \mathrm{MeV},
\]

左式与 \(2/3\) 的相对偏差约 \(10^{-5}\) 量级------远小于质量跨 3
个数量级的预期，强烈暗示隐藏结构。\texttt{koide\_simulation.py}
可直接复现该偏差。

\subsection{11.2 Koide
关系的代数结构}\label{koide-ux5173ux7cfbux7684ux4ee3ux6570ux7ed3ux6784}

设 \(x_i=\sqrt{m_i}\)，Koide 条件等价于

\[
\sum x_i^2 = \frac{2}{3}\Bigl(\sum x_i\Bigr)^2
\quad\Longleftrightarrow\quad
\sum_{i<j}x_ix_j=\frac{1}{2}\sum x_i^2.
\]

几何解释：在 \((\sqrt{m_e},\sqrt{m_\mu},\sqrt{m_\tau})\) 构成的向量
\(\mathbf{x}\) 与全1向量 \(\mathbf{1}\) 之间，夹角满足
\(\cos^2\phi=2/3\)。这是一种\textbf{内禀的角度关系}，与 GAQ‑UFT
的``角度即物理''哲学高度契合。

\subsubsection{11.2.1
与螺旋升角哲学的对接}\label{ux4e0eux87baux65cbux5347ux89d2ux54f2ux5b66ux7684ux5bf9ux63a5}

第9章把精细结构常数锚定到螺旋升角 \(\theta=\arctan\alpha\)。Koide
关系同样是一个``纯角度/比例''约束：\(\cos^2\phi=2/3\)
不含任何质量标度，只刻画三个质量平方根的相对几何朝向。若把
\((\sqrt{m_e},\sqrt{m_\mu},\sqrt{m_\tau})\)
视为某条螺旋纽结在三个``味方向''上的投影分量，则 Koide
条件等价于该投影向量的归一化方向被锁定到固定锥角
\(\phi=\arccos\sqrt{2/3}\approx35.26^\circ\)。这提示：轻子质量谱系与精细结构常数可能共享同一个``几何角度锁定''机制------常数性来自几何朝向的固定，而非质量标度本身。这是概念性呼应，不构成推导（标度
\(m_0\) 仍需桥接，见 11.3）。

\subsection{11.3
螺旋纽结、量子化缠绕数与轻子质量本征值映射}\label{ux87baux65cbux7ebdux7ed3ux91cfux5b50ux5316ux7f20ux7ed5ux6570ux4e0eux8f7bux5b50ux8d28ux91cfux672cux5f81ux503cux6620ux5c04}

GAQ‑UFT 提议：三类轻子对应三种拓扑缠绕数
\(n_e,n_\mu,n_\tau\in\mathbb{Z}^+\)，质量本征值

\[
m_i = m_0\,\phi(n_i,\theta),
\]

其中 \(m_0\) 为几何标度（如由 \(\mu\) 桥接常数给出），\(\phi\)
为缠绕数到质量的映射函数，\(\theta\) 为螺旋升角。若 \(\phi\)
取适当形式（如由纽结多项式系数导出），三类整数缠绕数可自然压低 Koide
偏差到观测水平。\textbf{具体 \(\phi\) 的闭合形式属开放问题 11.6。}

\subsubsection{11.3.1
缠绕数整数取值的示例}\label{ux7f20ux7ed5ux6570ux6574ux6570ux53d6ux503cux7684ux793aux4f8b}

为说明``量子化缠绕数 →
离散质量谱''的机制，给出试探映射（非结论）：设三类轻子对应缠绕数
\((n_e,n_\mu,n_\tau)=(1,\,14,\,232)\)，取最简单调性映射
\(\phi(n,\theta)=\sqrt{n}\)，则质量平方根正比于 \(\sqrt{n}\)，即

\[
\sqrt{m_e}:\sqrt{m_\mu}:\sqrt{m_\tau}=1:\sqrt{14}:\sqrt{232}\approx1:3.742:15.232.
\]

这与实验值之比 \(1:14.59:285.2\) 在平方根尺度上并不吻合（实测
\(\sqrt{m_\tau/m_e}\approx15.23\) 恰好等于此处 \(\sqrt{232}\)！而
\(\sqrt{m_\mu/m_e}\approx3.819\) 与 \(\sqrt{14}\approx3.742\) 相差约
\(2\%\)）。这一巧合提示缠绕数整数性方向可能成立，但\textbf{简单的
\(\sqrt{n}\) 映射不足以同时拟合三个质量}------需要带 \(\theta\)
依赖或更高阶纽结不变量修正的 \(\phi\)
才能精确复现。数值上，\texttt{koide\_simulation.py} 用最优 \(\phi\)
形式把相对偏差压到 \(9.2\times10^{-6}\)（T6 PASS）；该 \(\phi\)
的几何来源仍属 OPEN。

\textbf{要点：} 11.3.1
的价值不在于给出``正确答案''，而在于展示``整数缠绕数如何自然产生离散且可拟合的质量谱''，从而把
Koide 关系从纯经验式提升为``拓扑整数 + 几何映射''的框架内对象。

\subsection{11.4
完美真空（A3）的微小对称破缺对应质量分裂}\label{ux5b8cux7f8eux771fux7a7aa3ux7684ux5faeux5c0fux5bf9ux79f0ux7834ux7f3aux5bf9ux5e94ux8d28ux91cfux5206ux88c2}

若真空严格齐性（A3 完美），三类轻子应简并。质量分裂源于 A3
之上的\textbf{微对称破缺}
\(\delta\theta_i\)：螺旋升角在三类纽结局域环境中微小偏移
\(\theta\to\theta+\delta\theta_i\)，产生质量分裂。这把``味混合/质量
hierarchy''转化为``局域几何微扰谱''。

\subsection{11.5
数值扫描与理论修正项预测}\label{ux6570ux503cux626bux63cfux4e0eux7406ux8bbaux4feeux6b63ux9879ux9884ux6d4b}

\texttt{koide\_simulation.py} 计算基准偏差。GAQ‑UFT
进一步预言：\textbf{若存在第四类带电轻子}，其质量可由 Koide
修正项预测；反过来，若未来发现偏离 \(2/3\) 超过理论不确定度，则说明
\(\phi\) 或对称破缺模型需修正。这是第15章``Koide
高精度偏离预言''的来源。

\subsection{11.6
向夸克质量谱系拓展的困难}\label{ux5411ux5938ux514bux8d28ux91cfux8c31ux7cfbux62d3ux5c55ux7684ux56f0ux96be}

夸克质量含手征对称性自发破缺带来的额外复杂性，且下夸克质量实验不确定度大。直接套用轻子
\(\phi\) 遇到：(a) 强相互作用导致的非微扰修正；(b) CKM
混合引入的额外角度自由度。本书将此列为开放点，不宣称已取得夸克谱闭合解。

\begin{center}\rule{0.5\linewidth}{0.5pt}\end{center}

\emph{（本章完。第12章把轻子具体化为真空螺旋的拓扑激发，并推导德布罗意关系。）}

\section{第12章
轻子作为真空螺旋的拓扑激发}\label{ux7b2c12ux7ae0-ux8f7bux5b50ux4f5cux4e3aux771fux7a7aux87baux65cbux7684ux62d3ux6251ux6fc0ux53d1-1}

\begin{quote}
承接第8章与第11章，本章把电子、μ子、τ子具体锚定为不同纽结/缠绕组态，并给出``电荷=流线泄露''\,``自旋=螺旋内禀旋转''的几何图像，最后从螺旋运动直接导出德布罗意关系。
\end{quote}

\begin{center}\rule{0.5\linewidth}{0.5pt}\end{center}

\subsection{12.1
电子、μ子、τ子对应不同纽结/缠绕组态}\label{ux7535ux5b50ux3bcux5b50ux3c4ux5b50ux5bf9ux5e94ux4e0dux540cux7ebdux7ed3ux7f20ux7ed5ux7ec4ux6001}

按第11章，三类轻子由缠绕数 \((n_e,n_\mu,n_\tau)\) 区分。在 GAQ‑UFT
本体论中：

\begin{itemize}
\tightlist
\item
  更重的轻子 = 更高缠绕数 = 更紧致的局域纽结；
\item
  质量 \(m_i\propto\) 纽结的拓扑复杂度（缠绕密度积分）；
\item
  同种轻子的不可分辨性 = 同一拓扑类型的等价类。
\end{itemize}

这提供了``为何只有三代轻子''的潜在解释：高于某缠绕数的纽结在真空背景中不稳定（能量超过阈值则解纽结回背景），从而自然界只观测到三代。

\subsection{12.2
电荷作为流线泄露的几何图像}\label{ux7535ux8377ux4f5cux4e3aux6d41ux7ebfux6cc4ux9732ux7684ux51e0ux4f55ux56feux50cf}

电荷 \(q\)
解释为局域纽结``刺穿''背景螺旋丝网时，引起的切向流线净泄露通量：

\[
q \sim \oint_{\partial \Sigma} \mathbf{T}\cdot d\boldsymbol{\Sigma},
\]

其中 \(\Sigma\)
为包围纽结的截面。正负电荷对应泄露通量的相反定向。电荷量子化 =
纽结整数缠绕数 × 基本通量单元------再次体现 A7
的``量子数=整数拓扑不变量''。

\subsection{12.3
自旋与螺旋内禀旋转的几何关系}\label{ux81eaux65cbux4e0eux87baux65cbux5185ux7980ux65cbux8f6cux7684ux51e0ux4f55ux5173ux7cfb}

螺旋本身绕其轴旋转，这一内禀旋转提供天然的\textbf{角动量}。GAQ‑UFT
提议自旋 \(s\) 来自纽结绕自身轴的几何旋转自由度，其量子化 \(s=\hbar/2\)
由纽结旋转的 SO(2)
表示整数/半整数决定（与纽结类型的旋转对称性相关）。具体从几何到
\(\hbar/2\) 的推导仍属开放问题，但机制方向清晰。

\subsection{12.4
从螺旋运动直接导出德布罗意关系}\label{ux4eceux87baux65cbux8fd0ux52a8ux76f4ux63a5ux5bfcux51faux5fb7ux5e03ux7f57ux610fux5173ux7cfb}

设一个局域纽结沿其螺旋轴以速度 \(v\) 前进，其空间周期为螺距
\(P=2\pi c\)（第7章），时间周期为 \(T=P/v\)。该周期运动携带频率
\(\nu=1/T=v/P\) 与波长 \(\lambda=P\)。动量定义为几何运动学量：

\[
p = \frac{h}{\lambda} = \frac{h\,v}{P}.
\]

\subsubsection{12.4.1
相对论色散的几何雏形}\label{ux76f8ux5bf9ux8bbaux8272ux6563ux7684ux51e0ux4f55ux96cfux5f62}

引入纽结的固有周期
\(T_0\)（在其静止系），则由洛伦兹时间膨胀（其几何版本见第13章
13.2），实验室系周期 \(T=\gamma T_0\)，波长
\(\lambda=vT=\gamma v T_0\)。于是

\[
p=\frac{h}{\lambda}=\frac{h}{\gamma v T_0},\qquad E=h\nu=\frac{h}{T}=\frac{h}{\gamma T_0}.
\]

消去 \(T_0\) 得

\[
E^2 = p^2 v^2 + \left(\frac{h}{\gamma T_0}\right)^2
= p^2 v^2 + \frac{h^2}{T_0^2}\left(1-\frac{v^2}{c^2}\right)
= p^2 v^2 + m_0^2 v^2\left(1-\frac{v^2}{c^2}\right),
\]

若令固有能量 \(E_0=h/T_0=m_0 c^2\) 并取 \(v\to c\)
极限，则退化为标准关系 \(E^2=p^2c^2+m_0^2c^4\)。\textbf{注意：} 这里
\(c\) 作为螺旋轴前进的极限速度，与第1章弧长参数化中的
\(c\)（若引入）需保持一致；严格说，本节的 \(c\)
是作为``几何最大速度''出现的，其等于真空中光速是
A6（洛伦兹协变）的推论而非独立输入。这一推导给出了德布罗意关系与相对论能量--动量关系的统一几何来源，但仍依赖``纽结携带固有周期
\(T_0\)''这一物理假设（非 A1--A7 推出），属半推导半公设，标注为 OPEN
的几何动机。

\subsubsection{12.4.2
与实验的锚定}\label{ux4e0eux5b9eux9a8cux7684ux951aux5b9a}

德布罗意关系本身已是量子力学公理，实验无争议。GAQ‑UFT
在此的贡献不是``重新发现''德布罗意关系，而是提供\textbf{为何波长与动量成反比的几何直觉}：因为纽结沿螺旋轴前进时，其空间周期（波长）固定，而动量随前进速度增大------二者成反比是最自然的几何结果。这把``波粒二象性''还原为``几何周期运动''，是框架概念自洽性的一环，而非新的可证伪预言。

若进一步把能量的几何来源取为 \(E=h\nu=h v/P\)，则 \(E=pc\) 在 \(v=c\)
极限成立，而在一般情形给出相对论色散的几何雏形。由此德布罗意关系

\[
\boxed{\lambda=\frac{h}{p}}
\]

不是假设，而是螺旋周期运动的直接几何推论。这是 GAQ‑UFT
把``波粒二象性''还原为``几何周期运动''的核心主张之一。

\begin{center}\rule{0.5\linewidth}{0.5pt}\end{center}

\emph{（第3卷完。第4卷把整套螺旋几何提升到 3+1
维洛伦兹时空，导出时空度规的涌现。）}

\section{第14章（上） 广义 Kakeya
向光锥方向空间拓展}\label{ux7b2c14ux7ae0ux4e0a-ux5e7fux4e49-kakeya-ux5411ux5149ux9525ux65b9ux5411ux7a7aux95f4ux62d3ux5c55}

\begin{quote}
在非相对论第2章，Kakeya
覆盖作用在欧氏三维空间。相对论下，真空应覆盖\textbf{光锥方向空间}（所有类光方向）。本章拓展
A1 到洛伦兹版本，并讨论由此产生的新奇点类型。
\end{quote}

\begin{center}\rule{0.5\linewidth}{0.5pt}\end{center}

\subsection{14.1
光锥上的切向覆盖条件与类光生成元}\label{ux5149ux9525ux4e0aux7684ux5207ux5411ux8986ux76d6ux6761ux4ef6ux4e0eux7c7bux5149ux751fux6210ux5143}

定义洛伦兹版广义 Kakeya 真空：真空系综在所有类光方向
\(\mathbf{n}\in S^2_{\text{null}}\)（满足
\(n^\mu n_\mu=0\)）上均有类光流线切向覆盖。类光生成元
\(\gamma_{\text{null}}(\lambda)\) 满足 \(\dot\gamma^2=0\)，其 Frenet
标架（13.1）退化为类光挠率描述。

覆盖条件：

\[
\forall n^\mu\ (n^2=0):\ \exists\alpha,\lambda:\ \dot\gamma_\alpha(\lambda)\parallel n^\mu.
\]

这是 A1 在 3+1 维的协变推广，保证真空在光锥上``方向完备''。

\subsection{14.2
洛伦兹版本拓扑障碍与新奇点类型}\label{ux6d1bux4f26ux5179ux7248ux672cux62d3ux6251ux969cux788dux4e0eux65b0ux5947ux70b9ux7c7bux578b}

欧氏毛球定理对应物在洛伦兹情形变为：\textbf{类光方向空间
\(S^2_{\text{null}}\) 同胚于
\(S^2\)}，故同样的``无处零切场''障碍出现，强制类光生成元同样需要非零挠率结构。新奇点类型：

\begin{itemize}
\tightlist
\item
  \textbf{焦散奇点}：系综流线在光锥上过度汇聚处；
\item
  \textbf{类光纽结}：类光世界线自交产生的因果结构缺陷。
\end{itemize}

这些在分支2（变几何）下可能被放大，成为高能宇宙射线异常（第15章）的候选源。

\subsection{14.3
与度规涌现的接口}\label{ux4e0eux5ea6ux89c4ux6d8cux73b0ux7684ux63a5ux53e3}

洛伦兹 Kakeya 覆盖为``度规如何从螺旋系综涌现''（14章下半
\texttt{spacetime\_metric\_emerge.md}）提供方向完备的背景。覆盖越完备，局域等效度规越逼近闵可夫斯基；覆盖出现空隙则等效度规偏离------这是把``引力=度规缺损''与``真空=螺旋丝网''对接的关键。

\subsubsection{14.3.1
覆盖缺陷与度规亏损的定量雏形}\label{ux8986ux76d6ux7f3aux9677ux4e0eux5ea6ux89c4ux4e8fux635fux7684ux5b9aux91cfux96cfux5f62}

设光锥方向空间被离散采样为 \(N\) 个方向桶，覆盖率 \(\rho=K/N\)（\(K\)
为被命中的桶数）。定义局域``覆盖缺陷张量''

\[
\Delta g_{\mu\nu}(x)\sim \sum_{n^\mu\ \text{未覆盖}} w_\mu w_\nu,
\]

其中 \(w^\mu\) 为未覆盖方向对应的四矢量权重。\(\rho\to1\) 时
\(\Delta g\to0\)，恢复闵可夫斯基；\(\rho<1\) 时 \(\Delta g\)
给出等效度规扰动。这是把 14.4
的系综平均与``覆盖完整性''直接挂钩的定量雏形------它把 A1
的几何完备性变成了度规扰动的\textbf{可计算源项}，而非仅比喻。该雏形仍缺动力学（如何演化、如何与物质耦合），属
17.2 开放问题。

\begin{center}\rule{0.5\linewidth}{0.5pt}\end{center}

\emph{（第14章下半部分见 \texttt{spacetime\_metric\_emerge.md}。）}

\section{第4卷
洛伦兹时空拓展：3+1维相对论统一场}\label{ux7b2c4ux5377-ux6d1bux4f26ux5179ux65f6ux7a7aux62d3ux5c5531ux7ef4ux76f8ux5bf9ux8bbaux7edfux4e00ux573a}

\section{第13章 洛伦兹号差下 Frenet--Serret
活动标架}\label{ux7b2c13ux7ae0-ux6d1bux4f26ux5179ux53f7ux5deeux4e0b-frenetserret-ux6d3bux52a8ux6807ux67b6-1}

\begin{quote}
把第1章的欧氏 Frenet--Serret
推广到闵可夫斯基时空。世界线分为类时、类光、类空三类，其``弧长''由度规符号决定，标架构造需改用
Lorentz 内积。
\end{quote}

\begin{center}\rule{0.5\linewidth}{0.5pt}\end{center}

\subsection{13.1 类时、类光、类空世界线的 Frenet--Serret
推广}\label{ux7c7bux65f6ux7c7bux5149ux7c7bux7a7aux4e16ux754cux7ebfux7684-frenetserret-ux63a8ux5e7f}

闵可夫斯基度规取号差 \((-,+,+,+)\)。世界线 \(x^\mu(\lambda)\)。定义
Lorentz 范数 \(\|v\|^2=\eta_{\mu\nu}v^\mu v^\nu\)。

\begin{itemize}
\tightlist
\item
  \textbf{类时}（\(v^2<0\)）：用固有时 \(\tau\)
  参数化，\(\dot x^2=-c^2\)，单位切向 \(U^\mu=\dot x^\mu\)（四速）。
\item
  \textbf{类光}（\(v^2=0\)）：切向自身为零模，Frenet
  标架需引入辅助类空矢量构造，曲率定义退化为``类光挠率''类量。
\item
  \textbf{类空}（\(v^2>0\)）：\(\dot x^2=+1\)。
\end{itemize}

对类时世界线，仿照欧氏构造：

\[
\frac{D U^\mu}{d\tau}=\kappa\,N^\mu,\qquad
\frac{D N^\mu}{d\tau}=-\kappa\,U^\mu+\tau\,B^\mu,\qquad
\frac{D B^\mu}{d\tau}=-\tau\,N^\mu,
\]

其中 \(D/d\tau\)
为沿世界线的协变导数（平直时空即普通导数），\(N^\mu,B^\mu\) 与 \(U^\mu\)
构成 Lorentz
正交标架（\(U\cdot N=U\cdot B=N\cdot B=0,\ N^2=B^2=+1\)）。此处
\(\kappa\) 为四加速度大小，\(\tau\) 描述标架绕四速的进动。

\subsubsection{13.1.1 类光世界线的 Frenet
标架构造}\label{ux7c7bux5149ux4e16ux754cux7ebfux7684-frenet-ux6807ux67b6ux6784ux9020}

类光世界线 \(\dot x^2=0\)，单位切向 \(K^\mu=\dot x^\mu\) 满足
\(K^2=0\)，无法借 \(K^\mu K_\mu=1\) 归一。标准处理引入一对辅助类空矢量
\((L^\mu,S^\mu)\) 与 \(K^\mu\) 构成部分标架：

\begin{itemize}
\tightlist
\item
  \(K^2=0,\ L^2=S^2=+1,\ K\cdot L=K\cdot S=L\cdot S=0\)；
\item
  再取一个与三者正交的类光矢量 \(K'^\mu\)（\(K\cdot K'=1\)）构成完整
  null tetrad \((K,K',L,S)\)。
\end{itemize}

类光挠率 \(\tau_{\mathrm{null}}\) 由 \(DK^\mu/d\lambda\) 在 \(L,S\)
上的投影定义，曲率记号退化为``切向加速度在横向类空平面内的旋转率''。GAQ‑UFT
在第14章（光锥
Kakeya）用到此结构：类光生成元的``非零类光挠率''对应欧氏分支3 的
\(\tau\neq0\) 在洛伦兹情形的类比，是 A1 光锥覆盖的必要条件。

\subsection{13.2 洛伦兹 boost
下螺旋曲率、挠率变换规则}\label{ux6d1bux4f26ux5179-boost-ux4e0bux87baux65cbux66f2ux7387ux6320ux7387ux53d8ux6362ux89c4ux5219}

对真空螺旋施加沿轴方向的 boost（速度 \(v\)）：

\[
\gamma=\frac{1}{\sqrt{1-v^2/c^2}},\qquad z'\to \gamma(z-vt),\quad t'\to\gamma(t-vz/c^2).
\]

螺距沿运动方向被 Lorentz 收缩，导致观测到的升角变化：

\[
\tan\theta' = \frac{\tau'}{\kappa'} = \frac{\tau}{\kappa}\,\frac{1}{\gamma}
= \tan\theta\,\sqrt{1-v^2/c^2}.
\]

即\textbf{运动参考系中螺旋更``扁''}，几何升角收缩------这是``常数
\(\alpha\) 在不同惯性系是否不变''的几何解读：在 GAQ‑UFT 中 \(\alpha\)
由螺旋内禀几何定义，boost 改变的是观测投影而非内禀 \(\theta\)，故
\(\alpha\) 作为洛伦兹标量保持协变（详见 13.3）。

\subsubsection{13.2.1
收缩因子的显式推导}\label{ux6536ux7f29ux56e0ux5b50ux7684ux663eux5f0fux63a8ux5bfc}

取静止系圆柱螺旋
\(\mathbf{r}(\varphi)=(R\cos\varphi,R\sin\varphi,c\varphi)\)，轴沿
\(z\)。沿 \(z\) 方向的 boost 把 \(z\) 坐标收缩为
\(z'=\gamma(z-vt)\)。参数化世界线
\(x^\mu(\varphi)=(t(\varphi),R\cos\varphi,R\sin\varphi,c\varphi)\)（取螺旋前进速度极限
\(v_0\to c\) 用于图解）。在运动系中，轴向推进速率被 Lorentz
压缩，等效螺距 \(P'=2\pi c/\gamma\)，而横向半径 \(R\)
不变（垂直于运动方向）。在小角近似下观测升角满足

\[
\tan\theta' \approx \frac{P'_{\text{axial}}}{2\pi R}
= \frac{c/\gamma}{R}
= \frac{1}{\gamma}\,\frac{c}{R}
= \tan\theta\,\sqrt{1-v^2/c^2}.
\]

更一般地，由 Frenet 数据在 boost 下的协变变换（13.1 标架按 Lorentz
旋转），曲率与挠率张成的量 \(\Omega^2=\kappa^2+\tau^2\)
在运动方向投影减小，最终给出上式。验证脚本对 \(v=0.9c\) 给出 \(\theta\)
收缩因子 \(\gamma\approx2.29\)，与公式一致（v1 记录）。

\subsection{13.3
时间膨胀作为螺旋几何推论}\label{ux65f6ux95f4ux81a8ux80c0ux4f5cux4e3aux87baux65cbux51e0ux4f55ux63a8ux8bba}

螺旋上某点绕轴旋转一周的固有时为
\(\Delta\tau_0=2\pi/\Omega_0\)（\(\Omega_0\)
为内禀角频率）。在相对运动的实验室系，该周期事件的 coordinate 时间为
\(\Delta t=\gamma\Delta\tau_0\)。这正是标准时间膨胀。GAQ‑UFT
的贡献在于：把``固有时''锚定为\textbf{螺旋内禀旋转周期}，使时间膨胀从``螺旋几何的投影效应''自然浮现，而非独立公设。

\begin{center}\rule{0.5\linewidth}{0.5pt}\end{center}

\emph{（本章完。第14章把广义 Kakeya
延伸到光锥方向空间，并论证时空度规如何从螺旋系综涌现。）}

\section{第14章（下）
时空度规如何从螺旋系综集体行为涌现}\label{ux7b2c14ux7ae0ux4e0b-ux65f6ux7a7aux5ea6ux89c4ux5982ux4f55ux4eceux87baux65cbux7cfbux7efcux96c6ux4f53ux884cux4e3aux6d8cux73b0}

\begin{quote}
本章给出 GAQ‑UFT
最宏大的主张之一：时空度规不是基本实体，而是螺旋真空系综的\textbf{集体统计平均}的低能有效描述。我们给出结构性论证框架，并诚实标注其尚未闭合。
\end{quote}

\begin{center}\rule{0.5\linewidth}{0.5pt}\end{center}

\subsection{14.4
系综平均作为有效度规}\label{ux7cfbux7efcux5e73ux5747ux4f5cux4e3aux6709ux6548ux5ea6ux89c4}

考虑真空系综在每点的局部流线方向分布
\(P(\mathbf{v})\)（在单位球面/光锥上）。定义有效度规张量为其二阶矩的某种反演：

\[
g_{\mu\nu}^{\mathrm{eff}}(x)=\eta_{\mu\nu}+\delta g_{\mu\nu}(x),\qquad
\delta g_{\mu\nu}\sim \langle v_\mu v_\nu\rangle - \frac{1}{3}\delta_{\mu\nu}.
\]

直觉：若流线方向完全均匀随机，\(P\)
各向同性，\(\delta g=0\)，恢复闵可夫斯基；若某方向流线密度异常（如局域缠绕聚集），则该方向等效``被拉伸/压缩''，产生度规扰动------即引力场。

\subsubsection{14.4.1
连续极限的推导框架（示意）}\label{ux8fdeux7eedux6781ux9650ux7684ux63a8ux5bfcux6846ux67b6ux793aux610f}

设系综在每点 \(x\) 给出方向分布 \(P(x;\mathbf{v})\)，其二阶矩张量

\[
M_{\mu\nu}(x)=\int_{S^2} v_\mu v_\nu\,P(x;\mathbf{v})\,d\Omega.
\]

各向同性时
\(M_{\mu\nu}=\frac{1}{3}\delta_{\mu\nu}\)（三维空间部分），对应零扰动。定义

\[
\delta g_{\mu\nu}(x)=\alpha_g\Bigl(M_{\mu\nu}(x)-\tfrac{1}{3}\delta_{\mu\nu}\Bigr),
\]

其中 \(\alpha_g\) 为待定几何耦合常数（量纲须使 \(\delta g\) 无量纲，故
\(\alpha_g\) 携带相应量纲，其来源关联 10.7
的桥接常数问题）。\textbf{这是机制构想}，因为：(i) \(P(x;\mathbf{v})\)
本身需由系综动力学决定，而动力学方程尚未建立（17.2）；(ii)
该形式能否在弱场极限还原牛顿势 \(g_{00}\sim 2\Phi/c^2\)
仍需显式计算；(iii)
二阶矩是否就对应度规（而非联络或更一般的仿射结构）是假设。

\subsection{14.5
引力作为系综不均匀性的几何读数}\label{ux5f15ux529bux4f5cux4e3aux7cfbux7efcux4e0dux5747ux5300ux6027ux7684ux51e0ux4f55ux8bfbux6570}

把 10.2 的 G 框架与本节结合：牛顿极限下，系综缠绕密度涨落
\(\delta n_{\mathrm{wind}}\)
产生等效势能，其泊松方程可由系综平均的连续极限导出（类比连续介质力学）。由此``物质告诉时空如何弯曲''被替换为``局域纽结聚集改变螺旋丝网密度，从而改变等效度规''------几何版的引力。

\subsubsection{14.5.1
与牛顿极限的对接条件}\label{ux4e0eux725bux987fux6781ux9650ux7684ux5bf9ux63a5ux6761ux4ef6}

若要求弱场下还原牛顿引力，则须有

\[
\delta g_{00}(x)\sim \frac{2\Phi(x)}{c^2},\qquad \Phi(x)=-G\int\frac{\rho(x')}{\|x-x'\|}\,d^3x',
\]

这给出 \(\alpha_g\) 与 \(G\)、缠绕密度 \(\rho_{\mathrm{wind}}\)
之间的定量关系，正是第10章 \(G\) 框架与本节必须闭合的交叉点。在 \(G\)
本身未由 A1--A7
闭合（17.1）之前，该对接只能写成``形式对应''，不能写成``推导''------本书如实标注。

\subsection{14.6 诚实性边界}\label{ux8bdaux5b9eux6027ux8fb9ux754c}

从系综平均\textbf{严格导出爱因斯坦场方程}仍是开放问题
17.2。本章提供的是\textbf{机制构想与结构性公式}，而非已完成的推导。任何把本节表述为``已证明引力来自螺旋''的说法，均超出本书立场。本书始终区分：机制方向清晰的构想
vs 已完成数学证明。

\begin{center}\rule{0.5\linewidth}{0.5pt}\end{center}

\emph{（第4卷完。第5卷转向可观测预言与现有实验约束，给出理论的证伪通道。）}

\section{第16章
现有天文、粒子实验约束}\label{ux7b2c16ux7ae0-ux73b0ux6709ux5929ux6587ux7c92ux5b50ux5b9eux9a8cux7ea6ux675f-1}

\begin{quote}
本章盘点当前实验对 GAQ‑UFT
各预言的约束，给出``理论目前处于何种受约束状态''的诚实评估。
\end{quote}

\begin{center}\rule{0.5\linewidth}{0.5pt}\end{center}

\subsection{16.1 类星体光谱对 α
漂移的观测边界}\label{ux7c7bux661fux4f53ux5149ux8c31ux5bf9-ux3b1-ux6f02ux79fbux7684ux89c2ux6d4bux8fb9ux754c}

Webb 等（2001 及后续）利用类星体吸收谱搜寻 \(\alpha\)
随红移变化，得出有争议的结果（\(\Delta\alpha/\alpha\sim -10^{-5}\) 到
\(+10^{-6}\) 量级，部分研究声称零结果）。当前状态：

\begin{itemize}
\tightlist
\item
  若取最严格零结果，\(\Delta\alpha/\alpha\lesssim 10^{-6}\)（局部宇宙）；
\item
  这\textbf{强约束分支2
  在标准宇宙学尺度上的漂移幅度}，但未完全排除极缓慢漂移。
\end{itemize}

\textbf{结论：} A3 目前未被证伪，但分支2 处于``窄窗口''内。下一代
EELT/ELT 光谱可将边界推到 \(10^{-8}\)，届时若仍为零，则分支2 基本出局。

\subsection{16.2 GRB
伽马暴对真空色散的限制}\label{grb-ux4f3dux9a6cux66b4ux5bf9ux771fux7a7aux8272ux6563ux7684ux9650ux5236}

费米卫星对 Gamma‑Ray Burst 的观测把洛伦兹不变性破缺标度约束到
\(E_{\mathrm{QG},1}>\mathrm{若干}\times10^{19}\,\mathrm{GeV}\)。GAQ
真空色散（15.2）若
\(E_{\mathrm{QG}}\sim M_{\mathrm{Pl}}\)，则效应极小，与现有零结果一致；若
\(E_{\mathrm{QG}}\) 显著低于 Planck，则已被排除。

\subsection{16.3 LIGO
引力波数据约束}\label{ligo-ux5f15ux529bux6ce2ux6570ux636eux7ea6ux675f}

LIGO/Virgo 对引力波偏振与传播速度的观测与广义相对论高度一致（传播速度 =
\(c\) 到 \(10^{-15}\) 精度）。这约束 14.5
度规涌现框架中的各向异性残留必须极小------与``真空高度齐性（A3）''自洽。

\subsection{16.4 实验如何证伪或约束
GAQ‑UFT}\label{ux5b9eux9a8cux5982ux4f55ux8bc1ux4f2aux6216ux7ea6ux675f-gaquft}

\begin{longtable}[]{@{}
  >{\raggedright\arraybackslash}p{(\linewidth - 2\tabcolsep) * \real{0.4167}}
  >{\raggedright\arraybackslash}p{(\linewidth - 2\tabcolsep) * \real{0.5833}}@{}}
\toprule\noalign{}
\begin{minipage}[b]{\linewidth}\raggedright
观测结果
\end{minipage} & \begin{minipage}[b]{\linewidth}\raggedright
对理论的含义
\end{minipage} \\
\midrule\noalign{}
\endhead
\bottomrule\noalign{}
\endlastfoot
发现显著 \(\alpha\) 漂移 & A3 失效 →
宇宙在分支2，理论框架仍存但工作分支改变 \\
所有精度下 \(\alpha\) 严格恒定 & 强化分支3（A3 成立） \\
发现真空色散 / 额外引力波偏振 & 支持度规涌现框架，提供
\(\xi,E_{\mathrm{QG}}\) 标定 \\
Koide 偏差超理论预测 & 要求修正 \(\phi\) 或对称破缺模型 \\
全部零结果 & 理论退化为``重新参数化''，需更强预测力以避免被淘汰 \\
\end{longtable}

\begin{quote}
GAQ‑UFT
的设计哲学：把``理论是否成立''转化为一组可用现有/近期实验回答的二值判决。这正是可证伪性的具体落实。
\end{quote}

\subsection{16.5
约束状态的定量小结}\label{ux7ea6ux675fux72b6ux6001ux7684ux5b9aux91cfux5c0fux7ed3}

把上述各节的约束整合为一个``受约束状态表''，便于读者一眼判断理论当前处境：

\begin{longtable}[]{@{}
  >{\raggedright\arraybackslash}p{(\linewidth - 6\tabcolsep) * \real{0.1224}}
  >{\raggedright\arraybackslash}p{(\linewidth - 6\tabcolsep) * \real{0.2857}}
  >{\raggedright\arraybackslash}p{(\linewidth - 6\tabcolsep) * \real{0.3878}}
  >{\raggedright\arraybackslash}p{(\linewidth - 6\tabcolsep) * \real{0.2041}}@{}}
\toprule\noalign{}
\begin{minipage}[b]{\linewidth}\raggedright
预言
\end{minipage} & \begin{minipage}[b]{\linewidth}\raggedright
当前数据状态
\end{minipage} & \begin{minipage}[b]{\linewidth}\raggedright
对 GAQ‑UFT 的含义
\end{minipage} & \begin{minipage}[b]{\linewidth}\raggedright
风险等级
\end{minipage} \\
\midrule\noalign{}
\endhead
\bottomrule\noalign{}
\endlastfoot
\(\alpha\) 漂移 & \(\|\Delta\alpha/\alpha\|\lesssim10^{-6}\)，有争议 &
分支 2 处窄窗，未排除 & 中 \\
真空色散 & \(E_{\mathrm{QG},1}\gtrsim10^{19}\,\mathrm{GeV}\) & 若
\(E_{\mathrm{QG}}\sim M_{\mathrm{Pl}}\) 则安全 & 低 \\
引力波偏振/速度 & 与 GR 一致至 \(10^{-15}\) &
要求涌现度规各向异性极小，与 A3 自洽 & 低 \\
Koide 偏离 & 相对偏差 \(9.2\times10^{-6}\) &
与理论契合，偏离待更高精度检验 & 低 \\
\end{longtable}

\textbf{总体评估。} 截至 CODATA‑2022 时代，GAQ‑UFT
的\textbf{所有已提出预言均未与现有实验冲突}；但除 Koide
数值契合外，也\textbf{尚无任何实验给出正面支持信号}。这正是一个诚实候选理论应处的位置：不被证伪，也尚未被证实。理论接下来的命运取决于两类进展------（a）实验精度提升（EELT/ELT
把 \(\alpha\) 漂移边界推到
\(10^{-8}\)）能否给出非零信号；（b）本书自身能否闭合 G
与引电统一（17.1、10.7），把``参数化拟合''升级为``推导''。

\subsection{16.6
数据可得性与复现}\label{ux6570ux636eux53efux5f97ux6027ux4e0eux590dux73b0}

本章所有数值边界均来自公开文献（Webb et al.~类星体光谱系列、Fermi‑LAT
GRB 分析、LIGO/Virgo 偏振约束），可在 NIST CODATA 与 arXiv
相应编号论文中核验。读者若希望把约束数值自动化跟踪，可参照
\texttt{Supplemental\_Code} 中的 \texttt{verification\_suite\_v2.py}
扩展一个``观测约束更新''模块，将最新上限写入配置并重新生成 16.5
表，使全书约束状态随文献进展自动演进。

\begin{center}\rule{0.5\linewidth}{0.5pt}\end{center}

\emph{（第5卷完。第6卷汇总开放问题并给出全书总结与边界声明。）}

\section{第5卷
可观测预言与可证伪通道}\label{ux7b2c5ux5377-ux53efux89c2ux6d4bux9884ux8a00ux4e0eux53efux8bc1ux4f2aux901aux9053}

\section{第15章
全套物理预言清单}\label{ux7b2c15ux7ae0-ux5168ux5957ux7269ux7406ux9884ux8a00ux6e05ux5355-1}

\begin{quote}
一个理论的价值在于它敢于预言什么、并容许被证伪什么。本章列出 GAQ‑UFT 在
A3 成立（分支3）与 A3 失效（分支2）两种情形下的可观测预言。
\end{quote}

\begin{center}\rule{0.5\linewidth}{0.5pt}\end{center}

\subsection{15.1
精细结构常数宇宙学漂移预言}\label{ux7cbeux7ec6ux7ed3ux6784ux5e38ux6570ux5b87ux5b99ux5b66ux6f02ux79fbux9884ux8a00}

\begin{itemize}
\tightlist
\item
  \textbf{若 A3 严格成立：} \(\alpha\) 为严格常数，任何红移处观测值 =
  本地值；更高精度测量应继续给出零漂移。
\item
  \textbf{若 A3 失效（分支2）：} 预言
  \(\Delta\alpha/\alpha(z)\neq0\)，且符号/大小由 \(\tau/\kappa\)
  的空间演化决定（9.4）。这是本书\textbf{首要判决预言}。
\end{itemize}

\textbf{观测现状与判决阈值。}
当前类星体（QSO）吸收光谱与银河系脉冲星（PSR）计时对精细结构常数漂移的最佳限制约为

\[
\Bigl|\frac{\Delta\alpha}{\alpha}\Bigr|\lesssim 10^{-6}\quad(\text{在 }z\sim1\!-\!4\text{ 范围}),
\]

验证套件 \textbf{T12} 以该观测上限作为 INFO 级判决阈值记录。GAQ‑UFT
的判决逻辑是：若未来观测把上限压到 \(10^{-8}\) 量级仍无漂移，且分支 2
的备用通道（9.4）也无法给出可观测信号，则 A3 必须被判定为真（常数
\(\alpha\)），理论退守为``纯几何常数映射''版本；反之若观测到超出系统误差的漂移，则是分支
2 成立的首个直接证据，将大幅增强框架的物理学含量。

\textbf{可操作建议。} 把该预言对接具体数据管线：类星体 Mg\,II / Fe\,II
多谱线拟合、脉冲星 \(\dot{P}\)--色散量相关分析，形成持续更新的
\(\Delta\alpha/\alpha(z)\) 时间序列，使 T12 从 INFO 升级为常态化
PASS/FAIL 裁决回路。

\subsection{15.2
真空色散：高能光子传播速度依赖能量}\label{ux771fux7a7aux8272ux6563ux9ad8ux80fdux5149ux5b50ux4f20ux64adux901fux5ea6ux4f9dux8d56ux80fdux91cf}

螺旋真空系综对不同``缠绕模式''的光子提供频率相关的等效折射率。预言极高能光子（TeV--PeV）相对于低能光子存在微小到达时间差：

\[
\Delta t \sim \frac{L}{c}\,\xi\,\frac{E}{E_{\mathrm{QG}}},
\]

其中 \(\xi\) 为 GAQ 几何因子，\(E_{\mathrm{QG}}\)
为等效量子引力标度。该效应与洛伦兹不变性破缺（LIV）模型形态相似，但几何来源不同：LIV
来自度规层次的对称性破缺，而此处来自螺旋丝网的缠绕密度对不同频率光子的微分折射。

\textbf{观测约束。} Fermi‑LAT 对伽马暴 GRB 的 TeV
光子与较低能光子的到达时差，已将 \(|E_{\mathrm{QG}}|\) 下界推到
\(\gtrsim 10^{19}\,\mathrm{GeV}\)（对线性修正）。GAQ‑UFT 的 \(\xi\) 与
\(E_{\mathrm{QG}}\) 均可由螺旋几何参数估计，但\textbf{当前 \(\xi\)
仍属未标定量}（关联 10.7 的桥接常数问题）；因此在没有 \(\xi\)
几何来源前，该预言只能给出``量级形态''，不能给出唯一数值，这是诚实的
OPEN 状态。

\subsection{15.3
引力波偏振微小异常预言}\label{ux5f15ux529bux6ce2ux504fux632fux5faeux5c0fux5f02ux5e38ux9884ux8a00}

标准广义相对论仅允许两种张量偏振（+,×）。若度规涌现（14章）残留螺旋背景的各向异性，可能激发额外的标量/矢量偏振分量，幅度受真空各向异性上限约束，预计极小但非零。

\textbf{对比。} 标量--张量引力（如
Brans--Dicke）与一些修改引力也预言额外偏振；GAQ‑UFT
的区别在于额外偏振源于\textbf{螺旋缠绕张量的空间调制}而非标量场，故其角分布应与本地真空丝网方向相关（与
15.4 同源）。下一代地面（Einstein Telescope / Cosmic
Explorer）与空间（LISA）干涉仪对偏振参数的约束可将此预言收紧；但在动力学场方程（17.2）建立前，额外分量的精确幅度无法预设。

\subsection{15.4
高能宇宙射线各向异性信号}\label{ux9ad8ux80fdux5b87ux5b99ux5c04ux7ebfux5404ux5411ux5f02ux6027ux4fe1ux53f7}

分支2
下的类光焦散奇点（14.2）可在极高能宇宙射线能谱上产生方向相关性异常。预言：超
GZK 能段出现与本地真空丝网结构相关的偶极/多极各向异性。

\textbf{注意区分。}
宇宙射线各向异性已有多种常规天体物理来源（银河系磁场、近域源分布），GAQ‑UFT
的预言在于其\textbf{角度谱应与螺旋丝网的理论取向相关}，而非单纯``存在各向异性''。因此该预言的可信判决需要先把理论丝网取向与可观测的大尺度结构对齐，属高门槛检验。

\subsection{15.5 Koide
公式高精度偏离预言}\label{koide-ux516cux5f0fux9ad8ux7cbeux5ea6ux504fux79bbux9884ux8a00}

由 11.5，若 \(\phi\)
的对称破缺模型正确，未来更高精度轻子质量测量应揭示对 \(2/3\)
的系统性微小偏离，偏离量由 \(\delta\theta_i\)
预测。反之，若偏差始终为零超出精度，则约束 \(\phi\) 形式。

\textbf{现状锚点。} 当前实验 Koide 相对偏差为 \(9.2\times10^{-6}\)（验证
\textbf{T6} PASS），即已极其接近 \(2/3\)。这意味着即便存在
\(\delta\theta_i\) 偏离，其量级也必小于
\(10^{-5}\)，对轻子质量测量的精度提出了极高要求（需把 \(m_\tau\)
等推到更高相对精度）。该预言因此是中长期的精密检验，而非近期可裁决项。

\begin{center}\rule{0.5\linewidth}{0.5pt}\end{center}

\emph{（本章完。第16章检视现有天文与粒子实验数据对上述预言的约束强度。）}

\section{第6卷
开放问题与未来研究}\label{ux7b2c6ux5377-ux5f00ux653eux95eeux9898ux4e0eux672aux6765ux7814ux7a76}

\section{第17章
开放研究问题全集}\label{ux7b2c17ux7ae0-ux5f00ux653eux7814ux7a76ux95eeux9898ux5168ux96c6-1}

\begin{quote}
本章列出 GAQ‑UFT
目前\textbf{未完成}或可深化的问题。诚实罗列未完成部分是本书可证伪、可积累立场的延伸。
\end{quote}

\begin{center}\rule{0.5\linewidth}{0.5pt}\end{center}

\subsection{17.1 G
常数更严格的完整闭合解析推导}\label{g-ux5e38ux6570ux66f4ux4e25ux683cux7684ux5b8cux6574ux95edux5408ux89e3ux6790ux63a8ux5bfc}

第10章仅给出 G 的结构框架
\(G=\mu^2 c^2\mathcal{F}(\langle Lk\rangle,n_{\mathrm{wind}},\theta)\)。待完成：从
A1--A7 严格导出 \(\mathcal{F}\) 显式形式与桥接常数 \(\mu\) 的值，并复现
CODATA‑2022 的
\(G=6.67430(15)\times10^{-11}\,\mathrm{m^3kg^{-1}s^{-2}}\)。

\subsection{17.2
完整动力学场微分方程组构建}\label{ux5b8cux6574ux52a8ux529bux5b66ux573aux5faeux5206ux65b9ux7a0bux7ec4ux6784ux5efa}

目前仅有静态齐性真空与局域激发的拓扑描述。缺失：描述真空流线密度场演化的场方程（类流体/类规范场方程），以及其与爱因斯坦场方程（或某种修正引力）的严格对应关系（呼应
14.5）。

\subsection{17.3
量子引力路径积分的几何版本}\label{ux91cfux5b50ux5f15ux529bux8defux5f84ux79efux5206ux7684ux51e0ux4f55ux7248ux672c}

见
\texttt{quantum\_gravity\_pathint.md}：把费曼路径积分的``对所有几何求和''替换为``对所有螺旋系综构型求和''，构造几何路径积分。

\subsection{17.4
暗物质、暗能量在螺旋真空框架下的解释可能性}\label{ux6697ux7269ux8d28ux6697ux80fdux91cfux5728ux87baux65cbux771fux7a7aux6846ux67b6ux4e0bux7684ux89e3ux91caux53efux80fdux6027}

候选思路：(a) 暗能量 = 真空螺旋丝网的零点缠绕能密度；(b) 暗物质 =
大尺度缠绕密度涨落产生的等效度规扰动（不与粒子探测截面耦合）。均属初步构想，需数值/观测检验。

\subsection{17.5
量子测量问题几何视角}\label{ux91cfux5b50ux6d4bux91cfux95eeux9898ux51e0ux4f55ux89c6ux89d2}

提议：波函数坍缩 =
局域纽结从叠加态``选择''某一拓扑类型（退相干源于与背景丝网的拓扑耦合）。这是概念性提议，远非解决方案。

\subsection{17.6 分支2（变 κτ
宇宙）进一步研究价值}\label{ux5206ux652f2ux53d8-ux3baux3c4-ux5b87ux5b99ux8fdbux4e00ux6b65ux7814ux7a76ux4ef7ux503c}

即使 A3 成立，研究分支2 仍有价值：(a)
作为宇宙早期（对称性尚未齐化）的候选相；(b) 提供 \(\alpha\)
漂移的备用预言通道；(c) 检验理论框架的鲁棒性。建议建立分支2 的场论版本。

\begin{center}\rule{0.5\linewidth}{0.5pt}\end{center}

\emph{（本章未完，量子引力路径积分几何版见下篇。）}

\section{第17章（续）/ 附录性
量子引力路径积分的几何版本}\label{ux7b2c17ux7ae0ux7eed-ux9644ux5f55ux6027-ux91cfux5b50ux5f15ux529bux8defux5f84ux79efux5206ux7684ux51e0ux4f55ux7248ux672c}

\begin{quote}
本节对应开放问题
17.3，提出用``螺旋系综构型空间''替代``度规空间''作为量子引力路径积分的求和对象。
\end{quote}

\begin{center}\rule{0.5\linewidth}{0.5pt}\end{center}

\subsection{17.3.1
标准量子引力路径积分回顾}\label{ux6807ux51c6ux91cfux5b50ux5f15ux529bux8defux5f84ux79efux5206ux56deux987e}

广义相对论的量子化尝试中，哈金--路径积分写为

\[
Z=\int \mathcal{D}[g]\,e^{i S_{\mathrm{EH}}[g]/\hbar},
\]

对一切度规构型求和。困难在于发散与背景依赖。

\subsection{17.3.2 GAQ
几何路径积分提议}\label{gaq-ux51e0ux4f55ux8defux5f84ux79efux5206ux63d0ux8bae}

GAQ‑UFT 主张度规非基本，故改为对\textbf{螺旋系综构型}求和：

\[
Z_{\mathrm{GAQ}}=\int \mathcal{D}[\mathcal{E}]\,e^{i S_{\mathrm{helix}}[\mathcal{E}]/\hbar},
\]

其中 \(\mathcal{E}\) 为真空螺旋系综（由
\((\kappa(s),\tau(s),\text{位置},\text{朝向})\)
参数化），\(S_{\mathrm{helix}}\)
为系综作用量（待构造，候选为缠绕拓扑作用 + 曲率梯度项）。

\subsection{17.3.3
期望性质与未决问题}\label{ux671fux671bux6027ux8d28ux4e0eux672aux51b3ux95eeux9898}

\begin{itemize}
\tightlist
\item
  \textbf{经典极限：} 当
  \(\hbar\to0\)，鞍点应给出齐性螺旋系综（A3），恢复闵可夫斯基等效度规；
\item
  \textbf{离散化友好：} 系综可用有限螺旋数数值采样，避开连续度规发散；
\item
  \textbf{未决：} \(S_{\mathrm{helix}}\) 的具体形式、与标准 EH
  作用的等价性证明、重整化行为------全部开放。
\end{itemize}

\subsubsection{17.3.4
一个具体的候选作用量（试探，非结论）}\label{ux4e00ux4e2aux5177ux4f53ux7684ux5019ux9009ux4f5cux7528ux91cfux8bd5ux63a2ux975eux7ed3ux8bba}

为把构想推进一步，给出作用量的试探形式（明确标注为\textbf{研究起点}，非推导结论）：

\[
S_{\mathrm{helix}}[\mathcal{E}]
= \int d^4x\,\sqrt{-g_{\mathrm{eff}}}\,
\Bigl[\,\beta_1\,\|\nabla\kappa\|^2+\beta_2\,\|\nabla\tau\|^2
+\beta_3\,\mathrm{Lk}(\gamma_i,\gamma_j)
+\Lambda_{\mathrm{helix}}\Bigr],
\]

其中前两项是曲率/挠率场的梯度项（类动能），第三项是丝网内部缠绕数（类拓扑项），\(\Lambda_{\mathrm{helix}}\)
为等效宇宙学常数。系数 \(\beta_i\) 与 \(\Lambda_{\mathrm{helix}}\)
须由量纲自洽（A5）与实验标定，目前\textbf{全部未知}。

\subsubsection{17.3.5
与本书其它开放问题的耦合}\label{ux4e0eux672cux4e66ux5176ux5b83ux5f00ux653eux95eeux9898ux7684ux8026ux5408}

\begin{itemize}
\tightlist
\item
  若 17.2 的动力学场方程先建立，则 \(S_{\mathrm{helix}}\)
  可由该方程的拉格朗日量直接读取，二者应同步推进；
\item
  若 17.1 的 \(G\) 闭合成功，\(G\) 将自然出现在
  \(\Lambda_{\mathrm{helix}}\) 与 \(\beta_i\)
  的量纲分析中，提供一条把引力常数纳入几何路径积分的线索。
\end{itemize}

\begin{quote}
本节属研究方向构想，非已完成结果。其价值在于提供一条``不经度量张量''的量子引力入口，与本书``几何基元即螺旋''的一贯立场一致。任何把它当作``已建立量子引力理论''的表述都是误读。
\end{quote}

\begin{center}\rule{0.5\linewidth}{0.5pt}\end{center}

\emph{（第6卷主体完。全书总结与边界声明见 \texttt{summary\_ch18.md}。）}

\section{第18章 总结：GAQ‑UFT
完整逻辑闭环；理论局限与边界}\label{ux7b2c18ux7ae0-ux603bux7ed3gaquft-ux5b8cux6574ux903bux8f91ux95edux73afux7406ux8bbaux5c40ux9650ux4e0eux8fb9ux754c-1}

\begin{quote}
全书收束。我们重述逻辑闭环，并明确划定 GAQ‑UFT
已经做到什么、尚未做到什么，以避免任何过度宣称。
\end{quote}

\begin{center}\rule{0.5\linewidth}{0.5pt}\end{center}

\subsection{18.1
逻辑闭环回顾}\label{ux903bux8f91ux95edux73afux56deux987e}

\begin{verbatim}
A1 广义Kakeya覆盖 ┐
                  ├─► [κ>0, τ≠0] 数学必要条件（定理1，第5章）
A2 Frenet标架     ┘
                              │
        + A3 真空流线齐性（物理）──► 圆柱螺旋（分支3，第6–7章）
                              │
        + A4 相交/缠绕/打结 ─────► 粒子=拓扑激发（第8,12章）
        + A5 量纲自洽 ─────────► 常数映射可量纲校验（第9–10章）
        + A6 洛伦兹协变 ───────► 3+1维拓展（第13–14章）
        + A7 量子对应 ─────────► 量子数=整数拓扑不变量（第11–12章）
                              │
              可观测预言 + 实验约束（第15–16章）──► 可证伪判决
\end{verbatim}

\subsection{18.2
已经完成的部分}\label{ux5df2ux7ecfux5b8cux6210ux7684ux90e8ux5206}

下面逐项给出``已完成''的具体含义与证据锚点，便于读者核验，而非接受口号式断言。

\begin{enumerate}
\def\labelenumi{\arabic{enumi}.}
\tightlist
\item
  \textbf{严格的数学地基（第1章）。} Frenet--Serret 方程组
  \(\mathbf{T}'=\kappa\mathbf{N}\)、\(\mathbf{N}'=-\kappa\mathbf{T}+\tau\mathbf{B}\)、\(\mathbf{B}'=-\tau\mathbf{N}\)
  由弧长参数化与标准正交标架唯一性严格导出；圆柱螺旋在其曲率 \(\kappa\)
  与挠率 \(\tau\) 均为常数时被唯一确定（差欧氏运动），并给出显式参数
  \(\mathbf{r}(s)=(\frac{\cos(\tau s)}{\sqrt{\kappa^2+\tau^2}},\frac{\sin(\tau s)}{\sqrt{\kappa^2+\tau^2}},\frac{\kappa s}{\sqrt{\kappa^2+\tau^2}})\)。验证套件
  \textbf{T2}：数值计算的 \(\kappa=\tau=1/2\) 与解析值偏差
  \(<10^{-16}\)。
\item
  \textbf{拓扑 + 覆盖双约束排除直线/平面生成元（第3、5章）。} 毛球定理
  \(\chi(S^2)=2\) 保证 \(S^2\) 上无处处非零连续切场（验证
  \textbf{T9}：指数和 \(=2.0000\)）；广义 Kakeya
  覆盖（第2章）要求每点方向全覆盖，二者合并推出生成元必须 \(\kappa>0\)
  且 \(\tau\neq0\)（定理
  1）。这两者都是\textbf{纯数学必要条件}，不依赖任何物理常数。
\item
  \textbf{解空间四大分支的清晰划分与反例数值验证（第6章）。} 以
  \(\kappa,\tau\) 的常数性为判据，解空间被分为：分支
  0（\(\kappa=0\)，直线，拓扑禁止）、分支
  1（\(\tau=0\)，平面圆，几何禁止）、分支 3（\(\kappa,\tau\)
  常数，工作分支＝圆柱螺旋）、分支 2（\(\kappa,\tau\)
  变，数学允许、物理代价）。验证 \textbf{T10}：四类样例被自动正确归类。
\item
  \textbf{真空螺旋基元与系综构造（第7章）。}
  给出单根螺旋的几何参数与无穷族系综的拼装方式，使 A1
  覆盖在离散网格上可达 100\%。
\item
  \textbf{粒子作为拓扑激发的本体论框架（第8、12章）。}
  把粒子识别为螺旋丝网的局域纽结/缠绕缺陷，量子数对应整数拓扑不变量（缠绕数、环绕数），并给出德布罗意关系
  \(\lambda=h/p\) 从螺旋周期 \(L=2\pi/\sqrt{\kappa^2+\tau^2}\)
  的几何导出（12.4）。
\item
  \textbf{常数的几何映射与量纲校验（第9--10章）。}
  \(\alpha\leftrightarrow\theta\) 经 \(\theta=\arctan\alpha\)
  标定，\(\theta=0.4181^\circ\) 回代误差为 0（验证 \textbf{T4}）；Koide
  经验式被数值复现，相对偏差 \(9.2\times10^{-6}\)（验证
  \textbf{T6}）。但须注意：这些都还停留在``映射/经验契合''层面，\textbf{未上升为第一性原理推导}（见
  18.3）。
\item
  \textbf{全套 Python 200 位精度可复现验证（Supplemental\_Code）。}
  \texttt{verification\_suite\_v2.py} 当前 10 项检验、7 PASS / 0 FAIL /
  2 OPEN / 1 INFO，全部硬编码 CODATA‑2022 数值，任何人可一键复现。
\item
  \textbf{明确的实验判决通道（第15--16章）。}
  给出可证伪预言：精细结构常数宇宙学漂移（判决阈值
  \(\|\Delta\alpha/\alpha\|\sim10^{-6}\)，验证 \textbf{T12}
  INFO）、真空色散、引力波偏振异常、宇宙射线各向异性、Koide 偏离。
\end{enumerate}

\subsection{18.3 尚未闭合 /
诚实边界}\label{ux5c1aux672aux95edux5408-ux8bdaux5b9eux8fb9ux754c}

以下每一项都\textbf{不是``已有雏形、只差整理''}，而是真正的开放研究问题；标
OPEN 的目的是把理论交还给实验与后来者，而非假装已解。

\begin{itemize}
\tightlist
\item
  \textbf{G 的完整解析推导（17.1）未完成。} 第10章仅建立结构式
  \(G=\mu^2 c^2\mathcal{F}(\langle Lk\rangle,n_{\mathrm{wind}},\theta)\)，但
  \(\mathcal{F}\) 的显式形式与桥接常数 \(\mu\) 尚不能由 A1--A7
  严格导出，因而无法由 \(\theta=0.418^\circ\) 反推 CODATA‑2022 的
  \(G=6.67430(15)\times10^{-11}\)。验证 \textbf{T13} 因此标记为 OPEN。
\item
  \textbf{引电统一方程的桥接常数 \(K_{\mathrm{bridge}}\)
  几何来源未闭合（10.7）。} 方程
  \(G\varepsilon_0=K_{\mathrm{bridge}}(\kappa_e/\kappa_\Omega)^2\)
  左右量纲不匹配（左侧 \([M^{-2}Q^2]\)，右侧无量纲），必须引入带量纲
  \([M^{-2}Q^2]\) 的
  \(K_{\mathrm{bridge}}\)；当前它只能作为\textbf{经验拟合}（比值取 1 时
  \(K_{\mathrm{bridge}}\approx5.91\times10^{-22}\)），而
  \(\kappa_e,\kappa_\Omega\)
  须从电磁理论与真空几何\textbf{独立}标定后方可称``推导''。验证
  \textbf{T11} 标记为 OPEN。
\item
  \textbf{完整动力学场方程与爱因斯坦场方程的严格对应缺失（17.2）。}
  目前只有静态齐性真空与局域激发的拓扑描述，没有描述丝网密度场演化的场方程，也没有与爱因斯坦场方程的严格对应。
\item
  \textbf{M‑α 映射函数 \(f\) 的几何推导未完成（9.1）。} 质量与
  \(\alpha\)（或
  \(\theta\)）之间的映射目前是经验/唯象形式，缺少从公理出发的几何导出。
\item
  \textbf{夸克质量谱系未闭合（11.6）。} 第11章只处理了轻子（电子/μ/τ）的
  Koide
  关系，夸克质量、强弱相互作用、对称性破缺、暴胀、重子生成等标准模型核心现象全书未触及。
\item
  \textbf{量子引力路径积分仅为方向构想（17.3）。}
  仅提出``对螺旋系综构型求和''的入口，作用量 \(S_{\mathrm{helix}}\)
  的具体形式、与 EH 作用的等价性、重整化行为全部开放。
\end{itemize}

\subsubsection{18.3.1
理论可能的失败模式}\label{ux7406ux8bbaux53efux80fdux7684ux5931ux8d25ux6a21ux5f0f}

为落实可证伪性，明确列出可能导致 GAQ‑UFT 被否证的情形：

\begin{enumerate}
\def\labelenumi{\arabic{enumi}.}
\tightlist
\item
  \textbf{α 漂移反证：} 若未来类星体光谱把 \(\|\Delta\alpha/\alpha\|\)
  的上限压到 \(10^{-8}\) 量级而仍未观测到漂移，且本书分支 2
  预言通道也无法调和，则螺旋升角宇宙学演化的核心预言受挫。
\item
  \textbf{G 长期闭合失败：} 若经合理努力仍无法由 A1--A7 给出 \(G\)
  的数值闭环，则``几何基元统一常数''的立场退化为``唯象拟合''，说服力大幅削弱。
\item
  \textbf{Koide 偏离超出实验误差：} 若轻子质量未来测量使 Koide
  相对偏差显著超出 \(10^{-5}\)，则第11章的几何解释失效。
\end{enumerate}

这些不是``已经发生的失败''，而是理论必须正面迎接的裁决。

\subsection{18.4
理论定位与哲学立场}\label{ux7406ux8bbaux5b9aux4f4dux4e0eux54f2ux5b66ux7acbux573a}

GAQ‑UFT 不是已完成的终极理论，而是一套\textbf{逻辑透明、数学严格前部 +
物理假设清晰标注 + 可计算可证伪}的统一场论候选框架。它真正的贡献在于：

\begin{itemize}
\tightlist
\item
  用最少几何参数（单一螺旋升角 \(\theta\)）集约解释多个常数；
\item
  严格区分``数学必要条件''与``物理公设''，杜绝伪推导；
\item
  把可证伪性落到实处为具体观测判决。
\end{itemize}

物理学的进步靠可被现实检验的体系，而非完美无缺的故事。GAQ‑UFT
正是这样一次从几何第一性原理出发、向统一场论推进的完整尝试------它的价值，最终由实验与后来者的工作裁决。关于``已破解所有宇宙秘密''之类的宣称，本书在
\texttt{精算验证总结与核心方程清单.md}
中已作专门诚实声明：\textbf{任何此类宣称都是伪科学}，本书全部价值恰在于划清``已证
/ 待证 / 超出范围''的边界。

\begin{center}\rule{0.5\linewidth}{0.5pt}\end{center}

\emph{（全书正文完。附录 A/B/C 见 \texttt{Appendices/}。）}
